\documentclass[11pt]{article}
\usepackage[paperwidth=210mm,paperheight=297mm,margin=16mm]{geometry}
\usepackage{fontspec,amsmath,array,multirow,graphicx}
\usepackage{polyglossia}
\setmainlanguage{latin}\setotherlanguages{arabic,greek}
\setmainfont{Linux Libertine O}[Ligatures=TeX]
\newfontfamily\arabicfont{Amiri}[Script=Arabic]
\newfontfamily\greekfont{FreeSerif}[Script=Greek]
\newfontfamily\NallinoSigns{NallinoSigns.otf}[Path=./fonts/]
\newcommand{\AbjadZero}{{\NallinoSigns\symbol{"E000}}}
\newcommand{\AbjadOver}[1]{\leavevmode\vbox{\hrule height .45pt\kern1.2pt\hbox{#1}}}
\newcommand{\Name}[1]{{\addfontfeatures{LetterSpace=12}#1}}
\pagestyle{empty}\setlength{\parindent}{0pt}
\renewcommand{\arraystretch}{1.18}

\begin{document}
\clearpage
\noindent\makebox[\textwidth]{\textbf{144}\hfill{\small C. A. NALLINO}\hfill\textbf{}}\par\vspace{-1mm}\noindent\rule{\textwidth}{.4pt}\par\vspace{2mm}
\begin{center}{\small Fol. 226,v.-237,r. = Textus, t. III, p. 245-274.}\end{center}
\begin{center}{\large\bfseries Situs et magnitudines stellarum fixarum\\anno 1191 a Dhū ’l-qarnayn.}\end{center}
\begin{center}\rule{30mm}{.4pt}\end{center}
\begin{center}I.\end{center}
\begin{center}Constellationes quae sunt a septentrionibus zodiaci.\end{center}
\begin{center}\begin{tabular}{r||p{78mm}|r@{\hspace{2.2mm}}r|r@{\hspace{2.2mm}}r|c|c||l}\cline{2-8}
 & \centering Stellarum descriptio. & \multicolumn{2}{c|}{Longitudo.} & \multicolumn{2}{c|}{Latitudo.} & \parbox{9mm}{\centering\small Plaga Caeli.} & Magnitudo. & \\ \cline{2-8}
 & \multicolumn{7}{c||}{\parbox{150mm}{\centering\rule{0pt}{5mm}\textbf{Ursa minor} (\textit{ad–dubb al–aṣghar}).}} & \\ \cline{2-8}
1 & Stella quae est in extremitate caudae Ursae\newline\hspace*{4mm}minoris. & \raisebox{-1\normalbaselineskip}[0pt][0pt]{71°} & \raisebox{-1\normalbaselineskip}[0pt][0pt]{20′} & \raisebox{-1\normalbaselineskip}[0pt][0pt]{66°} & \raisebox{-1\normalbaselineskip}[0pt][0pt]{0′} & \raisebox{-1\normalbaselineskip}[0pt][0pt]{bor.} & \raisebox{-1\normalbaselineskip}[0pt][0pt]{3} & \raisebox{-1\normalbaselineskip}[0pt][0pt]{α} \\
2 & Stella quae est in radice caudae huius Ursae. & \textit{87} & \textit{10} & 74 & 20 & b & 4 & ε \\
3 & Australior in sequenti latere quadranguli. & 118 & 20 & 72 & 50 & b & 2 & β \\
4 & Borealior in hoc sequenti latere quadranguli. & 127 & 20 & 74 & 50 & b & 2 & γ \\
 & \multicolumn{7}{c||}{\parbox{150mm}{\centering\rule{0pt}{5mm}\textbf{Ursa maior} (\textit{ad–dubb al–akbar}).}} & \\ \cline{2-8}
1 & Stella quae est in anteriore parte rictus Ur-\newline\hspace*{4mm}sae maioris. & \raisebox{-1\normalbaselineskip}[0pt][0pt]{96} & \raisebox{-1\normalbaselineskip}[0pt][0pt]{30} & \raisebox{-1\normalbaselineskip}[0pt][0pt]{39} & \raisebox{-1\normalbaselineskip}[0pt][0pt]{50} & \raisebox{-1\normalbaselineskip}[0pt][0pt]{b} & \raisebox{-1\normalbaselineskip}[0pt][0pt]{4} & \raisebox{-1\normalbaselineskip}[0pt][0pt]{ο} \\
2 & Quae est in sinistro huius Ursae genu. & 111 & 50 & 35 & 0 & b & 3 & θ \\
3 & Septentrionalis quae est in extremitate pedis\newline\hspace*{4mm}anterioris sinistri. & \raisebox{-1\normalbaselineskip}[0pt][0pt]{106} & \raisebox{-1\normalbaselineskip}[0pt][0pt]{40} & \raisebox{-1\normalbaselineskip}[0pt][0pt]{29} & \raisebox{-1\normalbaselineskip}[0pt][0pt]{20} & \raisebox{-1\normalbaselineskip}[0pt][0pt]{b} & \raisebox{-1\normalbaselineskip}[0pt][0pt]{3} & \raisebox{-1\normalbaselineskip}[0pt][0pt]{ι} \\
4 & Australis in hoc pede. & 107 & 30 & 28 & 20 & b & 3 & κ \\
5 & Stella quae est in dorso Ursae in quadran-\newline\hspace*{4mm}gulo. & \raisebox{-1\normalbaselineskip}[0pt][0pt]{118} & \raisebox{-1\normalbaselineskip}[0pt][0pt]{50} & \raisebox{-1\normalbaselineskip}[0pt][0pt]{49} & \raisebox{-1\normalbaselineskip}[0pt][0pt]{0} & \raisebox{-1\normalbaselineskip}[0pt][0pt]{b} & \raisebox{-1\normalbaselineskip}[0pt][0pt]{2} & \raisebox{-1\normalbaselineskip}[0pt][0pt]{α} \\
\cline{2-8}\end{tabular}\end{center}
{\small
\par\hspace*{8mm}\textbf{Ursa minor} (1):
\par\hangindent=12mm\hangafter=1 1. Cod. « in extremitate genus » (\textarabic{ركبة} pro \textarabic{ذنب}).
\par\hangindent=12mm\hangafter=1 2. Long. cod. \textarabic{قا ك} 101° 20′ pro \textarabic{فز ى} 87° 10′ (ut Pt., S.) vel \textarabic{فا م} 81° 40′ (ut codex quidam Ptolemaei).
\par\hangindent=12mm\hangafter=1 3. Cod. « in secundo » \textarabic{الثاني} pro « in sequenti » \textarabic{التالي}; Pt. ἐν τῇ ἑπομένῃ πλευρᾷ, ut S. — Long. cod. \textarabic{قنج} 153° pro \textarabic{قيح} 118°; minuta congruunt cum L. et S. (sed H. et B. 30′). — Minuta lat. cod. \textarabic{ى} 10′ pro \textarabic{ن} 50′.
\par\hangindent=12mm\hangafter=1 4. Cod. « in secundo » ut nr. 3. — Long. cod. \textarabic{قنز} 157° pro \textarabic{قكز} 127°.
\par\hspace*{8mm}\textbf{Ursa maior}:
\par\hangindent=12mm\hangafter=1 1. Long. cod. \textarabic{قعو} 176° pro \textarabic{ضو} 96°. — Minuta lat. cod. \textarabic{كب} 22′ pro \textarabic{ن} 50′.
\par}
\par\vspace{2mm}\begin{center}\rule{40mm}{.4pt}\end{center}{\small (1) Conferantur animadversiones generales post finem tabularum omnium.}
\clearpage
\noindent\makebox[\textwidth]{\textbf{}\hfill{\small AL-BATTANI OPUS ASTRONOMICUM}\hfill\textbf{145}}\par\vspace{-1mm}\noindent\rule{\textwidth}{.4pt}\par\vspace{2mm}
\begin{center}\begin{tabular}{r||p{78mm}|r@{\hspace{2.2mm}}r|r@{\hspace{2.2mm}}r|c|c||l}\cline{2-8}
 & \centering Stellarum descriptio. & \multicolumn{2}{c|}{Longitudo.} & \multicolumn{2}{c|}{Latitudo.} & \parbox{9mm}{\centering\small Plaga Caeli.} & Magnitudo. & \\ \cline{2-8}
6 & Quae est in tenuiore parte ventris. & 123° & 20′ & 44° & 30′ & b & 2 & β \\
7 & Quae est in radice caudae. & 134 & \textit{20} & 51 & 0 & b & 3 & δ \\
8 & Quae est in radice poplitis sinistri poste-\newline\hspace*{4mm}rioris. & \raisebox{-1\normalbaselineskip}[0pt][0pt]{134} & \raisebox{-1\normalbaselineskip}[0pt][0pt]{10} & \raisebox{-1\normalbaselineskip}[0pt][0pt]{46} & \raisebox{-1\normalbaselineskip}[0pt][0pt]{30} & \raisebox{-1\normalbaselineskip}[0pt][0pt]{b} & \raisebox{-1\normalbaselineskip}[0pt][0pt]{2} & \raisebox{-1\normalbaselineskip}[0pt][0pt]{γ} \\
9 & Anterior quae est in extremitate pedis sini-\newline\hspace*{4mm}stri posterioris. & \raisebox{-1\normalbaselineskip}[0pt][0pt]{123} & \raisebox{-1\normalbaselineskip}[0pt][0pt]{50} & \raisebox{-1\normalbaselineskip}[0pt][0pt]{29} & \raisebox{-1\normalbaselineskip}[0pt][0pt]{20} & \raisebox{-1\normalbaselineskip}[0pt][0pt]{b} & \raisebox{-1\normalbaselineskip}[0pt][0pt]{3} & \raisebox{-1\normalbaselineskip}[0pt][0pt]{λ} \\
10 & Quae eam sequitur. & 125 & 20 & 28 & 15 & b & 3 & μ \\
11 & Borealis ex iis quae sunt in extremitate\newline\hspace*{4mm}pedis dexteri posterioris. & \raisebox{-1\normalbaselineskip}[0pt][0pt]{141} & \raisebox{-1\normalbaselineskip}[0pt][0pt]{0} & \raisebox{-1\normalbaselineskip}[0pt][0pt]{25} & \raisebox{-1\normalbaselineskip}[0pt][0pt]{50} & \raisebox{-1\normalbaselineskip}[0pt][0pt]{b} & \raisebox{-1\normalbaselineskip}[0pt][0pt]{3} & \raisebox{-1\normalbaselineskip}[0pt][0pt]{ν} \\
12 & Australis ex iis. & 141 & 30 & 25 & 0 & b & 3 & ξ \\
13 & Anterior e tribus quae sunt in cauda. & 143 & 20 & 53 & 30 & b & 2 & ε \\
14 & Media e tribus his. & 149 & 10 & 55 & 40 & b & 2 & ζ \\
15 & Tertia quae est in extremitate caudae. & 161 & 0 & 54 & 0 & b & 2 & η \\
 & \multicolumn{7}{l}{\hspace*{4mm}Et ex eius [stellis] quae in figura Ursae non sunt:} & \\
16 & Quae est sub cauda a parte australi. & 159 & 0 & 39 & 45 & b & 3 & 12 Canum \\
17 & Quae est inter pedem anteriorem Ursae et\newline\hspace*{4mm}caput Leonis. & \raisebox{-1\normalbaselineskip}[0pt][0pt]{116} & \raisebox{-1\normalbaselineskip}[0pt][0pt]{10} & \raisebox{-1\normalbaselineskip}[0pt][0pt]{17} & \raisebox{-1\normalbaselineskip}[0pt][0pt]{15} & \raisebox{-1\normalbaselineskip}[0pt][0pt]{b} & \raisebox{-1\normalbaselineskip}[0pt][0pt]{4} & \raisebox{-1\normalbaselineskip}[0pt][0pt]{40 Lync.} \\
18 & Obscura quae reliquas tres obscuras se-\newline\hspace*{4mm}quitur. & \raisebox{-1\normalbaselineskip}[0pt][0pt]{117} & \raisebox{-1\normalbaselineskip}[0pt][0pt]{20} & \raisebox{-1\normalbaselineskip}[0pt][0pt]{20} & \raisebox{-1\normalbaselineskip}[0pt][0pt]{0} & \raisebox{-1\normalbaselineskip}[0pt][0pt]{b} & \raisebox{-1\normalbaselineskip}[0pt][0pt]{obscura} & \raisebox{-1\normalbaselineskip}[0pt][0pt]{10 Leo. M.} \\
19 & Obscura hanc praecedens. & 113 & 20 & 22 & 50 & b & obscura &  \\
20 & Obscura ante hanc. & 112 & 20 & \textit{20} & \textit{20} & b & obscura &  \\
21 & Quae est ante hanc, inter eam et Geminos. & 101 & 10 & 22 & 15 & b & obscura & 31 Lync. \\
\cline{2-8}\end{tabular}\end{center}
{\small
\par\hangindent=12mm\hangafter=1 6. Gradus long. cod. \textarabic{قلج} 133° pro \textarabic{قكج} 123°; minuta ut L. et S. (sed B. et H. 40′). — Minuta lat. cod. \textarabic{ى} 10′ pro \textarabic{ل} 30′.
\par\hangindent=12mm\hangafter=1 7. Min. long. cod. \textarabic{\AbjadZero{}} 0′; L. et S. 20′; B. et H. 40′.
\par\hangindent=12mm\hangafter=1 8. Long. cod. \textarabic{قلج ل} 133° 30′ pro \textarabic{قلد ى} 134° 10′.
\par\hangindent=12mm\hangafter=1 9. Long. cod. \textarabic{ل} 30′ pro \textarabic{ن} 50′. — Lat. cod. \textarabic{مط} 49° pro \textarabic{كط} 29°.
\par\hangindent=12mm\hangafter=1 10. Lat. cod. \textarabic{نه} 55′ pro \textarabic{يه} 15′.
\par\hangindent=12mm\hangafter=1 11. Cod. « australis »! — Lat. cod. \textarabic{مه} 45° pro \textarabic{كه} 25°.
\par\hangindent=12mm\hangafter=1 12. Lat. cod. \textarabic{كح} 28° pro \textarabic{كه} 25°.
\par\hangindent=12mm\hangafter=1 13. Long. cod. \textarabic{قعج} 173° pro \textarabic{قمج} 143°.
\par\hangindent=12mm\hangafter=1 16. Lat. cod. \textarabic{نط} 59° pro \textarabic{لط} 39°.
\par\hangindent=12mm\hangafter=1 17. Lat. cod. \textarabic{نز} 57° pro \textarabic{يز} 17°.
\par\hangindent=12mm\hangafter=1 19. Lat. cod. \textarabic{كد} 24° pro \textarabic{كب} 22°; minuta ut B. et H., sed L. et S. 45′. — \Name{Baily}: « I cannot find any « star that corresponds with the position given by Ptolemy ».
\par\hangindent=12mm\hangafter=1 20. Lat. discrepans a B., H., S. (qui 23° 0′ habent) et a L. (23° 15′). Error fortasse iam in Graeco aliquo exemplari (κγ′ 20° $\tfrac{1}{3}$ pro κγ 23°). — \Name{Baily}: « I cannot find any star that corresponds with the « position given by Ptolemy ».
\par}
\par\vfill\hfill 19
\clearpage
\noindent\makebox[\textwidth]{\textbf{146}\hfill{\small C. A. NALLINO}\hfill\textbf{}}\par\vspace{-1mm}\noindent\rule{\textwidth}{.4pt}\par\vspace{2mm}
\begin{center}\begin{tabular}{r||p{78mm}|r@{\hspace{2.2mm}}r|r@{\hspace{2.2mm}}r|c|c||l}\cline{2-8}
 & \centering Stellarum descriptio. & \multicolumn{2}{c|}{Longitudo.} & \multicolumn{2}{c|}{Latitudo.} & \parbox{9mm}{\centering\small Plaga Caeli.} & Magnitudo. & \\ \cline{2-8}
 & \multicolumn{7}{c||}{\parbox{150mm}{\centering\rule{0pt}{5mm}\textbf{Draco} (\textit{tinnīn}).}} & \\ \cline{2-8}
1 & Stella quae est in extremitate linguae Dra-\newline\hspace*{4mm}conis. & \raisebox{-1\normalbaselineskip}[0pt][0pt]{217°} & \raisebox{-1\normalbaselineskip}[0pt][0pt]{50′} & \raisebox{-1\normalbaselineskip}[0pt][0pt]{76°} & \raisebox{-1\normalbaselineskip}[0pt][0pt]{30′} & \raisebox{-1\normalbaselineskip}[0pt][0pt]{b} & \raisebox{-1\normalbaselineskip}[0pt][0pt]{4} & \raisebox{-1\normalbaselineskip}[0pt][0pt]{μ} \\
2 & Stella quae est supra caput. & 250 & 50 & 75 & 30 & b & 3 & γ \\
3 & Stella quae est supra oculum Draconis. & 234 & 20 & 75 & 40 & b & 3 & β \\
4 & Stella ex triangulo quae est a parte occi-\newline\hspace*{4mm}dentali. & \raisebox{-1\normalbaselineskip}[0pt][0pt]{169} & \raisebox{-1\normalbaselineskip}[0pt][0pt]{30} & \raisebox{-1\normalbaselineskip}[0pt][0pt]{84} & \raisebox{-1\normalbaselineskip}[0pt][0pt]{50} & \raisebox{-1\normalbaselineskip}[0pt][0pt]{b} & \raisebox{-1\normalbaselineskip}[0pt][0pt]{3} & \raisebox{-1\normalbaselineskip}[0pt][0pt]{ζ} \\
5 & Borealis e duabus quae versus occidentem\newline\hspace*{4mm}sequuntur. & \raisebox{-1\normalbaselineskip}[0pt][0pt]{171} & \raisebox{-1\normalbaselineskip}[0pt][0pt]{10} & \raisebox{-1\normalbaselineskip}[0pt][0pt]{78} & \raisebox{-1\normalbaselineskip}[0pt][0pt]{0} & \raisebox{-1\normalbaselineskip}[0pt][0pt]{b} & \raisebox{-1\normalbaselineskip}[0pt][0pt]{3} & \raisebox{-1\normalbaselineskip}[0pt][0pt]{η} \\
6 & E duabus illis quae sequitur praecedentem\newline\hspace*{4mm}longinquam. & \raisebox{-1\normalbaselineskip}[0pt][0pt]{142} & \raisebox{-1\normalbaselineskip}[0pt][0pt]{20} & \raisebox{-1\normalbaselineskip}[0pt][0pt]{65} & \raisebox{-1\normalbaselineskip}[0pt][0pt]{30} & \raisebox{-1\normalbaselineskip}[0pt][0pt]{b} & \raisebox{-1\normalbaselineskip}[0pt][0pt]{3} & \raisebox{-1\normalbaselineskip}[0pt][0pt]{α} \\
7 & Quae hanc stellam sequitur. & 120 & 20 & 61 & 15 & b & 3 & κ \\
8 & Stella quae est in extremitate caudae Dra-\newline\hspace*{4mm}conis. & \raisebox{-1\normalbaselineskip}[0pt][0pt]{114} & \raisebox{-1\normalbaselineskip}[0pt][0pt]{20} & \raisebox{-1\normalbaselineskip}[0pt][0pt]{56} & \raisebox{-1\normalbaselineskip}[0pt][0pt]{15} & \raisebox{-1\normalbaselineskip}[0pt][0pt]{b} & \raisebox{-1\normalbaselineskip}[0pt][0pt]{3} & \raisebox{-1\normalbaselineskip}[0pt][0pt]{λ} \\
 & \multicolumn{7}{c||}{\parbox{150mm}{\centering\rule{0pt}{5mm}\textbf{Cephaeus} (\textit{Qīfāwus}, sive \textit{al–multahib} « exardescens »).}} & \\ \cline{2-8}
1 & Tangens in humero dextero. & 357 & 50 & 69 & 0 & b & 3 & α \\
2 & Tangens in cubito dextero. & 350 & 30 & 72 & 0 & b & 4 & η \\
3 & Stella quae est in brachio sinistro. & 18 & 40 & 62 & 30 & b & 4 magna & ι \\
4 & Australis e tribus quae sunt in tiara. & 357 & 30 & 60 & 15 & b & 5 & ε \\
\cline{2-8}\end{tabular}\end{center}
{\small
\par\hspace*{8mm}\textbf{Draco}:
\par\hangindent=12mm\hangafter=1 1. Long. cod. \textarabic{ل} 30′ pro \textarabic{ن} 50′.
\par\hangindent=12mm\hangafter=1 2. Long. cod. \textarabic{رنا ى} 251° 10′ pro \textarabic{رن ن} 250° 50′. — Lat. cod. \textarabic{ى} 10′ pro \textarabic{ل} 30′.
\par\hangindent=12mm\hangafter=1 3. Lat. cod. \textarabic{عد} 74° pro \textarabic{عه} 75°.
\par\hangindent=12mm\hangafter=1 4. Lat. cod. \textarabic{فه} 85° pro \textarabic{فد} 84°.
\par\hangindent=12mm\hangafter=1 7. Lat. cod. \textarabic{مه} 45′ pro \textarabic{ٮه} = \textarabic{يه} 15′; B., H. et S. 15′; L. 35′ (\textarabic{له}).
\par\hangindent=12mm\hangafter=1 8. Long. cod. \textarabic{م} 40′ pro \textarabic{ل} 20′.
\par\hspace*{8mm}\textbf{Cephaeus}:
\par\hangindent=12mm\hangafter=1 2. Lat. cod. \textarabic{يه} 15′ pro \textarabic{\AbjadZero{}} 0′. — Magnitudo cod. « 4 magna »; sed procul dubio « magna » pertinet ad stellam sequentem, ut Pts., L., S. et Arg. habent.
\par\hangindent=12mm\hangafter=1 3. Long. cod. \textarabic{صح يج} 68° 13′ pro \textarabic{يح م} 18° 40′. — Magnitudo cod. 4, sed cfr. adn. praecedentem.
\par\hangindent=12mm\hangafter=1 4. Cod. « media » \textarabic{المتوسط}; si hoc verum esset, long. 358° 30′, lat. 61° 15′ haberi deberent. Cfr. nr. 10 Bootis. — Long. cod. \textarabic{يج} 13′ pro \textarabic{ل} 30′; magnitudo cod. \textarabic{د} 4 (quod cum « media » conveniret) pro \textarabic{ه} 5.
\par}
\clearpage
\noindent\makebox[\textwidth]{\textbf{}\hfill{\small AL-BATTANI OPUS ASTRONOMICUM}\hfill\textbf{147}}\par\vspace{-1mm}\noindent\rule{\textwidth}{.4pt}\par\vspace{2mm}
\begin{center}\begin{tabular}{r||p{78mm}|r@{\hspace{2.2mm}}r|r@{\hspace{2.2mm}}r|c|c||l}\cline{2-8}
 & \centering Stellarum descriptio. & \multicolumn{2}{c|}{Longitudo.} & \multicolumn{2}{c|}{Latitudo.} & \parbox{9mm}{\centering\small Plaga Caeli.} & Magnitudo. & \\ \cline{2-8}
 & \multicolumn{7}{c||}{\parbox{150mm}{\centering\rule{0pt}{5mm}\textbf{Bootes} (\textit{al–ghūl} « daemon », sive \textit{ḥāris ash–shamāl} « custos septentrionis »,\\sive \textit{al–baqqār} « bubulcus »).}} & \\ \cline{2-8}
1 & Quae est in humero sinistro. & 180° & 50′ & 49° & 0′ & b & 3 & γ \\
2 & Stella quae est in capite. & 187 & 50 & 53 & 50 & b & 4 magna & β \\
3 & Stella quae est in humero dextero. & 196 & 50 & 48 & 40 & b & 4 magna & δ \\
4 & Stella quae est sub humero sinistro. & 198 & 50 & 46 & \textit{30} & b & 4 magna &  \\
5 & Quae est in poplite dextero, in cingulo et\newline\hspace*{4mm}vinculo. & \raisebox{-1\normalbaselineskip}[0pt][0pt]{191} & \raisebox{-1\normalbaselineskip}[0pt][0pt]{10} & \raisebox{-1\normalbaselineskip}[0pt][0pt]{40} & \raisebox{-1\normalbaselineskip}[0pt][0pt]{15} & \raisebox{-1\normalbaselineskip}[0pt][0pt]{b} & \raisebox{-1\normalbaselineskip}[0pt][0pt]{3} & \raisebox{-1\normalbaselineskip}[0pt][0pt]{ε} \\
6 & Praecedens e duabus quae sunt in cingulo. & 186 & 10 & 42 & 10 & b & 4 magna & ρ \\
7 & Stella quae est in talo dextero. & 196 & 30 & 28 & 0 & b & 3 & ζ \\
8 & Borealis ex tribus quae sunt in crure sini-\newline\hspace*{4mm}stro. & \raisebox{-1\normalbaselineskip}[0pt][0pt]{182} & \raisebox{-1\normalbaselineskip}[0pt][0pt]{30} & \raisebox{-1\normalbaselineskip}[0pt][0pt]{28} & \raisebox{-1\normalbaselineskip}[0pt][0pt]{0} & \raisebox{-1\normalbaselineskip}[0pt][0pt]{b} & \raisebox{-1\normalbaselineskip}[0pt][0pt]{3} & \raisebox{-1\normalbaselineskip}[0pt][0pt]{η} \\
9 & Media e tribus his. & 181 & 40 & 26 & 30 & b & 4 & τ \\
10 & Stella australis e tribus his. & 182 & 30 & 25 & 0 & b & 4 & υ \\
11 & \textit{as–Simāk ar–rāmiḥ} « praeeminens ha-\newline\hspace*{4mm}stam gerens », quae est inter duo fe-\newline\hspace*{4mm}mora daemonis, extra figuram. & \raisebox{-2\normalbaselineskip}[0pt][0pt]{188} & \raisebox{-2\normalbaselineskip}[0pt][0pt]{10} & \raisebox{-2\normalbaselineskip}[0pt][0pt]{31} & \raisebox{-2\normalbaselineskip}[0pt][0pt]{30} & \raisebox{-2\normalbaselineskip}[0pt][0pt]{b} & \raisebox{-2\normalbaselineskip}[0pt][0pt]{1} & \raisebox{-2\normalbaselineskip}[0pt][0pt]{α} \\
 & \multicolumn{7}{c||}{\parbox{150mm}{\centering\rule{0pt}{5mm}\textbf{Corona borealis} (\textit{al–fakkah}).}} & \\ \cline{2-8}
1 & Lucida e stellis Coronae. & 205 & 50 & 44 & 30 & b & 2 magna & α \\
2 & Praecedens e stellis Coronae. & 202 & 50 & 46 & 30 & b & 4 magna & β \\
\cline{2-8}\end{tabular}\end{center}
{\small
\par\hspace*{8mm}\textbf{Bootes}:
\par\hangindent=12mm\hangafter=1 2. Perperam lat. H. 57° 50′.
\par\hangindent=12mm\hangafter=1 3. Magnitudo cod. ut Pts. et S.; L. 4; Arg. 3.
\par\hangindent=12mm\hangafter=1 4. Pt. et S: « borealior e duabus quae sunt sub humero, in baculo ». — Lat. cod. ut H.; sed probabiliter pro \textarabic{ل} 30′ restituendum \textarabic{ى} 10′ ut B., L. et S. — Magnitudo ut B., H, L. et Pts.; S. « 5 magna ». — \Name{Baily}: « I am uncertain whether to designate this star as Flamsteed's 1 or 2 \textit{Coronae Borealis} »; \Name{Schjellerup} p. 65, adn. 1, proponit η Coronae, sed non sine haesitatione.
\par\hangindent=12mm\hangafter=1 5. Fortasse pro \textarabic{والرباط} « et vinculo » legendum \textarabic{وهي الرباط} « sive vinculo »; Pt. ὁ ἐπὶ τοῦ δεξιοῦ μηροῦ ἐν τῷ περιζώματι.
\par\hangindent=12mm\hangafter=1 6. Lat. cod. \textarabic{ل} 30′ pro \textarabic{ى} 10′; magnitudo ut Pts. et S. (L. sine « magna »).
\par\hangindent=12mm\hangafter=1 7. Long. cod. \textarabic{قصو} 166° pro \textarabic{قضو} 196°.
\par\hangindent=12mm\hangafter=1 8. Long. cod. \textarabic{قعب ى} 172° 10′ pro \textarabic{قفب ل} 182° 30′.
\par\hangindent=12mm\hangafter=1 10. Cod. « media »; error ut in nr. 4 Cephei. — Long. cod. \textarabic{ى} 10′ pro \textarabic{ل} 30′; magnitudo cod. \textarabic{ج} 3 pro \textarabic{د} 4.
\par\hspace*{8mm}\textbf{Corona australis}:
\par\hangindent=12mm\hangafter=1 1. Magnitudo ut L.; sed Pts., S. et Arg. sine « magna ».
\par\hangindent=12mm\hangafter=1 2. Magnitudo ut L. et Arg.; Pts. et S. sine « magna ».
\par}
\clearpage
\noindent\makebox[\textwidth]{\textbf{148}\hfill{\small C. A. NALLINO}\hfill\textbf{}}\par\vspace{-1mm}\noindent\rule{\textwidth}{.4pt}\par\vspace{2mm}
\begin{center}\begin{tabular}{r||p{78mm}|r@{\hspace{2.2mm}}r|r@{\hspace{2.2mm}}r|c|c||l}\cline{2-8}
 & \centering Stellarum descriptio. & \multicolumn{2}{c|}{Longitudo.} & \multicolumn{2}{c|}{Latitudo.} & \parbox{9mm}{\centering\small Plaga Caeli.} & Magnitudo. & \\ \cline{2-8}
 & \multicolumn{7}{c||}{\parbox{150mm}{\centering\rule{0pt}{5mm}\textbf{Hercules} (\textit{al–ǵāthī} « incumbens genibus »).}} & \\ \cline{2-8}
1 & Qnae est in capite Herculis. & 238° & 50′ & 37° & 30′ & b & 3 & α \\
2 & Quae est in humero dextero prope axillam. & 224 & 50 & 43 & 0 & b & 3 & β \\
3 & Quae est in brachio dextero. & 222 & 50 & 40 & 10 & b & 3 & γ \\
4 & Quae est in humero sinistro. & 237 & 50 & 48 & 0 & b & 3 & δ \\
5 & Quae est in brachio sinistro. & 243 & 10 & 49 & 30 & b & 4 magna & λ \\
6 & Quae est in cubito sinistro. & 248 & 50 & 52 & 0 & b & 4 magna & μ \\
7 & Australis ex tribus quae sunt in carpo si-\newline\hspace*{4mm}nistro. & \raisebox{-1\normalbaselineskip}[0pt][0pt]{252} & \raisebox{-1\normalbaselineskip}[0pt][0pt]{40} & \raisebox{-1\normalbaselineskip}[0pt][0pt]{53} & \raisebox{-1\normalbaselineskip}[0pt][0pt]{0} & \raisebox{-1\normalbaselineskip}[0pt][0pt]{b} & \raisebox{-1\normalbaselineskip}[0pt][0pt]{3} & \raisebox{-1\normalbaselineskip}[0pt][0pt]{ξ} \\
8 & Stella quae est in latere dextero. & 224 & 50 & 50 & 40 & b & 3 magna & ζ \\
9 & Quae est in radice poplitis sinistri. & 232 & 20 & 58 & 30 & b & 3 & c \\
10 & Quae eam sequitur ex tribus in poplite\newline\hspace*{4mm}sinistro. & \raisebox{-1\normalbaselineskip}[0pt][0pt]{237} & \raisebox{-1\normalbaselineskip}[0pt][0pt]{30} & \raisebox{-1\normalbaselineskip}[0pt][0pt]{61} & \raisebox{-1\normalbaselineskip}[0pt][0pt]{15} & \raisebox{-1\normalbaselineskip}[0pt][0pt]{b} & \raisebox{-1\normalbaselineskip}[0pt][0pt]{4 magna} & \raisebox{-1\normalbaselineskip}[0pt][0pt]{ρ} \\
11 & Quae est in radice poplitis dexteri. & 221 & 50 & 64 & 0 & b & 4 magna & η \\
12 & Quae est in genu dextero. & 206 & 50 & 65 & 30 & b & 4 magna & τ \\
 & \multicolumn{7}{c||}{\parbox{150mm}{\centering\rule{0pt}{5mm}\textbf{Lyra} (\textit{an–nasr al–wāqiʿ} « aquila ruens »).}} & \\ \cline{2-8}
1 & Lucida quae est in tiara tenentis lyram\newline\hspace*{4mm}(\textit{lūrah}); ipsaque est \textit{an–nasr} « aquila ». & \raisebox{-1\normalbaselineskip}[0pt][0pt]{268} & \raisebox{-1\normalbaselineskip}[0pt][0pt]{30} & \raisebox{-1\normalbaselineskip}[0pt][0pt]{62} & \raisebox{-1\normalbaselineskip}[0pt][0pt]{0} & \raisebox{-1\normalbaselineskip}[0pt][0pt]{b} & \raisebox{-1\normalbaselineskip}[0pt][0pt]{1} & \raisebox{-1\normalbaselineskip}[0pt][0pt]{α} \\
2 & Borealis e duabus ei propinquis. & 271 & 30 & 62 & 40 & b & 4 magna & ε \\
3 & Australis e duabus iis. & 271 & 30 & 61 & 0 & b & 4 magna & ζ \\
4 & Australis e duabus quae sunt in anteriore\newline\hspace*{4mm}[parte] praecedentis lancis librae. & \raisebox{-1\normalbaselineskip}[0pt][0pt]{272} & \raisebox{-1\normalbaselineskip}[0pt][0pt]{0} & \raisebox{-1\normalbaselineskip}[0pt][0pt]{55} & \raisebox{-1\normalbaselineskip}[0pt][0pt]{0} & \raisebox{-1\normalbaselineskip}[0pt][0pt]{b} & \raisebox{-1\normalbaselineskip}[0pt][0pt]{4 parva} & \raisebox{-1\normalbaselineskip}[0pt][0pt]{ν¹} \\
\cline{2-8}\end{tabular}\end{center}
{\small
\par\hspace*{8mm}\textbf{Hercules}:
\par\hangindent=12mm\hangafter=1 5. Magnitudo ut L. et Pts.; S. et Arg. 5; B. et H. 4.
\par\hangindent=12mm\hangafter=1 6. Magnitudo ut L. et Pts.; S. 4; Arg. 3 parva.
\par\hangindent=12mm\hangafter=1 7. Magnitudo L., Pts. et Arg. 4 magna; B. et H. 4; S. 4.
\par\hangindent=12mm\hangafter=1 8. Long. cod. \textarabic{ركز} 227° pro \textarabic{ركد} 224°. — Lat. cod. ut H.; B. 53° 0′; L. 56° 10′; S. 53° 10′. — Magnitudo ut L. et Arg.; Pts. et S. 3; B. et H. 4.
\par\hangindent=12mm\hangafter=1 9. Long. cod. \textarabic{رله} 235° pro \textarabic{رلب} 232°.
\par\hangindent=12mm\hangafter=1 10. Magnitudo ut L. et Pts.; S. et Arg. 4.
\par\hangindent=12mm\hangafter=1 11. Long. cod. \textarabic{ركج ى} 223° 10′ pro \textarabic{ركا ن} 221° 50′. — Lat. ut H. (ξδ); sed rectius L., B. et S. 60° 15′ (ξ ιε′). Magnitudo ut L. et Pts.; B., H. et S. 4; Arg. 3.
\par\hspace*{8mm}\textbf{Lyra}:
\par\hangindent=12mm\hangafter=1 1. Mira stellae descriptio; Pt. ὁ λαμπρὸς ἐπὶ τοῦ ὀστράκου καλούμενος λύρα « lucida in testa (testudinis); \textbf{Lyra} « vocata ».
\par\hangindent=12mm\hangafter=1 2. Magnitudo ut Pts. et S.; L. et Arg. 4.
\par\hangindent=12mm\hangafter=1 3. Magnitudo ut Pts. et S.; L. et Arg. 4.
\par\hangindent=12mm\hangafter=1 4. Mira descriptio; Pt.: ὁ νοτιώτερος αὐτῶν (i. e. τῶν ἐν τῷ ζυγώματι προηγουμένων β) « australior earum (i. « e. duarum praecedentium in zygomate lyrae) ».
\par}
\clearpage
\noindent\makebox[\textwidth]{\textbf{}\hfill{\small AL-BATTANI OPUS ASTRONOMICUM}\hfill\textbf{149}}\par\vspace{-1mm}\noindent\rule{\textwidth}{.4pt}\par\vspace{2mm}
\begin{center}\begin{tabular}{r||p{78mm}|r@{\hspace{2.2mm}}r|r@{\hspace{2.2mm}}r|c|c||l}\cline{2-8}
 & \centering Stellarum descriptio. & \multicolumn{2}{c|}{Longitudo.} & \multicolumn{2}{c|}{Latitudo.} & \parbox{9mm}{\centering\small Plaga Caeli.} & Magnitudo. & \\ \cline{2-8}
5 & Australis e duabus quae sunt in anteriore\newline\hspace*{4mm}[parte] posterioris lancis librae. & \raisebox{-1\normalbaselineskip}[0pt][0pt]{272°} & \raisebox{-1\normalbaselineskip}[0pt][0pt]{10′} & \raisebox{-1\normalbaselineskip}[0pt][0pt]{54°} & \raisebox{-1\normalbaselineskip}[0pt][0pt]{45′} & \raisebox{-1\normalbaselineskip}[0pt][0pt]{b} & \raisebox{-1\normalbaselineskip}[0pt][0pt]{4 parva} & \raisebox{-1\normalbaselineskip}[0pt][0pt]{λ} \\
6 & Borealis prior e duabus quae sunt in po-\newline\hspace*{4mm}steriore lance librae. & \raisebox{-1\normalbaselineskip}[0pt][0pt]{272} & \raisebox{-1\normalbaselineskip}[0pt][0pt]{10} & \raisebox{-1\normalbaselineskip}[0pt][0pt]{56} & \raisebox{-1\normalbaselineskip}[0pt][0pt]{10} & \raisebox{-1\normalbaselineskip}[0pt][0pt]{b} & \raisebox{-1\normalbaselineskip}[0pt][0pt]{3} & \raisebox{-1\normalbaselineskip}[0pt][0pt]{β} \\
7 & Borealis posterior e duabus quae sunt in\newline\hspace*{4mm}posteriore lance librae. & \raisebox{-1\normalbaselineskip}[0pt][0pt]{276} & \raisebox{-1\normalbaselineskip}[0pt][0pt]{20} & \raisebox{-1\normalbaselineskip}[0pt][0pt]{55} & \raisebox{-1\normalbaselineskip}[0pt][0pt]{20} & \raisebox{-1\normalbaselineskip}[0pt][0pt]{b} & \raisebox{-1\normalbaselineskip}[0pt][0pt]{3} & \raisebox{-1\normalbaselineskip}[0pt][0pt]{γ} \\
 & \multicolumn{7}{c||}{\parbox{150mm}{\centering\rule{0pt}{5mm}\textbf{Cygnus} (\textit{ad–daǵāǵah} « gallina »).}} & \\ \cline{2-8}
1 & Quae est in rostro. & 284 & 50 & 49 & 20 & b & 3 & β \\
2 & Quae est in medio colli. & 297 & 30 & 54 & 30 & b & 4 & η \\
3 & Quae est in pectore. & 309 & 40 & 57 & 20 & b & 3 & γ \\
4 & Stella lucida quae est in cauda. & 320 & 20 & 60 & 0 & b & 2 & α \\
5 & Stella quae est in cubito alae dexterae. & 300 & 30 & 64 & 40 & b & 3 & δ \\
6 & Media e tribus quae sunt in ala sinistra. & 302 & 20 & 71 & 30 & b & 4 magna & ι² \\
7 & Stella a septentrionibus eius, in extremitate\newline\hspace*{4mm}alae. & \raisebox{-1\normalbaselineskip}[0pt][0pt]{297} & \raisebox{-1\normalbaselineskip}[0pt][0pt]{50} & \raisebox{-1\normalbaselineskip}[0pt][0pt]{74} & \raisebox{-1\normalbaselineskip}[0pt][0pt]{0} & \raisebox{-1\normalbaselineskip}[0pt][0pt]{b} & \raisebox{-1\normalbaselineskip}[0pt][0pt]{4 magna} & \raisebox{-1\normalbaselineskip}[0pt][0pt]{κ} \\
8 & Quae est in cubito alae sinistrae. & 312 & 0 & 49 & 30 & b & 3 & ε \\
9 & Quae est in pede sinistro. & 321 & 10 & 55 & 10 & b & 4 magna & ν \\
10 & Quae est in genu sinistro. & 325 & 40 & 57 & 0 & b & 4 magna & ξ \\
 & \multicolumn{7}{c||}{\parbox{150mm}{\centering\rule{0pt}{5mm}\textbf{Cassiopea} (\textit{dhāt al–kursī} « [mulier] throno praedita »).}} & \\ \cline{2-8}
1 & Quae est in capite. & 19 & 0 & 45 & 20 & b & 3 & ζ \\
2 & Quae est in pectore. & 22 & 0 & 46 & 45 & b & 3 & α \\
\cline{2-8}\end{tabular}\end{center}
{\small
\par\hangindent=12mm\hangafter=1 5. Pt.: ὁ νοτιώτερος αὐτῶν (i. e. τῶν ἐν τῷ ζυγώματι ἑπομένων β) « australior earum (i. e. duarum sequentium « in zygomate lyrae) ». — Long. cod. ut H.; sed rectius B., L. et S. 275° 10′; qua re fortasse \textarabic{رعب} 272° error pro \textarabic{رعه} 275°. — Magnitudo ut L. et Pts. (H. et B. 4); S. et Arg. 5 parva.
\par\hangindent=12mm\hangafter=1 6. Pt.: τῶν ἐν τῷ ζυγώματι προηγουμένων β ὁ βορειότερος « borealior duarum praecedentium in zygomate lyrae ».— Lat. cod. \textarabic{م} 40′ pro \textarabic{ى} 10′.
\par\hangindent=12mm\hangafter=1 7. Pt.: τῶν ἐν τῷ ζυγώματι ἑπομένων β ὁ βορειότερος « borealior duarum sequentium in zygomate lyrae ».
\par\hspace*{8mm}\textbf{Cygnus}:
\par\hangindent=12mm\hangafter=1 1. Long. cod. \textarabic{رفز لـ} 287° 20′ pro \textarabic{رفد ن} 284° 50′, quae reposcit nr. 27 tabulae statuum fixarum ad annum 1211 (fol. 238,r.); Pt. et S. 285° 40′. — Lat. ut L. et S.; B. et H. 49° 0′.
\par\hangindent=12mm\hangafter=1 2. Magnitudo cod. \textarabic{ج} 3 pro \textarabic{د} 4; B. et H. 4; L. et Pts. 4 magna; S. 5; Arg. 4 parva.
\par\hangindent=12mm\hangafter=1 6. Lat. cod. \textarabic{عد ى} 74° 10′ pro \textarabic{عا ل} 71° 30′. — Magnitudo ut Pts. et L.; S. et Arg. 4.
\par\hangindent=12mm\hangafter=1 8. Cod. \textarabic{طرف} « extremitate » pro \textarabic{مرفق} « cubito » (ὁ ἐπὶ τοῦ ἀγκῶνος τῆς ἀριστερᾶς πτέρυγος).
\par\hangindent=12mm\hangafter=1 9. Magnitudo ut Pts. et L.; S. et Arg. 4.
\par\hangindent=12mm\hangafter=1 10. Long. cod. \textarabic{سكو لـ} 326° 20′ pro \textarabic{سكه م} 325° 40′. — Magnitudo ut Pts.; L., S. et Arg. 4.
\par\hspace*{8mm}\textbf{Cassiopea}:
\par\hangindent=12mm\hangafter=1 1. Magnitudo B., H. et Arg. 4; Pts., L. et S. 4 magna.
\par\hangindent=12mm\hangafter=1 2. Long. cod. \textarabic{نب} 52° pro \textarabic{كب} 22°.
\par}
\clearpage
\noindent\makebox[\textwidth]{\textbf{150}\hfill{\small C. A. NALLINO}\hfill\textbf{}}\par\vspace{-1mm}\noindent\rule{\textwidth}{.4pt}\par\vspace{2mm}
\begin{center}\begin{tabular}{r||p{78mm}|r@{\hspace{2.2mm}}r|r@{\hspace{2.2mm}}r|c|c||l}\cline{2-8}
 & \centering Stellarum descriptio. & \multicolumn{2}{c|}{Longitudo.} & \multicolumn{2}{c|}{Latitudo.} & \parbox{9mm}{\centering\small Plaga Caeli.} & Magnitudo. & \\ \cline{2-8}
3 & Quae est a septentrionibus eius, in zona. & 24° & 10′ & 47° & 50′ & b & 4 magna & η \\
4 & Stella quae est in genu. & 31 & 50 & 45 & 30 & b & 3 & δ \\
5 & Stella quae est in medio throni. & 19 & 0 & 51 & 40 & b & 3 & β \\
6 & Quae est supra pedem throni. & 26 & 10 & 52 & 40 & b & 4 parva & κ \\
 & \multicolumn{7}{c||}{\parbox{150mm}{\centering\rule{0pt}{5mm}\textbf{Perseus} (\textit{Firsāwus; al–fāris al–mumsik li–ra’s al–ghūl}\\« eques caput daemonis gestans »)}} & \\ \cline{2-8}
1 & Quae est in extremitate manus dexterae\newline\hspace*{4mm}. equitis. & 37 & 50 & 40 & 30 & b & nebulosa & χ \\
2 & Lucida quae est in latere dextero. & 46 & 0 & 30 & 0 & b & 2 & α \\
3 & Quae est in humero dextero. & 43 & 50 & 34 & 30 & b & 3 & γ \\
4 & Posterior ex tribus quae sunt in latere\newline\hspace*{4mm}dextero. & \raisebox{-1\normalbaselineskip}[0pt][0pt]{48} & \raisebox{-1\normalbaselineskip}[0pt][0pt]{50} & \raisebox{-1\normalbaselineskip}[0pt][0pt]{27} & \raisebox{-1\normalbaselineskip}[0pt][0pt]{20} & \raisebox{-1\normalbaselineskip}[0pt][0pt]{b} & \raisebox{-1\normalbaselineskip}[0pt][0pt]{4 magna} & \raisebox{-1\normalbaselineskip}[0pt][0pt]{δ} \\
5 & Quae est in femore sinistro. & 48 & 0 & 21 & 50 & b & 4 & ν \\
6 & Lucida ex iis quae sunt in capite daemonis. & 40 & 50 & 23 & 0 & b & 2 & β \\
7 & Stella quae est in talo sinistro. & 45 & 20 & 12 & 0 & b & 3 parva & ο \\
8 & Quae eam sequitur, in pede sinistro. & 47 & 30 & 11 & 0 & b & 3 & ζ \\
 & \multicolumn{7}{c||}{\parbox{150mm}{\centering\rule{0pt}{5mm}\textbf{Auriga} (\textit{dhū ’l–aʿinnah} « habenas tenens », idemque est \textit{al–ʿannāz al–ǵanūbī}\\« caprarius australis »).}} & \\ \cline{2-8}
1 & Stella quae est in capite aurigae. & 73 & 40 & 30 & 0 & b & 4 & δ \\
2 & Quae est in humero sinistro, et appellatur\newline\hspace*{4mm}\textit{al–ʿanz} « capra ». & \raisebox{-1\normalbaselineskip}[0pt][0pt]{66} & \raisebox{-1\normalbaselineskip}[0pt][0pt]{10} & \raisebox{-1\normalbaselineskip}[0pt][0pt]{22} & \raisebox{-1\normalbaselineskip}[0pt][0pt]{30} & \raisebox{-1\normalbaselineskip}[0pt][0pt]{b} & \raisebox{-1\normalbaselineskip}[0pt][0pt]{1} & \raisebox{-1\normalbaselineskip}[0pt][0pt]{α} \\
\cline{2-8}\end{tabular}\end{center}
{\small
\par\hangindent=12mm\hangafter=1 3. Long. cod. \textarabic{كج ن} 23° 50′ pro \textarabic{كد ى} 24° 10′. — Magnitudo ut Arg.; Pts., L et S. 4.
\par\hangindent=12mm\hangafter=1 5. Magnitudo cod. 4 parva; sed hoc manifeste pertinet ad stellam sequentem. Correxi iuxta Pt., S. et Arg.
\par\hangindent=12mm\hangafter=1 6. De magnitudine cfr. nr. 5; cod. « nebulosa » quod ad sequentem stellam pertinet. Correxi iuxta Pts., S. et Arg.; L. 4.
\par\hspace*{8mm}\textbf{Perseus}:
\par\hangindent=12mm\hangafter=1 1. In magnitudine cod. 4 parva, quod pertinet ad praecedentem stellam.
\par\hangindent=12mm\hangafter=1 2 et 3. Numeri magnitudinum in codice inter se permutati.
\par\hangindent=12mm\hangafter=1 4. Magnitudo Pt., S. et Arg. 3; sed emendare non audeo.
\par\hangindent=12mm\hangafter=1 5. Magnitudo cod. \textarabic{ب} 2 pro \textarabic{د} 4, ut habent B., S. et Arg.; L. 4 magna; H. 5.
\par\hangindent=12mm\hangafter=1 7. Long. cod. \textarabic{\AbjadZero{}} 0′ pro \textarabic{ل} 20′.
\par\hangindent=12mm\hangafter=1 8. Long. ut B., L. et S.; H. 40′. — Magnitudo ut Arg.; Pts. et S. 3 parva; L. 3 magna.
\par\hspace*{8mm}\textbf{Auriga}:
\par\hangindent=12mm\hangafter=1 1. Magnitudo cod. \textarabic{ج} 3 pro \textarabic{د} 4.
\par\hangindent=12mm\hangafter=1 2. Male in textu \textarabic{العيوق} \textit{al-ʿayyūq} (nomen Arabicum sideris illius) pro \textarabic{العير} cod. emendare proposui. Le-
\par}
\clearpage
\noindent\makebox[\textwidth]{\textbf{}\hfill{\small AL-BATTANI OPUS ASTRONOMICUM}\hfill\textbf{151}}\par\vspace{-1mm}\noindent\rule{\textwidth}{.4pt}\par\vspace{2mm}
\begin{center}\begin{tabular}{r||p{78mm}|r@{\hspace{2.2mm}}r|r@{\hspace{2.2mm}}r|c|c||l}\cline{2-8}
 & \centering Stellarum descriptio. & \multicolumn{2}{c|}{Longitudo.} & \multicolumn{2}{c|}{Latitudo.} & \parbox{9mm}{\centering\small Plaga Caeli.} & Magnitudo. & \\ \cline{2-8}
3 & Quae est in humero dextero. & 74° & 0′ & 20° & 0′ & b & 2 & β \\
4 & Stella quae est in malleolo sinistro. & 61 & 0 & 10 & 10 & b & 3 parva & ι \\
5 & Quae est in malleolo dextero. & 66 & 50 & 5 & 0 & b & 3 magna & γ \\
 & \multicolumn{7}{c||}{\parbox{150mm}{\centering\rule{0pt}{5mm}\textbf{Ophiuchus} (\textit{al–ḥawwā’ alladhī yamsik al–ḥayyah}\\« psyllus qui serpentem tenet »).}} & \\ \cline{2-8}
1 & Stella quae est in capite psylli. & 246 & 0 & 36 & 0 & b & 3 magna & α \\
2 & Praecedens e duabus quae sunt in humero\newline\hspace*{4mm}dextero. & \raisebox{-1\normalbaselineskip}[0pt][0pt]{249} & \raisebox{-1\normalbaselineskip}[0pt][0pt]{10} & \raisebox{-1\normalbaselineskip}[0pt][0pt]{27} & \raisebox{-1\normalbaselineskip}[0pt][0pt]{15} & \raisebox{-1\normalbaselineskip}[0pt][0pt]{b} & \raisebox{-1\normalbaselineskip}[0pt][0pt]{4} & \raisebox{-1\normalbaselineskip}[0pt][0pt]{β} \\
3 & Prima e duabus quae sunt in extremitate\newline\hspace*{4mm}manus sinistrae. & \raisebox{-1\normalbaselineskip}[0pt][0pt]{226} & \raisebox{-1\normalbaselineskip}[0pt][0pt]{10} & \raisebox{-1\normalbaselineskip}[0pt][0pt]{17} & \raisebox{-1\normalbaselineskip}[0pt][0pt]{0} & \raisebox{-1\normalbaselineskip}[0pt][0pt]{b} & \raisebox{-1\normalbaselineskip}[0pt][0pt]{3} & \raisebox{-1\normalbaselineskip}[0pt][0pt]{δ} \\
4 & Quae eam sequitur in extremitate manus. & 227 & 10 & 16 & 30 & b & 3 & ε \\
5 & Stella quae est in genu dextero. & 242 & 20 & 7 & 30 & b & 3 & η \\
6 & Stella quae est in crure dextero. & 244 & 50 & 2 & 15 & b & 3 magna & 40 Ophi \\
7 & Secunda e quatuor quae sunt in pede. & 245 & 30 & 1 & 30 & b & 4 magna &  \\
8 & Borealis quae est in linea trium borealium. & 233 & 20 & 11 & 50 & b & 3 magna & ζ \\
 & \multicolumn{7}{c||}{\parbox{150mm}{\centering\rule{0pt}{5mm}\textbf{Serpens} (\textit{al–ḥayyah allatī yamsik–hā al–ḥawwā’} « serpens quem psyllus tenet »).}} & \\ \cline{2-8}
1 & Quae est in tempore serpentis quem psyllus\newline\hspace*{4mm}tenet. & \raisebox{-1\normalbaselineskip}[0pt][0pt]{215} & \raisebox{-1\normalbaselineskip}[0pt][0pt]{30} & \raisebox{-1\normalbaselineskip}[0pt][0pt]{36} & \raisebox{-1\normalbaselineskip}[0pt][0pt]{0} & \raisebox{-1\normalbaselineskip}[0pt][0pt]{b} & \raisebox{-1\normalbaselineskip}[0pt][0pt]{3} & \raisebox{-1\normalbaselineskip}[0pt][0pt]{γ} \\
2 & Contigua naribus serpentis. & 212 & 50 & 40 & 0 & b & 4 & ρ \\
\cline{2-8}\end{tabular}\end{center}
{\small
\par\hangindent=12mm\hangafter=1 gendum erat \textarabic{العنز} \textit{al-ʿanz}, quae est versio nominis Ptolemaici αἴξ, ut cod. ipse praebet in nr. 3 tabulae statuum fixarum ad annum 1211. Cfr. \Name{Ṣūfī}, p. 93.
\par\hangindent=12mm\hangafter=1 5. Long. cod. \textarabic{صز ى} 67° 10′ pro \textarabic{صو ن} 66° 50′. — Magnitudo ut L. et Pts.; S. et Arg. 2.
\par\hspace*{8mm}\textbf{Ophiuchus}:
\par\hangindent=12mm\hangafter=1 1. Lat. cod. ut B., L. et S.; H. 36° 30′. — Magnitudo L., Pts. et S. 3; sed Arg. 2.
\par\hangindent=12mm\hangafter=1 2. Magnitudo ut Pts; L. 4 magna; S. 3 parva; Arg. 3.
\par\hangindent=12mm\hangafter=1 4. Magnitudo ut L.; Pts., S et Arg. 3 parva.
\par\hangindent=12mm\hangafter=1 5. Long. cod. \textarabic{م} 40′ pro \textarabic{لـ} 20′.
\par\hangindent=12mm\hangafter=1 6. Long. cod. ut L. et S.; minus bene B. et H. 247° 50′. — Magnitudo cod. ut L., sed vix recte; B. et H. 4; Pts. 4 magna; S. 4 parva vel 5 magna; Arg. 5.
\par\hangindent=12mm\hangafter=1 7. Lat. cod. \textarabic{له} 35′ pro \textarabic{ل} 30′; male H 4° 30′. — Plaga cod. ut L., codex quidam Graecus, et vetus translatio Arabica (apud \Name{Ṣūfī}, p. 102, adn. 2); sed melius B., H. et S. australis (tunc, iuxta \Name{Baily}, 42 Ophiuchi; iuxta \Name{Schjellerup}, ϑ Ophiuchi).
\par\hangindent=12mm\hangafter=1 8. Long. cod. \textarabic{رنج} 253° pro \textarabic{رلج} 233°. — Magnitudo ut Arg.; L, Pts. et S. 3.
\par\hspace*{8mm}\textbf{Serpens}:
\par\hangindent=12mm\hangafter=1 1. Long. ut B, L. et S.; male H. 212° 30′.
\par}
\clearpage
\noindent\makebox[\textwidth]{\textbf{152}\hfill{\small C. A. NALLINO}\hfill\textbf{}}\par\vspace{-1mm}\noindent\rule{\textwidth}{.4pt}\par\vspace{2mm}
\begin{center}\begin{tabular}{r||p{78mm}|r@{\hspace{2.2mm}}r|r@{\hspace{2.2mm}}r|c|c||l}\cline{2-8}
 & \centering Stellarum descriptio. & \multicolumn{2}{c|}{Longitudo.} & \multicolumn{2}{c|}{Latitudo.} & \parbox{9mm}{\centering\small Plaga Caeli.} & Magnitudo. & \\ \cline{2-8}
3 & Quae est prope initium colli. & 213° & 10′ & 34° & 15′ & b & 3 & β \\
4 & Quae est post vertebram anteriorem colli. & 212 & 50 & 29 & 15 & b & 3 & δ \\
5 & Media e tribus quae sunt post illam ver-\newline\hspace*{4mm}tebram. & \raisebox{-1\normalbaselineskip}[0pt][0pt]{215} & \raisebox{-1\normalbaselineskip}[0pt][0pt]{30} & \raisebox{-1\normalbaselineskip}[0pt][0pt]{25} & \raisebox{-1\normalbaselineskip}[0pt][0pt]{20} & \raisebox{-1\normalbaselineskip}[0pt][0pt]{b} & \raisebox{-1\normalbaselineskip}[0pt][0pt]{3} & \raisebox{-1\normalbaselineskip}[0pt][0pt]{α} \\
6 & Stella australis ex iis. & 217 & 30 & 24 & 0 & b & 3 magna & ε \\
7 & Stella australis quae est post femur psylli. & 244 & 50 & 10 & 30 & b & 4 magna & ν \\
8 & Quae sequitur primam e tribus in cauda\newline\hspace*{4mm}serpentis positis. & \raisebox{-1\normalbaselineskip}[0pt][0pt]{259} & \raisebox{-1\normalbaselineskip}[0pt][0pt]{50} & \raisebox{-1\normalbaselineskip}[0pt][0pt]{21} & \raisebox{-1\normalbaselineskip}[0pt][0pt]{15} & \raisebox{-1\normalbaselineskip}[0pt][0pt]{b} & \raisebox{-1\normalbaselineskip}[0pt][0pt]{4 magna} & \raisebox{-1\normalbaselineskip}[0pt][0pt]{η} \\
9 & Quae est in extremitate caudae serpentis. & 269 & 30 & 27 & 0 & b & 4 & θ \\
 & \multicolumn{7}{c||}{\parbox{150mm}{\centering\rule{0pt}{5mm}\textbf{Sagitta} (\textit{uwīsṭos} « ὀϊστός » id est \textit{an–naṣl} « cuspis sagittae »).}} & \\ \cline{2-8}
1 & Stella solitaria in sagitta. & 291 & 20 & 39 & 20 & b & 4 & γ \\
2 & Stella quae est in extremitate sagittae. & 284 & 30 & 38 & 40 & b & 5 & β \\
 & \multicolumn{7}{c||}{\parbox{150mm}{\centering\rule{0pt}{5mm}\textbf{Aquila} (\textit{an–nasr aṭ–ṭā’ir} « aquila volans »).}} & \\ \cline{2-8}
1 & Praecedens e duabus quae sunt in humero\newline\hspace*{4mm}sinistro. & \raisebox{-1\normalbaselineskip}[0pt][0pt]{284} & \raisebox{-1\normalbaselineskip}[0pt][0pt]{20} & \raisebox{-1\normalbaselineskip}[0pt][0pt]{31} & \raisebox{-1\normalbaselineskip}[0pt][0pt]{30} & \raisebox{-1\normalbaselineskip}[0pt][0pt]{b} & \raisebox{-1\normalbaselineskip}[0pt][0pt]{3} & \raisebox{-1\normalbaselineskip}[0pt][0pt]{γ} \\
2 & Quae sequitur stellam in medio capitis, et\newline\hspace*{4mm}est in collo. & \raisebox{-1\normalbaselineskip}[0pt][0pt]{286} & \raisebox{-1\normalbaselineskip}[0pt][0pt]{0} & \raisebox{-1\normalbaselineskip}[0pt][0pt]{27} & \raisebox{-1\normalbaselineskip}[0pt][0pt]{10} & \raisebox{-1\normalbaselineskip}[0pt][0pt]{b} & \raisebox{-1\normalbaselineskip}[0pt][0pt]{3} & \raisebox{-1\normalbaselineskip}[0pt][0pt]{β} \\
3 & \textit{an–nasr aṭ–ṭā’ir} « aquila volans », quae\newline\hspace*{4mm}est lucida inter humeros. & \raisebox{-1\normalbaselineskip}[0pt][0pt]{285} & \raisebox{-1\normalbaselineskip}[0pt][0pt]{0} & \raisebox{-1\normalbaselineskip}[0pt][0pt]{29} & \raisebox{-1\normalbaselineskip}[0pt][0pt]{10} & \raisebox{-1\normalbaselineskip}[0pt][0pt]{b} & \raisebox{-1\normalbaselineskip}[0pt][0pt]{2 magna} & \raisebox{-1\normalbaselineskip}[0pt][0pt]{α} \\
4 & Stella borealis propinqua \textit{an–nasr aṭ–ṭā’ir}. & 285 & 50 & 30 & 0 & b & 3 parva & ο \\
5 & Quae est sub \textit{an–nasr}, ab hac longe dis-\newline\hspace*{4mm}sita, in latere Viae lacteae. & \raisebox{-1\normalbaselineskip}[0pt][0pt]{273} & \raisebox{-1\normalbaselineskip}[0pt][0pt]{20} & \raisebox{-1\normalbaselineskip}[0pt][0pt]{36} & \raisebox{-1\normalbaselineskip}[0pt][0pt]{20} & \raisebox{-1\normalbaselineskip}[0pt][0pt]{b} & \raisebox{-1\normalbaselineskip}[0pt][0pt]{3} & \raisebox{-1\normalbaselineskip}[0pt][0pt]{ζ} \\
\cline{2-8}\end{tabular}\end{center}
{\small
\par\hangindent=12mm\hangafter=1 3. Lat. cod. \textarabic{ى} 10′ pro \textarabic{ٮه} = \textarabic{يه} 15′.
\par\hangindent=12mm\hangafter=1 4. Lat. cod. \textarabic{نط} 59° pro \textarabic{كط} 29°.
\par\hangindent=12mm\hangafter=1 6. Magnitudo L. et Pts. 3; S. et Arg. 3 parva.
\par\hangindent=12mm\hangafter=1 7. Long. cod. \textarabic{ى} 10′ pro \textarabic{ن} 50′; lat. cod. \textarabic{يز} 17° pro \textarabic{ي} 10°. — Magnitudo L, Pts. et S. 4; Arg. 4 parva.
\par\hangindent=12mm\hangafter=1 8. Minuta lat. ut S.; Pt. 10′.
\par\hspace*{8mm}\textbf{Sagitta}:
\par\hangindent=12mm\hangafter=1 1. Long. ut B., L. et S.; male H. 297° 10′.
\par\hangindent=12mm\hangafter=1 2. Lat. ut S.; L. 38° 45′; B. et H. 37° 40′. Magnitudo cod. \textarabic{د} 4 pro \textarabic{ه} 5.
\par\hspace*{8mm}\textbf{Aquila}:
\par\hangindent=12mm\hangafter=1 1. Long. cod. \textarabic{رضد} 294° pro \textarabic{رفد} 284°.
\par\hangindent=12mm\hangafter=1 2. Long. cod. \textarabic{رفد لـ} 284° 20′ pro \textarabic{رفو \AbjadZero{}} 286° 0′.
\par\hangindent=12mm\hangafter=1 3. Magnitudo cod. \textarabic{ج} 3 pro \textarabic{ب} 2.
\par\hangindent=12mm\hangafter=1 4. Long. cod. \textarabic{رفو ى} 286° 10′ pro \textarabic{رفه ن} 285° 50′.
\par}
\clearpage
\noindent\makebox[\textwidth]{\textbf{}\hfill{\small AL-BATTANI OPUS ASTRONOMICUM}\hfill\textbf{153}}\par\vspace{-1mm}\noindent\rule{\textwidth}{.4pt}\par\vspace{2mm}
\begin{center}\begin{tabular}{r||p{78mm}|r@{\hspace{2.2mm}}r|r@{\hspace{2.2mm}}r|c|c||l}\cline{2-8}
 & \centering Stellarum descriptio. & \multicolumn{2}{c|}{Longitudo.} & \multicolumn{2}{c|}{Latitudo.} & \parbox{9mm}{\centering\small Plaga Caeli.} & Magnitudo. & \\ \cline{2-8}
 & \multicolumn{7}{c||}{\parbox{150mm}{\centering\rule{0pt}{5mm}\textbf{Delphinus} (\textit{ad–dulfīn,} sive \textit{aṣ–ṣalīb} « crux »).}} & \\ \cline{2-8}
1 & Praecedens e tribus quae sunt in cauda. & 298 & 50 & 29 & 10 & b & 3 parva & ε \\
2 & Australis praecedens in primo latere. & 299 & 40 & 32 & 0 & b & 3 parva & β \\
3 & Contigua australis quae est in linea fulci-\newline\hspace*{4mm}menti rostri. & \raisebox{-1\normalbaselineskip}[0pt][0pt]{302} & \raisebox{-1\normalbaselineskip}[0pt][0pt]{30} & \raisebox{-1\normalbaselineskip}[0pt][0pt]{32} & \raisebox{-1\normalbaselineskip}[0pt][0pt]{0} & \raisebox{-1\normalbaselineskip}[0pt][0pt]{b} & \raisebox{-1\normalbaselineskip}[0pt][0pt]{3 parva} & \raisebox{-1\normalbaselineskip}[0pt][0pt]{δ} \\
4 & Borealis in latere posteriore. & 304 & 40 & 33 & 10 & b & 3 parva & γ \\
5 & Borealis in primo latere. & 301 & 20 & 33 & 50 & b & 3 parva & α \\
 & \multicolumn{7}{c||}{\parbox{150mm}{\centering\rule{0pt}{5mm}\textbf{Equuleus} (\textit{buruṭūmis} « [ἵππου] προτομή », sive \textit{al–faras} « equus »).}} & \\ \cline{2-8}
1 & Prima anterior e duabus quae sunt in ca-\newline\hspace*{4mm}pite \textit{buruṭūmis,} qui est equus. & \raisebox{-1\normalbaselineskip}[0pt][0pt]{307} & \raisebox{-1\normalbaselineskip}[0pt][0pt]{30} & \raisebox{-1\normalbaselineskip}[0pt][0pt]{20} & \raisebox{-1\normalbaselineskip}[0pt][0pt]{30} & \raisebox{-1\normalbaselineskip}[0pt][0pt]{b} & \raisebox{-1\normalbaselineskip}[0pt][0pt]{obscura} & \raisebox{-1\normalbaselineskip}[0pt][0pt]{α} \\
2 & Secunda, quae est e duabus his posterior. & 309 & 10 & 20 & 40 & b & obscura & β \\
3 & Praecedens e duabus quae sunt in ore. & 307 & 30 & 25 & 30 & b & obscura & γ \\
4 & Posterior ex iis. & 308 & 50 & 25 & 0 & b & obscura & δ \\
 & \multicolumn{7}{c||}{\parbox{150mm}{\centering\rule{0pt}{5mm}[\textbf{Pegasus}].}} & \\ \cline{2-8}
1 & Media ex iis quae sunt in capite Androme-\newline\hspace*{4mm}dae, quae est mulier maritum non ha-\newline\hspace*{4mm}bens. & \raisebox{-2\normalbaselineskip}[0pt][0pt]{358} & \raisebox{-2\normalbaselineskip}[0pt][0pt]{30} & \raisebox{-2\normalbaselineskip}[0pt][0pt]{26} & \raisebox{-2\normalbaselineskip}[0pt][0pt]{0} & \raisebox{-2\normalbaselineskip}[0pt][0pt]{b} & \raisebox{-2\normalbaselineskip}[0pt][0pt]{2 parva} & \raisebox{-2\normalbaselineskip}[0pt][0pt]{α Andr.} \\
2 & Quae est in dorso equi et in apice eius\newline\hspace*{4mm}humeri. & \raisebox{-1\normalbaselineskip}[0pt][0pt]{353} & \raisebox{-1\normalbaselineskip}[0pt][0pt]{20} & \raisebox{-1\normalbaselineskip}[0pt][0pt]{12} & \raisebox{-1\normalbaselineskip}[0pt][0pt]{30} & \raisebox{-1\normalbaselineskip}[0pt][0pt]{b} & \raisebox{-1\normalbaselineskip}[0pt][0pt]{2 parva} & \raisebox{-1\normalbaselineskip}[0pt][0pt]{γ} \\
\cline{2-8}\end{tabular}\end{center}
{\small
\par\hspace*{8mm}\textbf{Delphinus}:
\par\hangindent=12mm\hangafter=1 3. Pt., S. et \Name{Ulugh Beg}: « australis [duarum] in sequenti latere rhombi » (τῆς ἑπομένης τοῦ ῥόμβου πλευρᾶς ὁ νότιος). — Long. cod. \textarabic{ى} 10′ pro \textarabic{ل} 30′; lat cod. \textarabic{لج} 33° pro \textarabic{لب} 32°.
\par\hangindent=12mm\hangafter=1 4. Minuta long. ut L. et S.; B. et H. 20′.
\par\hangindent=12mm\hangafter=1 5. Minuta lat. ut B. et L.; H. 20′; S. 15′.
\par\hspace*{8mm}\textbf{Equuleus}:
\par\hangindent=12mm\hangafter=1 1. Male long. H. 307° 50′.
\par\hangindent=12mm\hangafter=1 2. Lat. cod. \textarabic{كه} 25° pro \textarabic{ك} 20°.
\par\hangindent=12mm\hangafter=1 4. Long. cod. \textarabic{\AbjadZero{}} 0′ pro \textarabic{ن} 50′. — Lat. cod. \textarabic{ل} 30′ pro \textarabic{\AbjadZero{}} 0′.
\par\hspace*{8mm}\textbf{Pegasus}:
\par\hangindent=12mm\hangafter=1 1. Long. cod. non amanuensium corruptela pro 359° 0′, quae Pt. et S. praebent, sed Albatenii correctio, [cuius numeri veritati sunt propinquiores quam Ptolemaici. — \textsc{Schiaparelli}].
\par\hangindent=12mm\hangafter=1 2. Pt., S. et \Name{Ulugh Beg}: « quae est in lumbo et in margine alae » (ὁ ἐπὶ τῆς ὀσφύος καὶ ἄκρου τοῦ πτεροῦ); versio Arabica Almagesti, quam al-Battānī prae oculis habuit, voci ἄκρον alterum sensum apicis (\textit{ra’s}) seu cacuminis tribuerat. Cfr. nr. 3 Andromedae.
\par}
\par\vfill\hfill 20
\clearpage
\noindent\makebox[\textwidth]{\textbf{154}\hfill{\small C. A. NALLINO}\hfill\textbf{}}\par\vspace{-1mm}\noindent\rule{\textwidth}{.4pt}\par\vspace{2mm}
\begin{center}\begin{tabular}{r||p{78mm}|r@{\hspace{2.2mm}}r|r@{\hspace{2.2mm}}r|c|c||l}\cline{2-8}
 & \centering Stellarum descriptio. & \multicolumn{2}{c|}{Longitudo.} & \multicolumn{2}{c|}{Latitudo.} & \parbox{9mm}{\centering\small Plaga Caeli.} & Magnitudo. & \\ \cline{2-8}
3 & \textit{Mankib al–faras} « humerus equi », quae\newline\hspace*{4mm}est in humero dextero ubi oritur pes. & \raisebox{-1\normalbaselineskip}[0pt][0pt]{343°} & \raisebox{-1\normalbaselineskip}[0pt][0pt]{20′} & \raisebox{-1\normalbaselineskip}[0pt][0pt]{31°} & \raisebox{-1\normalbaselineskip}[0pt][0pt]{0′} & \raisebox{-1\normalbaselineskip}[0pt][0pt]{b} & \raisebox{-1\normalbaselineskip}[0pt][0pt]{2 parva} & \raisebox{-1\normalbaselineskip}[0pt][0pt]{β} \\
4 & Quae est inter duos humeros in humero alae. & 337 & 50 & 19 & 40 & b & 2 parva & α \\
5 & Borealis e duabus quae sunt in genu dex-\newline\hspace*{4mm}tero. & \raisebox{-1\normalbaselineskip}[0pt][0pt]{340} & \raisebox{-1\normalbaselineskip}[0pt][0pt]{10} & \raisebox{-1\normalbaselineskip}[0pt][0pt]{35} & \raisebox{-1\normalbaselineskip}[0pt][0pt]{0} & \raisebox{-1\normalbaselineskip}[0pt][0pt]{b} & \raisebox{-1\normalbaselineskip}[0pt][0pt]{3} & \raisebox{-1\normalbaselineskip}[0pt][0pt]{η} \\
6 & Praecedens e duabus quae sunt in collo. & 330 & 0 & 18 & 0 & b & 3 & ζ \\
7 & Borealis e duabus quae sunt in capite. & 320 & 20 & 16 & 50 & b & 3 & θ \\
8 & Stella quae est in faucibus. & 316 & 30 & 22 & 30 & b & 3 magna & ε \\
9 & Quae est in malleolo dextero. & 334 & 50 & 41 & 10 & b & 4 magna & π² \\
10 & Quae est in genu sinistro. & 328 & 50 & 34 & 15 & b & 4 magna & ι \\
11 & Quae est in malleolo sinistro. & 323 & 30 & 36 & 50 & b & 4 magna & κ \\
 & \multicolumn{7}{c||}{\parbox{150mm}{\centering\rule{0pt}{5mm}\textbf{Andromeda} (\textit{Andrūmīdhis,} sive mulier quae maritum non vidit).}} & \\ \cline{2-8}
1 & Quae est inter duos humeros Andromedae. & 6 & 30 & 24 & 30 & b & 3 & δ \\
2 & Australis e tribus quae sunt supra eius\newline\hspace*{4mm}zonam. & \raisebox{-1\normalbaselineskip}[0pt][0pt]{15} & \raisebox{-1\normalbaselineskip}[0pt][0pt]{0} & \raisebox{-1\normalbaselineskip}[0pt][0pt]{26} & \raisebox{-1\normalbaselineskip}[0pt][0pt]{20} & \raisebox{-1\normalbaselineskip}[0pt][0pt]{b} & \raisebox{-1\normalbaselineskip}[0pt][0pt]{3} & \raisebox{-1\normalbaselineskip}[0pt][0pt]{β} \\
3 & Praecedens extra tres quae sunt in manu\newline\hspace*{4mm}dextera. & \raisebox{-1\normalbaselineskip}[0pt][0pt]{352} & \raisebox{-1\normalbaselineskip}[0pt][0pt]{50} & \raisebox{-1\normalbaselineskip}[0pt][0pt]{44} & \raisebox{-1\normalbaselineskip}[0pt][0pt]{0} & \raisebox{-1\normalbaselineskip}[0pt][0pt]{b} & \raisebox{-1\normalbaselineskip}[0pt][0pt]{3} & \raisebox{-1\normalbaselineskip}[0pt][0pt]{ο} \\
4 & Stella quae est supra pedem sinistrum. & 28 & 0 & 28 & 0 & b & 3 & γ \\
 & \multicolumn{7}{c||}{\parbox{150mm}{\centering\rule{0pt}{5mm}\textbf{Triangulum} (\textit{ṭrighānus} « τρίγονος » sive \textit{al–muthallath} « triangulum »).}} & \\ \cline{2-8}
1 & Stella quae est in summitate trianguli. & 22 & 10 & 16 & 30 & b & 3 & α \\
2 & Praecedens e tribus quae sunt in base. & 27 & 10 & 20 & 40 & b & 3 & β \\
3 & Postrema harum trium. & 28 & 0 & 19 & 0 & b & 3 & γ \\
\cline{2-8}\end{tabular}\end{center}
{\small
\par\hangindent=12mm\hangafter=1 7. Lat. cod. \textarabic{نو} 56° pro \textarabic{يو} 16°.
\par\hangindent=12mm\hangafter=1 8. Long. cod. \textarabic{سمو لج} 346° 33′ pro \textarabic{سٮو ل} = \textarabic{سيو ل} 316° 30′. — Magnitudo ut L.; Arg. 2 parva; Pts. et S. 3.
\par\hangindent=12mm\hangafter=1 9. Lat. cod. \textarabic{كا} 21° pro \textarabic{ما} 41°. — Magnitudo ut L. et Pts.; S. et Arg. 4.
\par\hangindent=12mm\hangafter=1 10. Long. cod. \textarabic{ل} 30′ pro \textarabic{ن} 50′. — Magnitudo ut L. et Pts.; S. et Arg. 4.
\par\hangindent=12mm\hangafter=1 11. Magnitudo ut Pts.; L., S. et Arg. 4.
\par\hspace*{8mm}\textbf{Andromeda}:
\par\hangindent=12mm\hangafter=1 2. Lat. ut L. et S.; B. et H. 25° 20′.
\par\hangindent=12mm\hangafter=1 3. Cod. « quae sunt in eius capite »; fortasse restituendum « quae sunt in apice [eius manus dexterae] »; nam Arabice \textit{ra’s} tam « caput » quam « apicem » significat. Graece ἐν τῷ δεξιῷ ἀκροχειρίῳ « in an- « teriore parte brachii dexteri »; ἄκρον = \textit{ra’s} ut in nr. 2 Pegasi. — Long. cod. \textarabic{سله} 335’ pro \textarabic{سنب} 352°.
\par\hspace*{8mm}\textbf{Triangulum}:
\par\hangindent=12mm\hangafter=1 1. Long. cod. \textarabic{فب} 82° pro \textarabic{كب} 22°.
\par\hangindent=12mm\hangafter=1 2. Long. cod. \textarabic{كج} 23° pro \textarabic{كز} 27°.
\par}
\clearpage
\noindent\makebox[\textwidth]{\textbf{}\hfill{\small AL-BATTANI OPUS ASTRONOMICUM}\hfill\textbf{155}}\par\vspace{-1mm}\noindent\rule{\textwidth}{.4pt}\par\vspace{2mm}
\begin{center}II.\end{center}
\begin{center}Constellationes zodiaci.\end{center}
\begin{center}\begin{tabular}{r||p{78mm}|r@{\hspace{2.2mm}}r|r@{\hspace{2.2mm}}r|c|c||l}\cline{2-8}
 & \centering Stellarum descriptio. & \multicolumn{2}{c|}{Longitudo.} & \multicolumn{2}{c|}{Latitudo.} & \parbox{9mm}{\centering\small Plaga Caeli.} & Magnitudo. & \\ \cline{2-8}
 & \multicolumn{7}{c||}{\parbox{150mm}{\centering\rule{0pt}{5mm}\textbf{Aries} (\textit{al–ḥamal}).}} & \\ \cline{2-8}
1 & Praecedens e duabus quae sunt in cornu. & 17° & 50′ & 7° & 20′ & b & 3 parva & γ \\
2 & Posterior e duabus iis. & 18 & 50 & 8 & 20 & b & 3 & β \\
3 & Borealis e duabus quae sunt in ore. & 22 & 10 & 7 & 40 & b & 5 & η \\
4 & Australis e duabus iis. & 22 & 40 & 6 & 0 & b & 5 & θ \\
5 & Stella quae est in collo Arietis. & 17 & 40 & 5 & 30 & b & 5 & ι \\
6 & Quae est in dorso. & 28 & \textit{40} & 6 & 0 & b & 6 & ν \\
7 & Quae est ubi oritur cauda. & 32 & 30 & 4 & 50 & b & 5 & ε \\
8 & Praecedens e tribus quae sunt in cauda. & 35 & 0 & 1 & 40 & b & 4 & δ \\
9 & Media ex iis. & 36 & 30 & 2 & 30 & b & 4 & ζ \\
10 & Postrema ex his tribus. & 38 & 10 & 1 & 50 & b & 4 & τ² \\
11 & Quae est post femur. & 30 & 50 & 1 & 10 & b & 5 & ρ¹ \\
12 & Quae est supra medium femoris. & 29 & 10 & 1 & 30 & austr. & 5 & σ \\
13 & Quae est in dorso pedis posterioris. & 28 & 10 & 5 & 15 & a & 4 \textit{magna} &  \\
 & \multicolumn{7}{l}{\hspace*{4mm}Et ex iis quae non sunt in figura:} & \\
14 & Quae est supra caput. & 21 & \textit{30} & 10 & 0 & b & 3 & α \\
 & \multicolumn{7}{c||}{\parbox{150mm}{\centering\rule{0pt}{5mm}\textbf{Taurus} (\textit{ath–thawr}).}} & \\ \cline{2-8}
1 & Borealis e quatuor quae sunt in sectione\newline\hspace*{4mm}Tauri. & \raisebox{-1\normalbaselineskip}[0pt][0pt]{37} & \raisebox{-1\normalbaselineskip}[0pt][0pt]{30} & \raisebox{-1\normalbaselineskip}[0pt][0pt]{6} & \raisebox{-1\normalbaselineskip}[0pt][0pt]{0} & \raisebox{-1\normalbaselineskip}[0pt][0pt]{a} & \raisebox{-1\normalbaselineskip}[0pt][0pt]{4} & \raisebox{-1\normalbaselineskip}[0pt][0pt]{\textit{f}} \\
\cline{2-8}\end{tabular}\end{center}
{\small
\par\hspace*{8mm}\textbf{Aries}:
\par\hangindent=12mm\hangafter=1 1. Long. cod. \textarabic{يد م} 14° 40′ pro \textarabic{يز ن} 17° 50′.
\par\hangindent=12mm\hangafter=1 2. Long. cod. \textarabic{م} 40′ pro \textarabic{ن} 50′.
\par\hangindent=12mm\hangafter=1 4. Long. cod. \textarabic{لـ} 20′ pro \textarabic{م} 40′.
\par\hangindent=12mm\hangafter=1 6. Long. cod. ut L., sed fortasse pro \textarabic{م} 40′ sunt \textarabic{ن} 50′ restituenda, ut B., H. et S. habent.
\par\hangindent=12mm\hangafter=1 7. Lat. cod. \textarabic{م} 40′ pro \textarabic{ن} 50′.
\par\hangindent=12mm\hangafter=1 8. Lat. cod. \textarabic{ن} 50′ pro \textarabic{م} 40′.
\par\hangindent=12mm\hangafter=1 11. Long. cod. \textarabic{ل} 30′ pro \textarabic{ن} 50′. — Minuta lat. ut L. et S.; B. et H. 30′.
\par\hangindent=12mm\hangafter=1 12. Lat. cod. \textarabic{ى} 10′ pro \textarabic{ل} 30′.
\par\hangindent=12mm\hangafter=1 13. Magnitudo ut L.; sed Pts. et S. 4, Arg. 5. Probabile est « magnam » pertinere ad stellam sequentem (Pts. et S. 3 magna; Arg. 2). — Iuxta \Name{Baily} est 87 Ceti; iuxta \Name{Delambre} et \Name{Schjellerup} μ Ceti.
\par\hangindent=12mm\hangafter=1 14. Long. ut H; sed probabiliter pro \textarabic{ل} 30′ restituenda \textarabic{ن} 50′ ut B., L. et S. habent. — Lat. ut L. et S; B. et H. 10° 30′. — De magnitudine cfr. nr. 13.
\par\hspace*{8mm}\textbf{Taurus}:
\par\hangindent=12mm\hangafter=1 1. Long. cod. \textarabic{لو} 36° pro \textarabic{لز} 37°.
\par}
\clearpage
\noindent\makebox[\textwidth]{\textbf{156}\hfill{\small C. A. NALLINO}\hfill\textbf{}}\par\vspace{-1mm}\noindent\rule{\textwidth}{.4pt}\par\vspace{2mm}
\begin{center}\begin{tabular}{r||p{78mm}|r@{\hspace{2.2mm}}r|r@{\hspace{2.2mm}}r|c|c||l}\cline{2-8}
 & \centering Stellarum descriptio. & \multicolumn{2}{c|}{Longitudo.} & \multicolumn{2}{c|}{Latitudo.} & \parbox{9mm}{\centering\small Plaga Caeli.} & Magnitudo. & \\ \cline{2-8}
2 & Secunda ex iis quatuor, quae sequitur. & 37° & 10′ & 7° & 15′ & a & 4 & \textit{s} \\
3 & Australis ex quatuor iis, quae sunt in se-\newline\hspace*{4mm}ctione. & \raisebox{-1\normalbaselineskip}[0pt][0pt]{35} & \raisebox{-1\normalbaselineskip}[0pt][0pt]{30} & \raisebox{-1\normalbaselineskip}[0pt][0pt]{9} & \raisebox{-1\normalbaselineskip}[0pt][0pt]{15} & \raisebox{-1\normalbaselineskip}[0pt][0pt]{a} & \raisebox{-1\normalbaselineskip}[0pt][0pt]{4} & \raisebox{-1\normalbaselineskip}[0pt][0pt]{ο} \\
4 & Stella quae est in pectore Tauri. & 44 & 50 & 8 & 0 & a & 3 & λ \\
5 & Quae est in humero dextero. & 40 & 50 & 9 & 30 & a & 5 & \textit{e} \\
6 & Quae est in vultu in naso, et est e stellis\newline\hspace*{4mm}\textit{ad–dabarān} appellatis. & \raisebox{-1\normalbaselineskip}[0pt][0pt]{50} & \raisebox{-1\normalbaselineskip}[0pt][0pt]{10} & \raisebox{-1\normalbaselineskip}[0pt][0pt]{5} & \raisebox{-1\normalbaselineskip}[0pt][0pt]{45} & \raisebox{-1\normalbaselineskip}[0pt][0pt]{a} & \raisebox{-1\normalbaselineskip}[0pt][0pt]{3 parva} & \raisebox{-1\normalbaselineskip}[0pt][0pt]{γ} \\
7 & Quae est inter hanc stellam et oculum bo-\newline\hspace*{4mm}realem. & \raisebox{-1\normalbaselineskip}[0pt][0pt]{51} & \raisebox{-1\normalbaselineskip}[0pt][0pt]{30} & \raisebox{-1\normalbaselineskip}[0pt][0pt]{4} & \raisebox{-1\normalbaselineskip}[0pt][0pt]{15} & \raisebox{-1\normalbaselineskip}[0pt][0pt]{a} & \raisebox{-1\normalbaselineskip}[0pt][0pt]{3 parva} & \raisebox{-1\normalbaselineskip}[0pt][0pt]{δ¹} \\
8 & Aldebaran (\textit{ad–dabarān} « sequens »), quae\newline\hspace*{4mm}est in oculo et sub cornu australi. & \raisebox{-1\normalbaselineskip}[0pt][0pt]{53} & \raisebox{-1\normalbaselineskip}[0pt][0pt]{50} & \raisebox{-1\normalbaselineskip}[0pt][0pt]{5} & \raisebox{-1\normalbaselineskip}[0pt][0pt]{10} & \raisebox{-1\normalbaselineskip}[0pt][0pt]{a} & \raisebox{-1\normalbaselineskip}[0pt][0pt]{1} & \raisebox{-1\normalbaselineskip}[0pt][0pt]{α} \\
9 & Quae est inter primam stellarum \textit{ad–daba-\newline\hspace*{4mm}rān} et oculum australem. & \raisebox{-1\normalbaselineskip}[0pt][0pt]{52} & \raisebox{-1\normalbaselineskip}[0pt][0pt]{0} & \raisebox{-1\normalbaselineskip}[0pt][0pt]{5} & \raisebox{-1\normalbaselineskip}[0pt][0pt]{50} & \raisebox{-1\normalbaselineskip}[0pt][0pt]{a} & \raisebox{-1\normalbaselineskip}[0pt][0pt]{3 parva} & \raisebox{-1\normalbaselineskip}[0pt][0pt]{θ¹} \\
10 & Sequens hanc, in oculo boreali. & 53 & 0 & 3 & 0 & a & 3 & ε \\
11 & Quae est in initio cornus, prope aurem au-\newline\hspace*{4mm}stralem. & \raisebox{-1\normalbaselineskip}[0pt][0pt]{58} & \raisebox{-1\normalbaselineskip}[0pt][0pt]{40} & \raisebox{-1\normalbaselineskip}[0pt][0pt]{4} & \raisebox{-1\normalbaselineskip}[0pt][0pt]{0} & \raisebox{-1\normalbaselineskip}[0pt][0pt]{a} & \raisebox{-1\normalbaselineskip}[0pt][0pt]{5} & \raisebox{-1\normalbaselineskip}[0pt][0pt]{\textit{i}} \\
12 & Quae est in medio cornu, ex duabus quae\newline\hspace*{4mm}sunt in cornu australi. & \raisebox{-1\normalbaselineskip}[0pt][0pt]{61} & \raisebox{-1\normalbaselineskip}[0pt][0pt]{30} & \raisebox{-1\normalbaselineskip}[0pt][0pt]{5} & \raisebox{-1\normalbaselineskip}[0pt][0pt]{0} & \raisebox{-1\normalbaselineskip}[0pt][0pt]{a} & \raisebox{-1\normalbaselineskip}[0pt][0pt]{5} & \raisebox{-1\normalbaselineskip}[0pt][0pt]{\textit{m}} \\
13 & Quae est in initio cornus borealis. & 56 & 50 & 4 & 0 & a & 4 & τ \\
14 & Quae est in extremitate cornus australis. & 68 & 50 & 2 & 30 & a & 3 & ζ \\
15 & Quae est in extremitate cornus borealis et\newline\hspace*{4mm}in pede Aurigae. & \raisebox{-1\normalbaselineskip}[0pt][0pt]{66} & \raisebox{-1\normalbaselineskip}[0pt][0pt]{50} & \raisebox{-1\normalbaselineskip}[0pt][0pt]{5} & \raisebox{-1\normalbaselineskip}[0pt][0pt]{0} & \raisebox{-1\normalbaselineskip}[0pt][0pt]{b} & \raisebox{-1\normalbaselineskip}[0pt][0pt]{3} & \raisebox{-1\normalbaselineskip}[0pt][0pt]{β} \\
\cline{2-8}\end{tabular}\end{center}
{\small
\par\hangindent=12mm\hangafter=1 2. Long. cod. \textarabic{لو ن} 36° 50′ pro \textarabic{لز ى} 37° 10′.
\par\hangindent=12mm\hangafter=1 4. Long. cod. \textarabic{مو ل} 46° 30′ pro \textarabic{مد ن} 44° 50′.
\par\hangindent=12mm\hangafter=1 5. Long. et lat. cod. \textarabic{ل} 30′ pro \textarabic{ن} 50′. — Magnitudo \textarabic{د} 4 pro \textarabic{ه} 5.
\par\hangindent=12mm\hangafter=1 6. Long. cod. \textarabic{\AbjadZero{}} 0′ pro \textarabic{ى} 10′.
\par\hangindent=12mm\hangafter=1 7. Lat. cod. \textarabic{ا} 1° pro \textarabic{د} 4°.
\par\hangindent=12mm\hangafter=1 8. Pt.: ὁ λαμπρὸς τῶν ὑάδων ἐπὶ τοῦ νοτίου ὀφθαλμοῦ ὑπόκιρρος « lucida Hyadum in oculo australi; subrufa ». Quae al-Battānī habet « et sub cornu » falsa sunt; fortasse versio Arabica Almagesti, qua ille usus est, pro ὑπόκιρρος « subrufa » in textu ὑπὸ κέρατος « sub cornu » legerat. Cfr. quae \Name{Ṣūfī} p. 37 de Albatenio narrat, et quae in adn. \textit{d} ad cap. LI dixi.
\par\hangindent=12mm\hangafter=1 9. Cod. « et alterum oculum australem ».
\par\hangindent=12mm\hangafter=1 10. Long. B., H. et S. 54°; sed L. 53° ut al-Battānī [quod veritatem magis appropinquat. \textsc{Schiaparelli}]. — Magnitudo cod. \textarabic{د} 4 pro \textarabic{ج} 3; L. 3; Pts. et S. 3 parva; Arg. 4 magna.
\par\hangindent=12mm\hangafter=1 11. Long. cod. \textarabic{نز} 57° pro \textarabic{نح} 58°. — Male H. lat. 0° 15′. — Magnitudo ut S.; Pt. 4.
\par\hangindent=12mm\hangafter=1 13. Latitudo falsa, sed e Graeco exemplari proveniens (δ 4° pro δ′ $\tfrac{1^\circ}{4}$), nam etiam apud L. et S.; S. autem addit: « latitudo, iuxta quod in Caelo conspicitur, debet esse 0° 0′ ». Vera Ptolemaei lectio 0° 15′, ut B. et H. praebent. De plaga \Name{Baily} dicit: « All the copies have the latitude \textit{south}, except L., where « it is described as \textit{north}, and which appears to be the most correct ».
\par\hangindent=12mm\hangafter=1 14. Lat. cod. \textarabic{ن} 50′ pro \textarabic{ل} 30′.
\par\hangindent=12mm\hangafter=1 15. Lat. cod. ut B., L. et S.; nam eadem est quae nr. 5 Aurigae; H. 6°. — Perperam magnitudo cod. « 3 « nebulosa ».
\par}
\clearpage
\noindent\makebox[\textwidth]{\textbf{}\hfill{\small AL-BATTANI OPUS ASTRONOMICUM}\hfill\textbf{157}}\par\vspace{-1mm}\noindent\rule{\textwidth}{.4pt}\par\vspace{2mm}
\begin{center}\begin{tabular}{r||p{78mm}|r@{\hspace{2.2mm}}r|r@{\hspace{2.2mm}}r|c|c||l}\cline{2-8}
 & \centering Stellarum descriptio. & \multicolumn{2}{c|}{Longitudo.} & \multicolumn{2}{c|}{Latitudo.} & \parbox{9mm}{\centering\small Plaga Caeli.} & Magnitudo. & \\ \cline{2-8}
16 & Borealis e Pleiadibus (\textit{ath–thurayyā}), quae\newline\hspace*{4mm}est in linea anteriore. & \raisebox{-1\normalbaselineskip}[0pt][0pt]{43°} & \raisebox{-1\normalbaselineskip}[0pt][0pt]{20′} & \raisebox{-1\normalbaselineskip}[0pt][0pt]{4°} & \raisebox{-1\normalbaselineskip}[0pt][0pt]{30′} & \raisebox{-1\normalbaselineskip}[0pt][0pt]{b} & \raisebox{-1\normalbaselineskip}[0pt][0pt]{5} &  \\
17 & Australis e Pleiadibus, quae est in extremi-\newline\hspace*{4mm}tate lateris anterioris. & \raisebox{-1\normalbaselineskip}[0pt][0pt]{43} & \raisebox{-1\normalbaselineskip}[0pt][0pt]{30} & \raisebox{-1\normalbaselineskip}[0pt][0pt]{3} & \raisebox{-1\normalbaselineskip}[0pt][0pt]{40} & \raisebox{-1\normalbaselineskip}[0pt][0pt]{b} & \raisebox{-1\normalbaselineskip}[0pt][0pt]{4} &  \\
18 & Parva posterior quae est in posteriore la-\newline\hspace*{4mm}tere Pleiadum. & \raisebox{-1\normalbaselineskip}[0pt][0pt]{44} & \raisebox{-1\normalbaselineskip}[0pt][0pt]{50} & \raisebox{-1\normalbaselineskip}[0pt][0pt]{3} & \raisebox{-1\normalbaselineskip}[0pt][0pt]{20} & \raisebox{-1\normalbaselineskip}[0pt][0pt]{b} & \raisebox{-1\normalbaselineskip}[0pt][0pt]{5} &  \\
19 & Parva exterior in parte boreali Pleiadum. & 44 & 50 & 5 & 0 & b & 4 &  \\
 & \multicolumn{7}{c||}{\parbox{150mm}{\centering\rule{0pt}{5mm}\textbf{Gemini} (\textit{at–taw’amāni}).}} & \\ \cline{2-8}
1 & Quae est in capite Gemini anterioris. & 94 & 30 & 9 & 30 & b & 2 & α \\
2 & Quae est in capite Gemini posterioris. & 97 & 50 & 6 & 15 & b & 2 & β \\
3 & Quae est ante pedem Gemini anterioris. & 77 & 40 & 1 & 30 & a & 4 magna & η \\
4 & Quae est in extremitate pedis Gemini an-\newline\hspace*{4mm}terioris. & \raisebox{-1\normalbaselineskip}[0pt][0pt]{81} & \raisebox{-1\normalbaselineskip}[0pt][0pt]{20} & \raisebox{-1\normalbaselineskip}[0pt][0pt]{3} & \raisebox{-1\normalbaselineskip}[0pt][0pt]{30} & \raisebox{-1\normalbaselineskip}[0pt][0pt]{a} & \raisebox{-1\normalbaselineskip}[0pt][0pt]{4 magna} & \raisebox{-1\normalbaselineskip}[0pt][0pt]{ν} \\
5 & Quae est in extremitate pedis sinistri Ge-\newline\hspace*{4mm}mini posterioris. & \raisebox{-1\normalbaselineskip}[0pt][0pt]{83} & \raisebox{-1\normalbaselineskip}[0pt][0pt]{10} & \raisebox{-1\normalbaselineskip}[0pt][0pt]{7} & \raisebox{-1\normalbaselineskip}[0pt][0pt]{30} & \raisebox{-1\normalbaselineskip}[0pt][0pt]{a} & \raisebox{-1\normalbaselineskip}[0pt][0pt]{3} & \raisebox{-1\normalbaselineskip}[0pt][0pt]{γ} \\
6 & Quae est in extremitate pedis dexteri Ge-\newline\hspace*{4mm}mini posterioris. & \raisebox{-1\normalbaselineskip}[0pt][0pt]{85} & \raisebox{-1\normalbaselineskip}[0pt][0pt]{50} & \raisebox{-1\normalbaselineskip}[0pt][0pt]{10} & \raisebox{-1\normalbaselineskip}[0pt][0pt]{30} & \raisebox{-1\normalbaselineskip}[0pt][0pt]{a} & \raisebox{-1\normalbaselineskip}[0pt][0pt]{4} & \raisebox{-1\normalbaselineskip}[0pt][0pt]{ξ} \\
7 & Quae est in sinistro genu Gemini poste-\newline\hspace*{4mm}rioris. & \raisebox{-1\normalbaselineskip}[0pt][0pt]{89} & \raisebox{-1\normalbaselineskip}[0pt][0pt]{\textit{40}} & \raisebox{-1\normalbaselineskip}[0pt][0pt]{2} & \raisebox{-1\normalbaselineskip}[0pt][0pt]{30} & \raisebox{-1\normalbaselineskip}[0pt][0pt]{a} & \raisebox{-1\normalbaselineskip}[0pt][0pt]{3} & \raisebox{-1\normalbaselineskip}[0pt][0pt]{ζ} \\
8 & Quae est in media flexura dextera huius\newline\hspace*{4mm}Gemini. & \raisebox{-1\normalbaselineskip}[0pt][0pt]{92} & \raisebox{-1\normalbaselineskip}[0pt][0pt]{\textit{50}} & \raisebox{-1\normalbaselineskip}[0pt][0pt]{6} & \raisebox{-1\normalbaselineskip}[0pt][0pt]{0} & \raisebox{-1\normalbaselineskip}[0pt][0pt]{a} & \raisebox{-1\normalbaselineskip}[0pt][0pt]{3} & \raisebox{-1\normalbaselineskip}[0pt][0pt]{λ} \\
9 & Quae est in latere sinistro huius Gemini. & 92 & \textit{50} & 0 & 30 & a & 3 & δ \\
10 & Quae est in genu sinistro Gemini ante-\newline\hspace*{4mm}rioris. & \raisebox{-1\normalbaselineskip}[0pt][0pt]{84} & \raisebox{-1\normalbaselineskip}[0pt][0pt]{10} & \raisebox{-1\normalbaselineskip}[0pt][0pt]{1} & \raisebox{-1\normalbaselineskip}[0pt][0pt]{30} & \raisebox{-1\normalbaselineskip}[0pt][0pt]{a} & \raisebox{-1\normalbaselineskip}[0pt][0pt]{4 magna} & \raisebox{-1\normalbaselineskip}[0pt][0pt]{ε} \\
\cline{2-8}\end{tabular}\end{center}
{\small
\par\hangindent=12mm\hangafter=1 16. Lat. cod. \textarabic{ه} 5° pro \textarabic{د} 4°. — Magnitudo cod. 3 nebulosa.
\par\hangindent=12mm\hangafter=1 17. Long. ut B.; H., L. et S. 40′. — Magnitudo cod. 4 nebulosa.
\par\hangindent=12mm\hangafter=1 18. Lat. ut S. et nonnulli codices Graeci; L. 5° 20′; B. et H. 3° 40′. — Magnitudo cod. 3 nebulosa.
\par\hangindent=12mm\hangafter=1 19. Magnitudo cod. 4 nebulosa.
\par\hspace*{8mm}\textbf{Gemini}:
\par\hangindent=12mm\hangafter=1 1. Lat. cod. \textarabic{لـ} 20′ pro \textarabic{ل} 30′; B., H. et S. 9° 30′; L. 9° 40′.
\par\hangindent=12mm\hangafter=1 4. Long. cod. \textarabic{قا} 101° pro \textarabic{فا} 81°. — Magnitudo ut Pts.; L. 4; S. 3 parva.
\par\hangindent=12mm\hangafter=1 5. Long. cod. \textarabic{قج} 103° pro \textarabic{فج} 83°.
\par\hangindent=12mm\hangafter=1 6. Long. cod. \textarabic{قه} 105° pro \textarabic{فه} 85°; male H. 82°.
\par\hangindent=12mm\hangafter=1 7. Long. cod. \textarabic{قط} 109° pro \textarabic{فط} 89°; minuta incerta sunt, nam B. et H. 25′, S. 10′, L. 20′.
\par\hangindent=12mm\hangafter=1 8. Long. cod. ut trium codicum Graecorum; B., H., L. et S. 92° 30′. — Lat. ut L. et S.; B. et H. 6° 10′.
\par\hangindent=12mm\hangafter=1 9. Long ut S.; sed pro \textarabic{ن} 50′ fortasse \textarabic{ل} 30′ restituenda, quae B., H. et L. praebent.
\par\hangindent=12mm\hangafter=1 10. Long. cod. \textarabic{قد} 104° pro \textarabic{فد} 84°. — Lat. cod. \textarabic{ى} 10′ pro \textarabic{ل} 30′. — Magnitudo L. 3; Pts. et S. 3 parva; Arg. 4 magna.
\par}
\clearpage
\noindent\makebox[\textwidth]{\textbf{158}\hfill{\small C. A. NALLINO}\hfill\textbf{}}\par\vspace{-1mm}\noindent\rule{\textwidth}{.4pt}\par\vspace{2mm}
\begin{center}\begin{tabular}{r||p{78mm}|r@{\hspace{2.2mm}}r|r@{\hspace{2.2mm}}r|c|c||l}\cline{2-8}
 & \centering Stellarum descriptio. & \multicolumn{2}{c|}{Longitudo.} & \multicolumn{2}{c|}{Latitudo.} & \parbox{9mm}{\centering\small Plaga Caeli.} & Magnitudo. & \\ \cline{2-8}
 & \multicolumn{7}{c||}{\parbox{150mm}{\centering\rule{0pt}{5mm}\textbf{Cancer} (\textit{as–saraṭān}).}} & \\ \cline{2-8}
1 & Medium praesepis, nebulosa quae est in pe-\newline\hspace*{4mm}ctore. & \raisebox{-1\normalbaselineskip}[0pt][0pt]{111°} & \raisebox{-1\normalbaselineskip}[0pt][0pt]{30′} & \raisebox{-1\normalbaselineskip}[0pt][0pt]{0°} & \raisebox{-1\normalbaselineskip}[0pt][0pt]{\textit{40}′} & \raisebox{-1\normalbaselineskip}[0pt][0pt]{b} & \raisebox{-1\normalbaselineskip}[0pt][0pt]{nebulosa} & \raisebox{-1\normalbaselineskip}[0pt][0pt]{ε} \\
2 & Borealis e duabus quae sunt prope quadri-\newline\hspace*{4mm}laterum nebulosum in tenebris. & \raisebox{-1\normalbaselineskip}[0pt][0pt]{108} & \raisebox{-1\normalbaselineskip}[0pt][0pt]{50} & \raisebox{-1\normalbaselineskip}[0pt][0pt]{1} & \raisebox{-1\normalbaselineskip}[0pt][0pt]{15} & \raisebox{-1\normalbaselineskip}[0pt][0pt]{b} & \raisebox{-1\normalbaselineskip}[0pt][0pt]{4 parva} & \raisebox{-1\normalbaselineskip}[0pt][0pt]{η} \\
3 & Australis e duabus iis, et appellantur qua-\newline\hspace*{4mm}drilaterum. & \raisebox{-1\normalbaselineskip}[0pt][0pt]{109} & \raisebox{-1\normalbaselineskip}[0pt][0pt]{10} & \raisebox{-1\normalbaselineskip}[0pt][0pt]{1} & \raisebox{-1\normalbaselineskip}[0pt][0pt]{10} & \raisebox{-1\normalbaselineskip}[0pt][0pt]{a} & \raisebox{-1\normalbaselineskip}[0pt][0pt]{4 parva} & \raisebox{-1\normalbaselineskip}[0pt][0pt]{θ} \\
4 & Quae est in chela boreal Cancri. & 109 & 30 & 11 & 50 & b & 4 magna & ι \\
5 & Quae est in chela australi. & 117 & 40 & 5 & 30 & a & 4 & α \\
6 & Quae est in extremitate pedis posterioris\newline\hspace*{4mm}australis. & \raisebox{-1\normalbaselineskip}[0pt][0pt]{108} & \raisebox{-1\normalbaselineskip}[0pt][0pt]{20} & \raisebox{-1\normalbaselineskip}[0pt][0pt]{7} & \raisebox{-1\normalbaselineskip}[0pt][0pt]{30} & \raisebox{-1\normalbaselineskip}[0pt][0pt]{a} & \raisebox{-1\normalbaselineskip}[0pt][0pt]{4} & \raisebox{-1\normalbaselineskip}[0pt][0pt]{β} \\
7 & Quae est in extremitate pedis posterioris\newline\hspace*{4mm}borealis. & \raisebox{-1\normalbaselineskip}[0pt][0pt]{103} & \raisebox{-1\normalbaselineskip}[0pt][0pt]{50} & \raisebox{-1\normalbaselineskip}[0pt][0pt]{1} & \raisebox{-1\normalbaselineskip}[0pt][0pt]{0} & \raisebox{-1\normalbaselineskip}[0pt][0pt]{b} & \raisebox{-1\normalbaselineskip}[0pt][0pt]{5} & \raisebox{-1\normalbaselineskip}[0pt][0pt]{μ²} \\
8 & \textit{Australis e duabus quas memoravimus.} & \textit{108} & \textit{20} & \textit{7} & \textit{30} & \textit{a} & \textit{4 magna} & β \\
9 & Borealis e duabus quae sunt prope nebu-\newline\hspace*{4mm}losam, et appellantur « duo asini » (\textit{al–\newline\hspace*{4mm}–ḥimārāni}). & \raisebox{-2\normalbaselineskip}[0pt][0pt]{111} & \raisebox{-2\normalbaselineskip}[0pt][0pt]{30} & \raisebox{-2\normalbaselineskip}[0pt][0pt]{2} & \raisebox{-2\normalbaselineskip}[0pt][0pt]{40} & \raisebox{-2\normalbaselineskip}[0pt][0pt]{b} & \raisebox{-2\normalbaselineskip}[0pt][0pt]{4 magna} & \raisebox{-2\normalbaselineskip}[0pt][0pt]{γ} \\
10 & Australis e duabus iis. & 112 & 30 & 0 & 10 & a & 4 & δ \\
 & \multicolumn{7}{c||}{\parbox{150mm}{\centering\rule{0pt}{5mm}\textbf{Leo} (\textit{al–asad}).}} & \\ \cline{2-8}
1 & Stella quae est in extremitate narium. & 119 & 30 & 10 & 0 & b & 4 & κ \\
2 & Borealis e duabus quae sunt in capite. & 125 & 30 & 12 & 0 & b & 3 & μ \\
3 & Media e tribus quae sunt in collo. & 133 & 20 & 8 & 30 & b & 2 & γ \\
4 & Praecedens, eaque borealis, e tribus quae\newline\hspace*{4mm}sunt in collo. & \raisebox{-1\normalbaselineskip}[0pt][0pt]{131} & \raisebox{-1\normalbaselineskip}[0pt][0pt]{20} & \raisebox{-1\normalbaselineskip}[0pt][0pt]{11} & \raisebox{-1\normalbaselineskip}[0pt][0pt]{0} & \raisebox{-1\normalbaselineskip}[0pt][0pt]{b} & \raisebox{-1\normalbaselineskip}[0pt][0pt]{3} & \raisebox{-1\normalbaselineskip}[0pt][0pt]{ζ} \\
\cline{2-8}\end{tabular}\end{center}
{\small
\par\hspace*{8mm}\textbf{Cancer}:
\par\hangindent=12mm\hangafter=1 1. Lat. ut L.; sed fortasse \textarabic{م} 40′ error cod pro \textarabic{لـ} 20′ quae B., H. et S. praebent.
\par\hangindent=12mm\hangafter=1 2. Rectius Pt. et S.: τοῦ περὶ τὸ νεφέλιον τετραπλεύρου τῶν προηγουμένων β ὁ βορειότατος « borealior duarum « praecedentium quadrilateri [quod est] circa nebulosam ». — Long. cod. \textarabic{قيا} 111° pro \textarabic{قح} 108°.
\par\hangindent=12mm\hangafter=1 3. Long. cod. \textarabic{قيب} 112° pro \textarabic{قط} 109°.
\par\hangindent=12mm\hangafter=1 4. Magnitudo L., Pts., S. et Arg. 4, sine « magna ».
\par\hangindent=12mm\hangafter=1 5. Long. cod. \textarabic{قيا} 111° pro \textarabic{قيز} 117°.
\par\hangindent=12mm\hangafter=1 6. Plaga cod. « borealis »; sed cfr. nr. 8. — Magnitudo cod. \textarabic{ه} 5 pro \textarabic{د} 4; L. et Arg. 4 magna; Pts. et S. 4 (cfr. nr. 8).
\par\hangindent=12mm\hangafter=1 7. Magnitudo cod. \textarabic{د} 4 pro \textarabic{ه} 5; ergo permutatio cum praecedenti.
\par\hangindent=12mm\hangafter=1 8. Est eadem stella quae sub nr. 6.
\par\hangindent=12mm\hangafter=1 9. Long. ut B., L. et S.; male H. 114° 10′. — Magnitudo ut L. et Pts.; S. 4; Arg. 4 parva.
\par\hspace*{8mm}\textbf{Leo}:
\par\hangindent=12mm\hangafter=1 2. Long. cod. \textarabic{\AbjadZero{}} 0′ pro \textarabic{ل} 30′.
\par\hangindent=12mm\hangafter=1 4. Long. cod. \textarabic{قله} 135° pro \textarabic{قلا} 131°. Magnitudo \textarabic{د} pro \textarabic{ـحـ} = \textarabic{ـجـ} 3.
\par}
\clearpage
\noindent\makebox[\textwidth]{\textbf{}\hfill{\small AL-BATTANI OPUS ASTRONOMICUM}\hfill\textbf{159}}\par\vspace{-1mm}\noindent\rule{\textwidth}{.4pt}\par\vspace{2mm}
\begin{center}\begin{tabular}{r||p{78mm}|r@{\hspace{2.2mm}}r|r@{\hspace{2.2mm}}r|c|c||l}\cline{2-8}
 & \centering Stellarum descriptio. & \multicolumn{2}{c|}{Longitudo.} & \multicolumn{2}{c|}{Latitudo.} & \parbox{9mm}{\centering\small Plaga Caeli.} & Magnitudo. & \\ \cline{2-8}
5 & Australis ex his tribus quae sunt in collo. & 131° & 50′ & 4° & 30′ & b & 3 & η \\
6 & \textit{Qalb al–asad} « cor Leonis », et vocatur\newline\hspace*{4mm}etiam \textit{al–malakī} (?) « regia ». & \raisebox{-1\normalbaselineskip}[0pt][0pt]{134} & \raisebox{-1\normalbaselineskip}[0pt][0pt]{0} & \raisebox{-1\normalbaselineskip}[0pt][0pt]{0} & \raisebox{-1\normalbaselineskip}[0pt][0pt]{10} & \raisebox{-1\normalbaselineskip}[0pt][0pt]{b} & \raisebox{-1\normalbaselineskip}[0pt][0pt]{1} & \raisebox{-1\normalbaselineskip}[0pt][0pt]{α} \\
7 & Praecedens e duabus quae sunt in dorso. & 142 & 30 & 12 & 15 & b & 5 & \textit{b} \\
8 & Stella quae hanc sequitur. & 145 & 20 & 13 & 40 & b & 2 parva & δ \\
9 & Quae est in femore posteriore. & 151 & 30 & 5 & 50 & b & 3 & ι \\
10 & Secunda quae etiam est in femore posteriore. & 152 & 50 & 1 & 15 & b & 4 & σ \\
11 & Quae etiam est in medio femoris posterioris. & 158 & 40 & 3 & 0 & a & 5 & υ \\
12 & \textit{aṣ–Ṣarfah} « conversio », stella quae est\newline\hspace*{4mm}in extremitate caudae. & \raisebox{-1\normalbaselineskip}[0pt][0pt]{155} & \raisebox{-1\normalbaselineskip}[0pt][0pt]{40} & \raisebox{-1\normalbaselineskip}[0pt][0pt]{11} & \raisebox{-1\normalbaselineskip}[0pt][0pt]{50} & \raisebox{-1\normalbaselineskip}[0pt][0pt]{b} & \raisebox{-1\normalbaselineskip}[0pt][0pt]{1} & \raisebox{-1\normalbaselineskip}[0pt][0pt]{β} \\
13 & Borealis e duabus quae sunt in postica parte\newline\hspace*{4mm}Leonis. & \raisebox{-1\normalbaselineskip}[0pt][0pt]{145} & \raisebox{-1\normalbaselineskip}[0pt][0pt]{30} & \raisebox{-1\normalbaselineskip}[0pt][0pt]{11} & \raisebox{-1\normalbaselineskip}[0pt][0pt]{30} & \raisebox{-1\normalbaselineskip}[0pt][0pt]{b} & \raisebox{-1\normalbaselineskip}[0pt][0pt]{5} & \raisebox{-1\normalbaselineskip}[0pt][0pt]{71 Leo.} \\
14 & Australis e duabus iis. & 147 & 30 & 9 & 40 & b & 3 & θ \\
15 & Stella quae est in dextero genu. & 128 & 30 & 0 & 0 & 0 & 5 & ψ \\
 & \multicolumn{7}{c||}{Stellae Comae, quae non sunt in figura Leonis:} & \\
16 & Prima earum \textit{Bulūqāmos} « πλόκαμος (co-\newline\hspace*{4mm}ma) », quae est stella inter caudam Leo-\newline\hspace*{4mm}nis et \textit{as–Simāk ar–rāmiḥ} (α Bootis). & \raisebox{-2\normalbaselineskip}[0pt][0pt]{156} & \raisebox{-2\normalbaselineskip}[0pt][0pt]{0} & \raisebox{-2\normalbaselineskip}[0pt][0pt]{30} & \raisebox{-2\normalbaselineskip}[0pt][0pt]{\textit{15}} & \raisebox{-2\normalbaselineskip}[0pt][0pt]{b} & \raisebox{-2\normalbaselineskip}[0pt][0pt]{conspicua} &  \\
17 & Praecedens magna quae est in crinibus ple-\newline\hspace*{4mm}xis et vocatur \textit{ʿurf al–asad} « iuba\newline\hspace*{4mm}Leonis ». & \raisebox{-2\normalbaselineskip}[0pt][0pt]{155} & \raisebox{-2\normalbaselineskip}[0pt][0pt]{30} & \raisebox{-2\normalbaselineskip}[0pt][0pt]{25} & \raisebox{-2\normalbaselineskip}[0pt][0pt]{0} & \raisebox{-2\normalbaselineskip}[0pt][0pt]{b} & \raisebox{-2\normalbaselineskip}[0pt][0pt]{obscura} & \raisebox{-2\normalbaselineskip}[0pt][0pt]{4 Comae} \\
18 & Stella quae eam sequitur in crinibus plexis. & 159 & 40 & 25 & 30 & b & obscura & 21 Comae \\
 & \multicolumn{7}{c||}{Hae tres appellantur \textit{adh–dhawā’ib} « comae ».} & \\
\cline{2-8}\end{tabular}\end{center}
{\small
\par\hangindent=12mm\hangafter=1 6. Cod. \textarabic{الماكن} \textit{al-mākin} « potestatem habens »; sed nomen mihi aliunde ignotum, qua re puto ab amanuensibus esse corruptum e notissimo \textarabic{الملكى} \textit{al–malakī} « [stella] regia » i. e. βασιλ΄σκος, Regulus. — Long. cod. est quam ipse al-Battānī observaverat (cfr. t. I, p. 124, cap. LI); Pt. et S 133° 40′.
\par\hangindent=12mm\hangafter=1 7. Magnitudo ut L. et Pts.; S. 5 magna; Arg. 4 parva; B. et H. 6.
\par\hangindent=12mm\hangafter=1 8. Long. cod. \textarabic{قمد} 144° pro \textarabic{قمه} 145°. — Magnitudo cod. \textarabic{د} 4 pro \textarabic{ٮ} = \textarabic{ب} 2; Arg. 2 parva, ut cod.; L., Pts. et S. 2.
\par\hangindent=12mm\hangafter=1 9. Long. cod. \textarabic{قنب} 152° pro \textarabic{قنا} 151°. — Lat. cod. \textarabic{د} 4° pro \textarabic{ه} 5°.
\par\hangindent=12mm\hangafter=1 11. Lat. ut L. et S.; H. 3° 10′; B. 3° 12′.
\par\hangindent=12mm\hangafter=1 12. Long. cod. \textarabic{قنه} 155° pro \textarabic{قند} 154°. — Lat. cod. \textarabic{ل} 30′ pro \textarabic{ن} 50′.
\par\hangindent=12mm\hangafter=1 13. Lat. cod. ut L.; B. et H. 10′; S. 20′. — \Name{Schjellerup} 72 Leonis, aliam stellam considerans.
\par\hangindent=12mm\hangafter=1 15. Long. cod. \textarabic{لـ} 20′ pro \textarabic{ل} 30′.
\par\hangindent=12mm\hangafter=1 16. Magnitudo in cod. definitur \textarabic{مشهور} « conspicua, nota, celebris, vulgata »; L. « luminosa »; versio Trapezuntii (1), e Graeco exemplari facta, « splendida »; e contra B et H. ἀμαυρός « obscura ». \Name{Ṣūfī}, qui eam quintae magnitudinis existimat, p. 154 ait eandem Ptolemaeum ex « obscuris » fecisse. — Lat. cod. paulo differt a Pt. et S. (30° 0′). — \Name{Baily}: « I cannot identify this star. ..... The lati- « tude accords with 31 Comae Berenices; but the longitude differs as much as 8° ». Iuxta \Name{Schjellerup} est Fl. 15 Comae.
\par}
\par\vspace{2mm}\begin{center}\rule{40mm}{.4pt}\end{center}{\small (1) \Name{Claudii Ptolemaei Pheludiensis Alexandrini} \textit{Almagestum, seu magnae constructionis mathematicae opus plane divinum, Latina donatum lingua ab Georgio Trapezuntio}, Venetiis 1528.}
\clearpage
\noindent\makebox[\textwidth]{\textbf{160}\hfill{\small C. A. NALLINO}\hfill\textbf{}}\par\vspace{-1mm}\noindent\rule{\textwidth}{.4pt}\par\vspace{2mm}
\begin{center}\begin{tabular}{r||p{78mm}|r@{\hspace{2.2mm}}r|r@{\hspace{2.2mm}}r|c|c||l}\cline{2-8}
 & \centering Stellarum descriptio. & \multicolumn{2}{c|}{Longitudo.} & \multicolumn{2}{c|}{Latitudo.} & \parbox{9mm}{\centering\small Plaga Caeli.} & Magnitudo. & \\ \cline{2-8}
 & \multicolumn{7}{c||}{\parbox{150mm}{\centering\rule{0pt}{5mm}\textbf{Virgo} (\textit{al–ʿadhrā’} « virgo » sive \textit{as–sunbulah} « spica »).}} & \\ \cline{2-8}
1 & Australis e duabus quae sunt in capite Vir-\newline\hspace*{4mm}ginis. & \raisebox{-1\normalbaselineskip}[0pt][0pt]{157°} & \raisebox{-1\normalbaselineskip}[0pt][0pt]{30′} & \raisebox{-1\normalbaselineskip}[0pt][0pt]{4°} & \raisebox{-1\normalbaselineskip}[0pt][0pt]{15′} & \raisebox{-1\normalbaselineskip}[0pt][0pt]{b} & \raisebox{-1\normalbaselineskip}[0pt][0pt]{5} & \raisebox{-1\normalbaselineskip}[0pt][0pt]{ν} \\
2 & Borealis e duabus iis. & 158 & 10 & 5 & 40 & b & 5 & ξ \\
3 & Quae est in dorso prope alam sinistram. & 160 & 10 & 0 & 10 & b & 3 & β \\
4 & Praecedens e quatuor quae sunt in ala si-\newline\hspace*{4mm}nistra. & \raisebox{-1\normalbaselineskip}[0pt][0pt]{169} & \raisebox{-1\normalbaselineskip}[0pt][0pt]{25} & \raisebox{-1\normalbaselineskip}[0pt][0pt]{1} & \raisebox{-1\normalbaselineskip}[0pt][0pt]{10} & \raisebox{-1\normalbaselineskip}[0pt][0pt]{b} & \raisebox{-1\normalbaselineskip}[0pt][0pt]{3} & \raisebox{-1\normalbaselineskip}[0pt][0pt]{η} \\
5 & Stella quae hanc sequitur. & 174 & 20 & 2 & 50 & b & 3 & γ \\
6 & Quae est in latere dextero sub mammis. & 175 & 30 & 8 & 30 & b & 3 & δ \\
7 & Postrema ex his quatuor memoratis. & 182 & 10 & 1 & 40 & b & 4 & θ \\
8 & Quae est in extremitate pedis sinistri. & 201 & 10 & 0 & 30 & b & 4 & λ \\
9 & Quae est in extremitate pedis dexteri bo-\newline\hspace*{4mm}realis. & \raisebox{-1\normalbaselineskip}[0pt][0pt]{203} & \raisebox{-1\normalbaselineskip}[0pt][0pt]{50} & \raisebox{-1\normalbaselineskip}[0pt][0pt]{9} & \raisebox{-1\normalbaselineskip}[0pt][0pt]{50} & \raisebox{-1\normalbaselineskip}[0pt][0pt]{b} & \raisebox{-1\normalbaselineskip}[0pt][0pt]{4} & \raisebox{-1\normalbaselineskip}[0pt][0pt]{μ} \\
10 & Quae est in cursu caudae vestis. & 197 & 50 & 7 & 30 & b & 4 & ι \\
11 & Quae est e tribus in ala boreali, [appel-\newline\hspace*{4mm}lata] \textit{al–mutaqaddim li-’l-qaṭāf} « pro-\newline\hspace*{4mm}« cedens ad vindemiandum ». & \raisebox{-2\normalbaselineskip}[0pt][0pt]{173} & \raisebox{-2\normalbaselineskip}[0pt][0pt]{20} & \raisebox{-2\normalbaselineskip}[0pt][0pt]{15} & \raisebox{-2\normalbaselineskip}[0pt][0pt]{10} & \raisebox{-2\normalbaselineskip}[0pt][0pt]{b} & \raisebox{-2\normalbaselineskip}[0pt][0pt]{3} & \raisebox{-2\normalbaselineskip}[0pt][0pt]{ε} \\
12 & Lucida quae est in extremitate manus sini-\newline\hspace*{4mm}strae, et appellatur \textit{as-sunbulah} « spica »\newline\hspace*{4mm}sive \textit{as–simāk al–aʿzal} « praeeminens\newline\hspace*{4mm}« inermis ». & \raisebox{-3\normalbaselineskip}[0pt][0pt]{187} & \raisebox{-3\normalbaselineskip}[0pt][0pt]{50} & \raisebox{-3\normalbaselineskip}[0pt][0pt]{2} & \raisebox{-3\normalbaselineskip}[0pt][0pt]{0} & \raisebox{-3\normalbaselineskip}[0pt][0pt]{a} & \raisebox{-3\normalbaselineskip}[0pt][0pt]{1} & \raisebox{-3\normalbaselineskip}[0pt][0pt]{α} \\
13 & Quae est in zona et nate dextera. & 186 & 0 & 8 & 40 & b & 3 & ζ \\
\cline{2-8}\end{tabular}\end{center}
{\small
\par\hspace*{8mm}\textbf{Virgo}:
\par\hangindent=12mm\hangafter=1 1. Long. ut L. et S.; minus bene B. et H. 156° 30′. — Male H. lat. 1° 15′.
\par\hangindent=12mm\hangafter=1 2. Lat. cod. \textarabic{ى} 10′ pro \textarabic{م} 40′.
\par\hangindent=12mm\hangafter=1 3. In cod. lat. 1° 10′ ex permutatione cum lat. stellae sequentis.
\par\hangindent=12mm\hangafter=1 4. Long. cod. \textarabic{قصا} 161° pro \textarabic{قصط} 169°. — Lat. cod. 0° 10′; cfr. nr. 3. — Magnitudo cod. \textarabic{د} 4 pro \textarabic{ـحـ} = \textarabic{ـجـ} 3.
\par\hangindent=12mm\hangafter=1 5. Long. cod. \textarabic{قعب} 172° pro \textarabic{قعد} 174°.
\par\hangindent=12mm\hangafter=1 6. Long. cod. \textarabic{قعب} 172° pro \textarabic{قعه} 175°.
\par\hangindent=12mm\hangafter=1 7. Long. cod. \textarabic{قعب} 172° pro \textarabic{قفب} 182°.
\par\hangindent=12mm\hangafter=1 9. Cod.: « Borealis quae est in extremitate pedis dexteri ». — Male H. long. 203° 30′, et lat. 0° 50′. — Magnitudo ut L. et Arg.; Pts. et S. 4 magna; B. et H. 3.
\par\hangindent=12mm\hangafter=1 10. Long. ut L. et S.; B. et H. 197° 30′.
\par\hangindent=12mm\hangafter=1 11. Long. cod. \textarabic{رج} 203° pro \textarabic{قعج} 173°. — Lat. cod. \textarabic{ل} 30′ pro \textarabic{ى} 10′, ut L. et S.; B. et H. 16° 0′. — Magnitudo ut L. et S.; Pts. 3 parva; Arg. 3 magna; male B. et H. 5.
\par\hangindent=12mm\hangafter=1 12. Long. cod. \textarabic{رز ل} 207° 30′ pro \textarabic{قفز ن} 187° 50′. — Plaga cod.: borealis (e litterarum permutatione cum stella sequenti).
\par\hangindent=12mm\hangafter=1 13. Long. cod. \textarabic{رو} 206° pro \textarabic{قفو} 186°. — Plaga cod.: australis (e permutatione cum plaga praecedentis stellae).
\par}
\clearpage
\noindent\makebox[\textwidth]{\textbf{}\hfill{\small AL-BATTANI OPUS ASTRONOMICUM}\hfill\textbf{161}}\par\vspace{-1mm}\noindent\rule{\textwidth}{.4pt}\par\vspace{2mm}
\begin{center}\begin{tabular}{r||p{78mm}|r@{\hspace{2.2mm}}r|r@{\hspace{2.2mm}}r|c|c||l}\cline{2-8}
 & \centering Stellarum descriptio. & \multicolumn{2}{c|}{Longitudo.} & \multicolumn{2}{c|}{Latitudo.} & \parbox{9mm}{\centering\small Plaga Caeli.} & Magnitudo. & \\ \cline{2-8}
 & \multicolumn{7}{c||}{\parbox{150mm}{\centering\rule{0pt}{5mm}\textbf{Libra} (\textit{al–mīzān}).}} & \\ \cline{2-8}
1 & Stella lucida e stellis lancis australis. & 209° & 10′ & 2° & 0′ & b & 2 & α \\
2 & Obscura a septentrionibus huius stellae. & 208 & 10 & 2 & 30 & b & 5 & μ \\
3 & Lucida ex iis quae sunt in extremitate lan-\newline\hspace*{4mm}cis borealis. & \raisebox{-1\normalbaselineskip}[0pt][0pt]{213} & \raisebox{-1\normalbaselineskip}[0pt][0pt]{20} & \raisebox{-1\normalbaselineskip}[0pt][0pt]{8} & \raisebox{-1\normalbaselineskip}[0pt][0pt]{20} & \raisebox{-1\normalbaselineskip}[0pt][0pt]{b} & \raisebox{-1\normalbaselineskip}[0pt][0pt]{2} & \raisebox{-1\normalbaselineskip}[0pt][0pt]{β} \\
4 & Media ex iis quae sunt in lance australi. & 215 & 10 & 1 & 40 & a & 4 & ι¹ \\
5 & Stella hanc praecedens. & 212 & 30 & 1 & 15 & b & 4 & ν¹ \\
6 & Stella media in lance boreali. & 219 & 0 & 4 & 45 & b & 4 & γ \\
 & \multicolumn{7}{c||}{\parbox{150mm}{\centering\rule{0pt}{5mm}\textbf{Scorpius} (\textit{al–ʿaqrab}).}} & \\ \cline{2-8}
1 & Borealis e tribus quae sunt inter oculos. & 227 & 40 & 1 & 20 & b & 3 & β \\
2 & Media ex his tribus. & 226 & 50 & 1 & 40 & a & 3 & δ \\
3 & Australis ex his tribus. & 226 & 50 & 5 & 0 & a & 3 & π \\
4 & Praecedens e tribus lucidis quae sunt in\newline\hspace*{4mm}pectore. & \raisebox{-1\normalbaselineskip}[0pt][0pt]{231} & \raisebox{-1\normalbaselineskip}[0pt][0pt]{30} & \raisebox{-1\normalbaselineskip}[0pt][0pt]{3} & \raisebox{-1\normalbaselineskip}[0pt][0pt]{45} & \raisebox{-1\normalbaselineskip}[0pt][0pt]{a} & \raisebox{-1\normalbaselineskip}[0pt][0pt]{3} & \raisebox{-1\normalbaselineskip}[0pt][0pt]{σ} \\
5 & \textit{Qalb al–ʿaqrab} « cor Scorpii », earum me-\newline\hspace*{4mm}dia, rubra. & \raisebox{-1\normalbaselineskip}[0pt][0pt]{233} & \raisebox{-1\normalbaselineskip}[0pt][0pt]{50} & \raisebox{-1\normalbaselineskip}[0pt][0pt]{4} & \raisebox{-1\normalbaselineskip}[0pt][0pt]{0} & \raisebox{-1\normalbaselineskip}[0pt][0pt]{a} & \raisebox{-1\normalbaselineskip}[0pt][0pt]{2} & \raisebox{-1\normalbaselineskip}[0pt][0pt]{α} \\
6 & Posterior ex his tribus. & 235 & 40 & 5 & 30 & a & 3 & τ \\
7 & Quae hanc sequitur in prima vertebra. & 239 & 40 & 11 & 0 & a & 3 & ε \\
8 & Quae est in secunda vertebra. & 240 & 0 & 15 & 0 & a & 3 & μ \\
9 & Borealis duplae quae est in tertia vertebra. & 241 & 10 & 18 & 40 & a & 4 & ζ¹ \\
10 & Quae eam sequitur in quarta vertebra. & 244 & 20 & 19 & 30 & a & 3 & η \\
11 & Eas sequens in quinta vertebra. & 249 & 20 & 18 & 50 & a & 3 & θ \\
\cline{2-8}\end{tabular}\end{center}
{\small
\par\hspace*{8mm}\textbf{Libra}:
\par\hangindent=12mm\hangafter=1 1. Lat. cod. inserviit quoque ad supputanda elementa nr. 45 tabulae statuum fixarum ad annum 1211; Pt. et S. 0° 40′.
\par\hangindent=12mm\hangafter=1 3. Long. cod. \textarabic{ه} 5′ pro \textarabic{لـ} 20′. — Lat. cod. confirmatur elementis nr. 46 tabulae statuum fixarum ad annum 1211; B., H. et S. 8° 50′; L. 8° 30′.
\par\hangindent=12mm\hangafter=1 4. Plaga cod.: borealis.
\par\hangindent=12mm\hangafter=1 5. Cod.: « sequens ».
\par\hspace*{8mm}\textbf{Scorpius}:
\par\hangindent=12mm\hangafter=1 1. Minuta long. cod. incerta; t. I, p. 124, cap. LI, ait se eam in 227° 45′ (Plato 20′) observavisse, quod comparatione cum Menelai longitudine confirmatur (cfr. I, 124, adn. 6). At Pt. et S. 227° 30′.
\par\hangindent=12mm\hangafter=1 2. Long. cod. \textarabic{ل} 30′ pro \textarabic{ن} 50′. — Lat. cod. \textarabic{لـ} 20′ pro \textarabic{م} 40′.
\par\hangindent=12mm\hangafter=1 3. Long. cod. \textarabic{ل} 30′ pro \textarabic{ن} 50′.
\par\hangindent=12mm\hangafter=1 5 Long. cod. \textarabic{ل} 30′ pro \textarabic{ن} 50′.
\par\hangindent=12mm\hangafter=1 6. Long. cod. \textarabic{ل} 30′ pro \textarabic{م} 40′.
\par\hangindent=12mm\hangafter=1 9. Cod.: « dupla borealis ..... »; Pt.: τοῦ ἐν τῷ τρίτῳ σπονδύλῳ διπλοῦ ὁ βόρειος.
\par\hangindent=12mm\hangafter=1 10. Long. cod. \textarabic{رمه} 245° pro \textarabic{رمد} 244°. — Magnitudo cod. \textarabic{د} 4 pro \textarabic{ـحـ} = \textarabic{ـجـ} 3.
\par}
\par\vfill\hfill 21
\clearpage
\noindent\makebox[\textwidth]{\textbf{162}\hfill{\small C. A. NALLINO}\hfill\textbf{}}\par\vspace{-1mm}\noindent\rule{\textwidth}{.4pt}\par\vspace{2mm}
\begin{center}\begin{tabular}{r||p{78mm}|r@{\hspace{2.2mm}}r|r@{\hspace{2.2mm}}r|c|c||l}\cline{2-8}
 & \centering Stellarum descriptio. & \multicolumn{2}{c|}{Longitudo.} & \multicolumn{2}{c|}{Latitudo.} & \parbox{9mm}{\centering\small Plaga Caeli.} & Magnitudo. & \\ \cline{2-8}
12 & Quae eas sequitur in sexta vertebra. & 251° & 40′ & 16° & 40′ & a & 3 & ι \\
13 & Quae est in septima vertebra prope acu-\newline\hspace*{4mm}leum. & \raisebox{-1\normalbaselineskip}[0pt][0pt]{250} & \raisebox{-1\normalbaselineskip}[0pt][0pt]{10} & \raisebox{-1\normalbaselineskip}[0pt][0pt]{15} & \raisebox{-1\normalbaselineskip}[0pt][0pt]{10} & \raisebox{-1\normalbaselineskip}[0pt][0pt]{a} & \raisebox{-1\normalbaselineskip}[0pt][0pt]{3} & \raisebox{-1\normalbaselineskip}[0pt][0pt]{ϰ} \\
14 & Posterior e duabus quae sunt in aculeo. & 248 & 40 & 13 & 20 & a & 3 & λ \\
15 & Praecedens e duabus iis. & 248 & 10 & 13 & 30 & a & \textit{4} & υ \\
 & \multicolumn{7}{l}{\hspace*{4mm}Et ex iis quae non sunt in figura:} & \\
16 & Borealis nebulosa quae aculeum sequitur. & 252 & 20 & 13 & 15 & a & nebulosa &  \\
 & \multicolumn{7}{c||}{\parbox{150mm}{\centering\rule{0pt}{5mm}\textbf{Sagittarius} (\textit{ar–rāmī} « iaculator » sive \textit{al–qaws} « arcus » et \textit{as–sahm}\\« sagitta » simul).}} & \\ \cline{2-8}
1 & Stella quae est in cuspide sagittae. & 255 & 40 & 6 & 20 & a & 3 & γ \\
2 & Quae est in initio manus sinistrae. & 258 & 50 & 6 & 30 & a & 3 & δ \\
3 & Quae est in parte australi arcus. & 259 & 10 & 10 & 50 & a & 3 & ε \\
4 & Australis e duabus quae sunt in parte bo-\newline\hspace*{4mm}reali arcus. & \raisebox{-1\normalbaselineskip}[0pt][0pt]{260} & \raisebox{-1\normalbaselineskip}[0pt][0pt]{10} & \raisebox{-1\normalbaselineskip}[0pt][0pt]{1} & \raisebox{-1\normalbaselineskip}[0pt][0pt]{30} & \raisebox{-1\normalbaselineskip}[0pt][0pt]{a} & \raisebox{-1\normalbaselineskip}[0pt][0pt]{3} & \raisebox{-1\normalbaselineskip}[0pt][0pt]{λ} \\
5 & Borealis e duabus iis, in extremitate arcus. & 257 & 50 & 2 & 50 & b & 3 & μ \\
6 & Quae est in humero Sagittarii. & 266 & 30 & 3 & 10 & a & 3 & σ \\
7 & Quae est ante hanc in sagitta. & 264 & 10 & 3 & 45 & a & 4 \textit{parva} & φ \\
8 & Dupla nebulosa quae est in oculo Sagittarii. & 266 & 20 & 0 & 45 & b & nebulosa & ν² \\
9 & Quae est in malleolo anteriore sinistro. & 268 & \textit{50} & 23 & 0 & a & 3 & β \\
\cline{2-8}\end{tabular}\end{center}
{\small
\par\hangindent=12mm\hangafter=1 15. Magnitudo cod. \textarabic{ـجـ} 3 pro \textarabic{د} 4; Pt. et Arg. 4; S. 3 parva.
\par\hangindent=12mm\hangafter=1 16. Long. cod. \textarabic{\AbjadZero{}} 0′ pro \textarabic{لـ} 20′. — Lat. cod. \textarabic{لـ} 20′ pro \textarabic{يه} 15′.
\par\hspace*{8mm}\textbf{Sagittarius}:
\par\hangindent=12mm\hangafter=1 1. Male H. long. 260° 40′. — Magnitudo cod.: « 3 nebulosa ».
\par\hangindent=12mm\hangafter=1 2. Long. cod. \textarabic{ل} 30′ pro \textarabic{ن} 50′.
\par\hangindent=12mm\hangafter=1 3. Lat. ut B., L. et S.; male H. 20° 20′. — Magnitudo cod. \textarabic{د} 4 pro \textarabic{ـحـ} = \textarabic{ـجـ} 3.
\par\hangindent=12mm\hangafter=1 6. Cod. plagam et magnitudinem cum nr. 8 permutavit (\textit{b}; nebulosa).
\par\hangindent=12mm\hangafter=1 7. Minuta lat. discrepant a B. et H. (30′), atque a L. et S. (50′). — Plaga cod. borealis. — Magnitudo « 4 parva » error videtur; L. 4; Pts., S. et Arg. 4 magna.
\par\hangindent=12mm\hangafter=1 8. Magnitudo cod. 3, e confusione cum nr. 6.
\par\hangindent=12mm\hangafter=1 9. Minuta long. in cod. \textarabic{ل} 30′ pro \textarabic{ن} 50′; correxi secundum Pt. et S., sed nescio an recte. Nam etiam in catalogi Albateniani exemplari quod \Name{Ṣūfī} prae oculis habuit minuta 30′ erant; qua re \Name{Ṣūfī}, p. 30, maligne et procul dubio iniuria dicit: « Il [i. e. al-Battānī] dit ensuite qu’il a observé les étoiles du « Sagittaire (1), et qu’il a trouvé l’étoile située sur le jarret de la jambe gauche de devant, dans « 28° 30′ du Sagittaire. Il apporte aussi dans son livre qu’il a trouvé, au temps de son observation, « que l’addition à faire à chaque étoile rapportée dans l’\textit{al-madjistī}, était de 11° 10′. Or, d’après ce « qu’il a dit sur l’addition à faire à chaque étoile, cette étoile devait être dans 28° 50′ du Sagittaire; « puisque son lieu rapporté dans l’\textit{al-madjistī} est dans 17° 40′, d’où il suit qu’il en a retranché 20
\par}
\par\vspace{2mm}\begin{center}\rule{40mm}{.4pt}\end{center}{\small (1) Falsum est; id nullo loco al–Battānī narrat.}
\clearpage
\noindent\makebox[\textwidth]{\textbf{}\hfill{\small AL-BATTANI OPUS ASTRONOMICUM}\hfill\textbf{163}}\par\vspace{-1mm}\noindent\rule{\textwidth}{.4pt}\par\vspace{2mm}
\begin{center}\begin{tabular}{r||p{78mm}|r@{\hspace{2.2mm}}r|r@{\hspace{2.2mm}}r|c|c||l}\cline{2-8}
 & \centering Stellarum descriptio. & \multicolumn{2}{c|}{Longitudo.} & \multicolumn{2}{c|}{Latitudo.} & \parbox{9mm}{\centering\small Plaga Caeli.} & Magnitudo. & \\ \cline{2-8}
10 & Media, in humero, e tribus quae sunt in\newline\hspace*{4mm}dorso. & \raisebox{-1\normalbaselineskip}[0pt][0pt]{268°} & \raisebox{-1\normalbaselineskip}[0pt][0pt]{50′} & \raisebox{-1\normalbaselineskip}[0pt][0pt]{4°} & \raisebox{-1\normalbaselineskip}[0pt][0pt]{30′} & \raisebox{-1\normalbaselineskip}[0pt][0pt]{a} & \raisebox{-1\normalbaselineskip}[0pt][0pt]{4} & \raisebox{-1\normalbaselineskip}[0pt][0pt]{τ} \\
11 & Sequens quae est sub axilla et est proxima\newline\hspace*{4mm}mediae ex iis quae sunt in dorso. & \raisebox{-1\normalbaselineskip}[0pt][0pt]{267} & \raisebox{-1\normalbaselineskip}[0pt][0pt]{30} & \raisebox{-1\normalbaselineskip}[0pt][0pt]{6} & \raisebox{-1\normalbaselineskip}[0pt][0pt]{45} & \raisebox{-1\normalbaselineskip}[0pt][0pt]{a} & \raisebox{-1\normalbaselineskip}[0pt][0pt]{3} & \raisebox{-1\normalbaselineskip}[0pt][0pt]{ζ} \\
12 & Quae est in genu sinistro. & 268 & 10 & 18 & 0 & a & 3 & α \\
13 & Quae est in malleolo pedis anterioris. & 257 & 50 & 13 & 0 & a & 3 & η \\
14 & Quae est in femore sinistro. & 278 & 30 & 13 & 30 & a & 3 & θ \\
15 & Quae est in crure dextero posteriore. & \textit{277} & \textit{30} & 20 & 10 & a & 3 & ι \\
16 & Borealis e quatuor quae sunt in initio cau-\newline\hspace*{4mm}dae; et dicitur \textit{ʿurqūb ar–rāmī} « mal-\newline\hspace*{4mm}leolus Sagittarii » [sic!]. & \raisebox{-2\normalbaselineskip}[0pt][0pt]{274} & \raisebox{-2\normalbaselineskip}[0pt][0pt]{0} & \raisebox{-2\normalbaselineskip}[0pt][0pt]{6} & \raisebox{-2\normalbaselineskip}[0pt][0pt]{30} & \raisebox{-2\normalbaselineskip}[0pt][0pt]{a} & \raisebox{-2\normalbaselineskip}[0pt][0pt]{1} & \raisebox{-2\normalbaselineskip}[0pt][0pt]{ω} \\
17 & Quae eam sequitur in latere boreali. & 280 & 0 & 4 & 50 & a & 5 & A \\
 & \multicolumn{7}{c||}{\parbox{150mm}{\centering\rule{0pt}{5mm}\textbf{Capricornus} (\textit{al–ǵady} « haedus »).}} & \\ \cline{2-8}
1 & Praecedens e tribus quae sunt in cornu po-\newline\hspace*{4mm}steriore. & \raisebox{-1\normalbaselineskip}[0pt][0pt]{288} & \raisebox{-1\normalbaselineskip}[0pt][0pt]{30} & \raisebox{-1\normalbaselineskip}[0pt][0pt]{7} & \raisebox{-1\normalbaselineskip}[0pt][0pt]{20} & \raisebox{-1\normalbaselineskip}[0pt][0pt]{b} & \raisebox{-1\normalbaselineskip}[0pt][0pt]{3} & \raisebox{-1\normalbaselineskip}[0pt][0pt]{α¹} \\
\cline{2-8}\end{tabular}\end{center}
{\small
\par\hangindent=12mm\hangafter=1 « pour faire croire qu’il a observé cette étoile (1). La preuve qu’il ne l’a ni observée ni connue, « et qu’elle n’a pas été plus connue des autres astronomes, depuis que l’on a composé des tables, « fait des sphères et dessiné dessus les étoiles, c’est qu’ils ont marqué dans leurs livres et sur leurs « sphères cette étoile comme de la deuxième grandeur, tandis qu’elle est des moindres de la qua- « trième ». Sed magnitudo cod. 3, quod verum est; B., H., L. et Pts. 2; fortasse Albatenii emendatio, ut in nr. 12.
\par\hangindent=12mm\hangafter=1 10. Long. cod. \textarabic{ل} 30′ pro \textarabic{ن} 50′. — Magnitudo cod. \textarabic{ـجـ} 3 fortasse Albatenii emendatio; L., Pts., S. et Arg. 4 magna.
\par\hangindent=12mm\hangafter=1 12. Magnitudo 3 videtur ab al-Battānī ipso provenire; \Name{Ṣūfī} p. 30: « Il a aussi trouvé dans l’\textit{al-madjistī} « que l’étoile au genou de la même jambe de devant était des petites de la deuxième grandeur, et « néanmoins il l’a marquée de la troisième. C’est l’étoile qui se trouve très-près de la 6ᵉ étoile de « la Couronne, parce qu’elles ont la même longitude; l’étoile du genou est seulement plus méridio- « nale de 50′. Elles sont toutes les deux de la quatrième grandeur. Peut-être les traducteurs ou les « copistes ont-ils marqué la grandeur de ces deux étoiles par un \textarabic{د} \textit{dâl}, qu’un copiste a pris pour « un \textarabic{ب} \textit{bâ}, et ainsi les a marqué de la deuxième grandeur ». Tamen non amanuenses, sed Ptolemaeus ipse erraverat et stellam duxerat secundae magnitudinis; L. 2 parva.
\par\hangindent=12mm\hangafter=1 13. Long. cod. \textarabic{ل} 30′ pro \textarabic{ن} 50′.
\par\hangindent=12mm\hangafter=1 15. Long. cod. incerta; B. et H. 274° 0′; L. 277° 50′; S. 278° 0′.
\par\hangindent=12mm\hangafter=1 16. \textit{ʿUrqūb ar-rāmī} « malleolus Sagittarii » apud reliquos Arabes est β Sagittarii, scilicet nr. 9 Albatenii. — Long. cod. est quae, iuxta B. et H., ad nr. 15 pertineret; lat. et magnitudo quoque valde a Ptolemaicis differunt. Exspectandum erat: long. 278° 30′ (L. 279° 50′, S. 278° 50′), lat. 4° 50′, magnitudo 5. Sed numeris falsis codicis al–Battānī supputavit elementa nr. 48 tabulae statuum fixarum ad annum 1211.
\par\hangindent=12mm\hangafter=1 17. Long. cod. \textarabic{لـ} 20′ pro \textarabic{\AbjadZero{}} 0′.
\par\hspace*{8mm}\textbf{Capricornus}:
\par\hangindent=12mm\hangafter=1 1. Cod. « praecedens », sed rectius « borealis ».
\par}
\par\vspace{2mm}\begin{center}\rule{40mm}{.4pt}\end{center}{\small (1) Sed cfr. quae dixi in adn. \textit{e} ad LI caput.}
\clearpage
\noindent\makebox[\textwidth]{\textbf{164}\hfill{\small C. A. NALLINO}\hfill\textbf{}}\par\vspace{-1mm}\noindent\rule{\textwidth}{.4pt}\par\vspace{2mm}
\begin{center}\begin{tabular}{r||p{78mm}|r@{\hspace{2.2mm}}r|r@{\hspace{2.2mm}}r|c|c||l}\cline{2-8}
 & \centering Stellarum descriptio. & \multicolumn{2}{c|}{Longitudo.} & \multicolumn{2}{c|}{Latitudo.} & \parbox{9mm}{\centering\small Plaga Caeli.} & Magnitudo. & \\ \cline{2-8}
2 & Media e tribus his. & 288° & 50′ & 6° & 40′ & b & 6 & ν \\
3 & Australis ex his tribus memoratis. & 288 & 30 & 5 & 0 & b & 3 & β \\
4 & Australis e tribus quae sunt in ore. & 290 & 10 & 0 & 45 & b & 6 & ο \\
5 & Praecedens duarum reliquarum ex tribus\newline\hspace*{4mm}quae sunt in ore. & \raisebox{-1\normalbaselineskip}[0pt][0pt]{289} & \raisebox{-1\normalbaselineskip}[0pt][0pt]{50} & \raisebox{-1\normalbaselineskip}[0pt][0pt]{1} & \raisebox{-1\normalbaselineskip}[0pt][0pt]{45} & \raisebox{-1\normalbaselineskip}[0pt][0pt]{b} & \raisebox{-1\normalbaselineskip}[0pt][0pt]{6} & \raisebox{-1\normalbaselineskip}[0pt][0pt]{π} \\
6 & Tertia stella quae eas duas sequitur in ore. & 290 & 0 & 1 & 30 & b & 6 & ρ \\
7 & Quae est sub genu dextero. & 292 & 0 & 6 & 30 & a & 4 & ψ \\
8 & Quae est in genu sinistro. & 292 & 50 & 8 & 40 & a & 4 & ω \\
9 & Praecedens ex duabus sibi invicem proximis\newline\hspace*{4mm}quae sunt in ventre. & \raisebox{-1\normalbaselineskip}[0pt][0pt]{301} & \raisebox{-1\normalbaselineskip}[0pt][0pt]{20} & \raisebox{-1\normalbaselineskip}[0pt][0pt]{6} & \raisebox{-1\normalbaselineskip}[0pt][0pt]{50} & \raisebox{-1\normalbaselineskip}[0pt][0pt]{a} & \raisebox{-1\normalbaselineskip}[0pt][0pt]{4} & \raisebox{-1\normalbaselineskip}[0pt][0pt]{ζ} \\
10 & Stella hanc sequens in ventre. & 301 & 30 & 6 & 0 & a & 5 &  \\
11 & Praecedens e duabus quae sunt prope\newline\hspace*{4mm}caudam. & \raisebox{-1\normalbaselineskip}[0pt][0pt]{305} & \raisebox{-1\normalbaselineskip}[0pt][0pt]{0} & \raisebox{-1\normalbaselineskip}[0pt][0pt]{2} & \raisebox{-1\normalbaselineskip}[0pt][0pt]{10} & \raisebox{-1\normalbaselineskip}[0pt][0pt]{a} & \raisebox{-1\normalbaselineskip}[0pt][0pt]{3} & \raisebox{-1\normalbaselineskip}[0pt][0pt]{γ} \\
12 & Stella quae eam sequitur. & 307 & 30 & 2 & 0 & a & 3 & δ \\
13 & Australis e tribus his reliquis. & 309 & 50 & 0 & 0 & 0 & 5 & μ \\
14 & Media e tribus iis. & 308 & 50 & 2 & 50 & b & 5 & λ \\
15 & Borealis ex iis quae sunt in extremitate\newline\hspace*{4mm}caudae. & 309 & 50 & 4 & 20 & b & 5 & \textit{c}¹ \\
16 & Praecedens e duabus quae sunt in spina dorsi. & 304 & 30 & 4 & 45 & a & 4 & ε \\
17 & Praecedens e duabus quae sunt in dorso. & 297 & 50 & 0 & 0 & a & 4 & θ \\
18 & Australis e quatuor quae sunt in parte bo-\newline\hspace*{4mm}reali caudae. & \raisebox{-1\normalbaselineskip}[0pt][0pt]{303} & \raisebox{-1\normalbaselineskip}[0pt][0pt]{0} & \raisebox{-1\normalbaselineskip}[0pt][0pt]{0} & \raisebox{-1\normalbaselineskip}[0pt][0pt]{0} & \raisebox{-1\normalbaselineskip}[0pt][0pt]{a} & \raisebox{-1\normalbaselineskip}[0pt][0pt]{5} & \raisebox{-1\normalbaselineskip}[0pt][0pt]{μ} \\
 & \multicolumn{7}{c||}{\parbox{150mm}{\centering\rule{0pt}{5mm}\textbf{Aquarius} (\textit{as–sāqī} « aquator » sive \textit{ad–dalw} « situla »).}} & \\ \cline{2-8}
1 & Stella quae est in capitê Aquarii. & 311 & 30 & 15 & 45 & b & 5 & \textit{d} \\
2 & Lucida e duabus quae sunt in humero\newline\hspace*{4mm}dextero. & \raisebox{-1\normalbaselineskip}[0pt][0pt]{317} & \raisebox{-1\normalbaselineskip}[0pt][0pt]{30} & \raisebox{-1\normalbaselineskip}[0pt][0pt]{11} & \raisebox{-1\normalbaselineskip}[0pt][0pt]{0} & \raisebox{-1\normalbaselineskip}[0pt][0pt]{b} & \raisebox{-1\normalbaselineskip}[0pt][0pt]{3} & \raisebox{-1\normalbaselineskip}[0pt][0pt]{α} \\
\cline{2-8}\end{tabular}\end{center}
{\small
\par\hangindent=12mm\hangafter=1 2. Long. cod. \textarabic{ل} 30′ pro \textarabic{ن} 50′.
\par\hangindent=12mm\hangafter=1 4. Lat. cod. \textarabic{مه مه} 45° 45′ pro \textarabic{\AbjadZero{} مه} 0° 45′.
\par\hangindent=12mm\hangafter=1 7. Male lat. H. 6° 20′. — Plaga cod. borealis.
\par\hangindent=12mm\hangafter=1 10. Iuxta \Name{Delambre} et \Name{Schjellerup} \textit{b} Capricorni; iuxta \Name{Baily} Fl. 35 Capricorni.
\par\hangindent=12mm\hangafter=1 11. Long. cod. \textarabic{سط} 309°, quod in scriptura Arabico-maghrebina facile oritur ex \textarabic{سه} 305°. Hunc numerum igitur recepi, qui etiam cum observationibus hodiernis, \Name{Schiaparelli} teste, melius quam alii congruit. B., L. et S. 306°; H. 303°. — Lat. cod. \textarabic{ل} 30′ pro \textarabic{ى} 10′.
\par\hangindent=12mm\hangafter=1 12-14. Numeri cod. sine erroribus, sed nomina stellarum inter se commutata. Nam duodecima stella in cod. est: « Praecedens e quatuor quae sunt in parte boreali caudae », quod est sidus in hoc catalogo Albateniano omissum; item cod. nomen stellae 14 loco 13, 15 loco 14, 12 loco 15 ponit, et 13 omittit. — In long. nr. 13 male H. 301° 50′.
\par\hangindent=12mm\hangafter=1 18. Est mera repetitio nr. 13, mutata long. iuxta falsam lectionem quam etiam H. praebet. — Lat. cod. \textarabic{ج} 3° pro \textarabic{\AbjadZero{}} 0°. — Magnitudo cod. \textarabic{د} 4 pro \textarabic{ه} 5.
\par\hspace*{8mm}\textbf{Aquarius}:
\par\hangindent=12mm\hangafter=1 1. Lat. cod. \textarabic{ه} 5° pro \textarabic{يه} 15°.
\par}
\clearpage
\noindent\makebox[\textwidth]{\textbf{}\hfill{\small AL-BATTANI OPUS ASTRONOMICUM}\hfill\textbf{165}}\par\vspace{-1mm}\noindent\rule{\textwidth}{.4pt}\par\vspace{2mm}
\begin{center}\begin{tabular}{r||p{78mm}|r@{\hspace{2.2mm}}r|r@{\hspace{2.2mm}}r|c|c||l}\cline{2-8}
 & \centering Stellarum descriptio. & \multicolumn{2}{c|}{Longitudo.} & \multicolumn{2}{c|}{Latitudo.} & \parbox{9mm}{\centering\small Plaga Caeli.} & Magnitudo. & \\ \cline{2-8}
3 & Obscura quae est sub ea. & 316° & 20′ & 9° & 40′ & b & 5 & ο \\
4 & Quae est in humero sinistro. & 307 & 40 & 8 & 50 & b & 3 & β \\
5 & Stella sequens e tribus quae sunt in manu\newline\hspace*{4mm}sinistra. & \raisebox{-1\normalbaselineskip}[0pt][0pt]{298} & \raisebox{-1\normalbaselineskip}[0pt][0pt]{50} & \raisebox{-1\normalbaselineskip}[0pt][0pt]{5} & \raisebox{-1\normalbaselineskip}[0pt][0pt]{30} & \raisebox{-1\normalbaselineskip}[0pt][0pt]{b} & \raisebox{-1\normalbaselineskip}[0pt][0pt]{3} & \raisebox{-1\normalbaselineskip}[0pt][0pt]{ν} \\
6 & Media e tribus his. & 297 & 20 & 8 & 0 & b & 4 & μ \\
7 & Praecedens ex his tribus. & 295 & 50 & 8 & 40 & b & 3 & ε \\
8 & Quae est in brachio dextero. & 320 & 40 & 8 & 45 & b & 3 & γ \\
9 & Borealis e tribus quae sunt in extremitate\newline\hspace*{4mm}manus. & \raisebox{-1\normalbaselineskip}[0pt][0pt]{322} & \raisebox{-1\normalbaselineskip}[0pt][0pt]{50} & \raisebox{-1\normalbaselineskip}[0pt][0pt]{10} & \raisebox{-1\normalbaselineskip}[0pt][0pt]{45} & \raisebox{-1\normalbaselineskip}[0pt][0pt]{b} & \raisebox{-1\normalbaselineskip}[0pt][0pt]{3} & \raisebox{-1\normalbaselineskip}[0pt][0pt]{π} \\
10 & Praecedens e duabus australibus reliquis. & 323 & 10 & 9 & 0 & b & 3 & ζ \\
11 & Stella hanc sequens. & 324 & 30 & 8 & 30 & b & 3 & η \\
12 & Australis e duabus quae sunt in crure dextero. & 322 & 50 & 7 & 30 & a & 3 & δ \\
13 & Praecedens e quinque quae sunt in loco exi-\newline\hspace*{4mm}tus aquae. & \raisebox{-1\normalbaselineskip}[0pt][0pt]{326} & \raisebox{-1\normalbaselineskip}[0pt][0pt]{10} & \raisebox{-1\normalbaselineskip}[0pt][0pt]{2} & \raisebox{-1\normalbaselineskip}[0pt][0pt]{0} & \raisebox{-1\normalbaselineskip}[0pt][0pt]{b} & \raisebox{-1\normalbaselineskip}[0pt][0pt]{4} &  \\
14 & Quae eam sequitur ad austrum. & 326 & 0 & 0 & 10 & b & 4 & λ \\
15 & Quae hanc sequitur post ansam. & 328 & 50 & 1 & 10 & a & 4 & \textit{h}¹ \\
16 & Quae eam etiam sequitur. & 331 & 10 & 0 & 30 & a & 4 & φ \\
17 & Australis ex iis quae sunt in ansa. & 331 & 40 & 1 & 40 & a & 4 & χ \\
18 & Postrema earum quae sunt in loco finis\newline\hspace*{4mm}aquae, et appellatur \textit{fam al–ḥūt al–ǵa-\newline\hspace*{4mm}nūbī} « os piscis australis ». & \raisebox{-2\normalbaselineskip}[0pt][0pt]{318} & \raisebox{-2\normalbaselineskip}[0pt][0pt]{10} & \raisebox{-2\normalbaselineskip}[0pt][0pt]{20} & \raisebox{-2\normalbaselineskip}[0pt][0pt]{20} & \raisebox{-2\normalbaselineskip}[0pt][0pt]{a} & \raisebox{-2\normalbaselineskip}[0pt][0pt]{1} & \raisebox{-2\normalbaselineskip}[0pt][0pt]{α Piscis Austr.} \\
19 & Borealis e duabus quae sunt ab austro stel-\newline\hspace*{4mm}larum ansae. & \raisebox{-1\normalbaselineskip}[0pt][0pt]{330} & \raisebox{-1\normalbaselineskip}[0pt][0pt]{10} & \raisebox{-1\normalbaselineskip}[0pt][0pt]{3} & \raisebox{-1\normalbaselineskip}[0pt][0pt]{30} & \raisebox{-1\normalbaselineskip}[0pt][0pt]{a} & \raisebox{-1\normalbaselineskip}[0pt][0pt]{4} & \raisebox{-1\normalbaselineskip}[0pt][0pt]{ψ¹} \\
20 & Praecedens e duabus sibi invicem proximis\newline\hspace*{4mm}quae sunt in crure [dextero]. & \raisebox{-1\normalbaselineskip}[0pt][0pt]{\textit{322}} & \raisebox{-1\normalbaselineskip}[0pt][0pt]{\textit{50}} & \raisebox{-1\normalbaselineskip}[0pt][0pt]{7} & \raisebox{-1\normalbaselineskip}[0pt][0pt]{\textit{30}} & \raisebox{-1\normalbaselineskip}[0pt][0pt]{b} & \raisebox{-1\normalbaselineskip}[0pt][0pt]{\textit{3}} & \raisebox{-1\normalbaselineskip}[0pt][0pt]{θ} \\
21 & Quae est in nate dextera Aquarii. & 319 & 50 & 0 & 50 & a & 5 & σ \\
\cline{2-8}\end{tabular}\end{center}
{\small
\par\hangindent=12mm\hangafter=1 3. Lat. cod. \textarabic{ن} 50′ pro \textarabic{م} 40′.
\par\hangindent=12mm\hangafter=1 4. Long. cod. \textarabic{لـ} 20′ pro \textarabic{م} 40′.
\par\hangindent=12mm\hangafter=1 5. Cod. « sequens tres .....», omissa scilicet particula \textarabic{من} « ex ».
\par\hangindent=12mm\hangafter=1 7. Lat. cod. \textarabic{لـ} 20′ pro \textarabic{م} 40′.
\par\hangindent=12mm\hangafter=1 8. Cod. \textarabic{فخذ} « femore » pro \textarabic{عضد} « brachio ». — Lat. cod. \textarabic{ر} 7′ pro \textarabic{مه} 45′.
\par\hangindent=12mm\hangafter=1 9. Cod. \textarabic{ذنبه} « caudae » pro \textarabic{يده} « manus ».
\par\hangindent=12mm\hangafter=1 12. Plaga cod. borealis.
\par\hangindent=12mm\hangafter=1 13. \Name{Baily}: « This star is probably [Flamsteed] 67 Aquarii; but the longitude should be increased nearly « 2° ». Item \Name{Schjellerup}: Fl. 67.
\par\hangindent=12mm\hangafter=1 15. Lat. ut B. et H.; L. et S. 1° 0′.
\par\hangindent=12mm\hangafter=1 16. Long. cod. \textarabic{م} 40′ pro \textarabic{ى} 10′. — Lat. cod. \textarabic{ـجـ \AbjadZero{}} 3° 0′ pro \textarabic{\AbjadZero{} ل} 0° 30′.
\par\hangindent=12mm\hangafter=1 18. Lat cod. ut B.; male H., L. et S. 23° 0′ (ϰγ pro ϰγ′).
\par\hangindent=12mm\hangafter=1 20. Plaga excepta, numeri omnes falsi; ad nr. 12 pertinent. Veri erant: 317° 20′, 3° 0′, b, 4.
\par\hangindent=12mm\hangafter=1 21. Intra nr. 20 et 21 cod. perperam denuo stellam 19 inserit. — Lat. cod. 3° 20′, e confusione cum lat. stellae 19 falso in cod. repetitae. — Magnitudo cod. ut L.; Pts. 4; S. 4 parva; Arg. 5 magna.
\par}
\clearpage
\noindent\makebox[\textwidth]{\textbf{166}\hfill{\small C. A. NALLINO}\hfill\textbf{}}\par\vspace{-1mm}\noindent\rule{\textwidth}{.4pt}\par\vspace{2mm}
\begin{center}\begin{tabular}{r||p{78mm}|r@{\hspace{2.2mm}}r|r@{\hspace{2.2mm}}r|c|c||l}\cline{2-8}
 & \centering Stellarum descriptio. & \multicolumn{2}{c|}{Longitudo.} & \multicolumn{2}{c|}{Latitudo.} & \parbox{9mm}{\centering\small Plaga Caeli.} & Magnitudo. & \\ \cline{2-8}
22 & Media e tribus quae sunt in prima curva-\newline\hspace*{4mm}tura ansae. & \raisebox{-1\normalbaselineskip}[0pt][0pt]{333°} & \raisebox{-1\normalbaselineskip}[0pt][0pt]{20′} & \raisebox{-1\normalbaselineskip}[0pt][0pt]{14°} & \raisebox{-1\normalbaselineskip}[0pt][0pt]{45′} & \raisebox{-1\normalbaselineskip}[0pt][0pt]{a} & \raisebox{-1\normalbaselineskip}[0pt][0pt]{5} & \raisebox{-1\normalbaselineskip}[0pt][0pt]{A²} \\
23 & Prima e tribus quae sunt in secunda cur-\newline\hspace*{4mm}vatura ansae. & \raisebox{-1\normalbaselineskip}[0pt][0pt]{323} & \raisebox{-1\normalbaselineskip}[0pt][0pt]{0} & \raisebox{-1\normalbaselineskip}[0pt][0pt]{14} & \raisebox{-1\normalbaselineskip}[0pt][0pt]{45} & \raisebox{-1\normalbaselineskip}[0pt][0pt]{a} & \raisebox{-1\normalbaselineskip}[0pt][0pt]{4} & \raisebox{-1\normalbaselineskip}[0pt][0pt]{\textit{c}¹} \\
24 & Australis e duabus quae sunt in nate sinistra. & 312 & 50 & 1 & 40 & a & 4 & ι \\
25 & Media e tribus aliis quae sunt in prima\newline\hspace*{4mm}curvatura ansae. & \raisebox{-1\normalbaselineskip}[0pt][0pt]{328} & \raisebox{-1\normalbaselineskip}[0pt][0pt]{40} & \raisebox{-1\normalbaselineskip}[0pt][0pt]{15} & \raisebox{-1\normalbaselineskip}[0pt][0pt]{0} & \raisebox{-1\normalbaselineskip}[0pt][0pt]{a} & \raisebox{-1\normalbaselineskip}[0pt][0pt]{4} & \raisebox{-1\normalbaselineskip}[0pt][0pt]{\textit{b}²} \\
26 & Borealis e duabus quae sunt in crure sini-\newline\hspace*{4mm}stro sub genu. & \raisebox{-1\normalbaselineskip}[0pt][0pt]{319} & \raisebox{-1\normalbaselineskip}[0pt][0pt]{0} & \raisebox{-1\normalbaselineskip}[0pt][0pt]{9} & \raisebox{-1\normalbaselineskip}[0pt][0pt]{0} & \raisebox{-1\normalbaselineskip}[0pt][0pt]{a} & \raisebox{-1\normalbaselineskip}[0pt][0pt]{5} & \raisebox{-1\normalbaselineskip}[0pt][0pt]{\textit{g}¹} \\
 & \multicolumn{7}{l}{\hspace*{4mm}Et ex iis quae non sunt in figura:} & \\
27 & Praecedens e tribus quae sequuntur ansam\newline\hspace*{4mm}amphorae. & \raisebox{-1\normalbaselineskip}[0pt][0pt]{337} & \raisebox{-1\normalbaselineskip}[0pt][0pt]{50} & \raisebox{-1\normalbaselineskip}[0pt][0pt]{15} & \raisebox{-1\normalbaselineskip}[0pt][0pt]{30} & \raisebox{-1\normalbaselineskip}[0pt][0pt]{a} & \raisebox{-1\normalbaselineskip}[0pt][0pt]{4 magna} & \raisebox{-1\normalbaselineskip}[0pt][0pt]{2 Ceti.} \\
28 & Borealis e duabus reliquis. & 340 & 50 & 14 & 40 & a & 4 magna & 6 Ceti. \\
29 & Stella eas sequens post ansam amphorae. & 340 & 10 & 18 & 15 & a & 4 magna & 7 Ceti. \\
 & \multicolumn{7}{c||}{\parbox{150mm}{\centering\rule{0pt}{5mm}\textbf{Pisces} (\textit{as–samakatāni} « duo pisces »).}} & \\ \cline{2-8}
1 & Quae est in ore Piscis anterioris sive au-\newline\hspace*{4mm}stralis. & \raisebox{-1\normalbaselineskip}[0pt][0pt]{332} & \raisebox{-1\normalbaselineskip}[0pt][0pt]{50} & \raisebox{-1\normalbaselineskip}[0pt][0pt]{9} & \raisebox{-1\normalbaselineskip}[0pt][0pt]{15} & \raisebox{-1\normalbaselineskip}[0pt][0pt]{b} & \raisebox{-1\normalbaselineskip}[0pt][0pt]{4 magna} & \raisebox{-1\normalbaselineskip}[0pt][0pt]{β} \\
2 & Australis e duabus quae sunt in capite hu-\newline\hspace*{4mm}ius Piscis. & \raisebox{-1\normalbaselineskip}[0pt][0pt]{335} & \raisebox{-1\normalbaselineskip}[0pt][0pt]{20} & \raisebox{-1\normalbaselineskip}[0pt][0pt]{7} & \raisebox{-1\normalbaselineskip}[0pt][0pt]{30} & \raisebox{-1\normalbaselineskip}[0pt][0pt]{b} & \raisebox{-1\normalbaselineskip}[0pt][0pt]{4} & \raisebox{-1\normalbaselineskip}[0pt][0pt]{γ} \\
3 & Borealis e duabus his quae sunt in capite. & 337 & 10 & 9 & 20 & b & 4 & \textit{b} \\
4 & Praecedens e duabus quae sunt in dorso. & 339 & 20 & 9 & 30 & b & 4 & θ \\
5 & Secunda e duabus quae sunt in dorso. & 341 & 50 & 7 & 30 & b & 4 & ι \\
6 & Praecedens e duabus quae sunt in ventre\newline\hspace*{4mm}Piscis australis. & \raisebox{-1\normalbaselineskip}[0pt][0pt]{337} & \raisebox{-1\normalbaselineskip}[0pt][0pt]{10} & \raisebox{-1\normalbaselineskip}[0pt][0pt]{4} & \raisebox{-1\normalbaselineskip}[0pt][0pt]{30} & \raisebox{-1\normalbaselineskip}[0pt][0pt]{b} & \raisebox{-1\normalbaselineskip}[0pt][0pt]{4} & \raisebox{-1\normalbaselineskip}[0pt][0pt]{ϰ} \\
7 & Posterior e duabus iis. & 340 & 50 & 3 & 30 & b & 4 & λ \\
8 & Quae est in cauda huius Piscis. & 347 & 10 & 6 & 20 & b & 4 & ω \\
\cline{2-8}\end{tabular}\end{center}
{\small
\par\hangindent=12mm\hangafter=1 22. Long. cod. \textarabic{\AbjadZero{}} 0′ pro \textarabic{لـ} 20′ (B., H. et S.); L. 50′. — Plaga cod. borealis. — Magnitudo cod. \textarabic{د} 4 pro \textarabic{ه} 5.
\par\hangindent=12mm\hangafter=1 23. Gradus long. cod. \textarabic{سكح} 328° pro \textarabic{سكج} 323°; minuta long. ut B., H. et S.; L. 5′. — Lat. ut B. et H.; L. et S. 14° 50′.
\par\hangindent=12mm\hangafter=1 24. Long. cod. \textarabic{سير} 317° pro \textarabic{سيب} 312°. — Lat. cod. \textarabic{يا} 11° pro \textarabic{ا} 1°.
\par\hangindent=12mm\hangafter=1 26. Plaga cod. borealis ut L.
\par\hangindent=12mm\hangafter=1 27. Long. cod. \textarabic{ل} 30′ pro \textarabic{ن} 50′. — Magnitudo cod. \textarabic{ب} 2 pro \textarabic{د} 4; L. 4; Pts. et S. 4 magna; Arg. 4 parva.
\par\hangindent=12mm\hangafter=1 28 et 29. Magnitudo ut Pts. et S.; L. 4; Arg. 5 magna.
\par\hspace*{8mm}\textbf{Pisces}:
\par\hangindent=12mm\hangafter=1 1. Magnitudo ut Pts. et L.; S. 4; Arg. 5 magna.
\par\hangindent=12mm\hangafter=1 2. Long. cod. \textarabic{سلد} 334° pro \textarabic{سله} 335°. — Magnitudo cod. \textarabic{ـجـ} 3 pro \textarabic{د} 4.
\par\hangindent=12mm\hangafter=1 5. Lat. cod. \textarabic{لـ} 20° pro \textarabic{ر} 7°.
\par\hangindent=12mm\hangafter=1 6. Long. cod. \textarabic{سم} 340° pro \textarabic{سلر} 337°; fortasse e confusione cum sequenti stella. — Lat. cod. \textarabic{ـجـ} 3° pro \textarabic{د} 4°.
\par\hangindent=12mm\hangafter=1 7. Long. cod. \textarabic{سمٮ} 342° vel \textarabic{سمر} 347° pro \textarabic{سم} 340°.
\par}
\clearpage
\noindent\makebox[\textwidth]{\textbf{}\hfill{\small AL-BATTANI OPUS ASTRONOMICUM}\hfill\textbf{167}}\par\vspace{-1mm}\noindent\rule{\textwidth}{.4pt}\par\vspace{2mm}
\begin{center}\begin{tabular}{r||p{78mm}|r@{\hspace{2.2mm}}r|r@{\hspace{2.2mm}}r|c|c||l}\cline{2-8}
 & \centering Stellarum descriptio. & \multicolumn{2}{c|}{Longitudo.} & \multicolumn{2}{c|}{Latitudo.} & \parbox{9mm}{\centering\small Plaga Caeli.} & Magnitudo. & \\ \cline{2-8}
9 & Praecedens ex iis quae sunt in cauda in\newline\hspace*{4mm}filo lini. & \raisebox{-1\normalbaselineskip}[0pt][0pt]{352°} & \raisebox{-1\normalbaselineskip}[0pt][0pt]{10′} & \raisebox{-1\normalbaselineskip}[0pt][0pt]{5°} & \raisebox{-1\normalbaselineskip}[0pt][0pt]{45′} & \raisebox{-1\normalbaselineskip}[0pt][0pt]{b} & \raisebox{-1\normalbaselineskip}[0pt][0pt]{6} & \raisebox{-1\normalbaselineskip}[0pt][0pt]{\textit{d}} \\
10 & Quae hanc sequitur. & 354 & 10 & 3 & 45 & b & 6 & Fl. 51. \\
11 & Praecedens e tribus lucidis quae sunt post\newline\hspace*{4mm}eas quarum praecessit mentio. & \raisebox{-1\normalbaselineskip}[0pt][0pt]{358} & \raisebox{-1\normalbaselineskip}[0pt][0pt]{20} & \raisebox{-1\normalbaselineskip}[0pt][0pt]{2} & \raisebox{-1\normalbaselineskip}[0pt][0pt]{15} & \raisebox{-1\normalbaselineskip}[0pt][0pt]{b} & \raisebox{-1\normalbaselineskip}[0pt][0pt]{4} & \raisebox{-1\normalbaselineskip}[0pt][0pt]{δ} \\
12 & Media ex iis. & 1 & 40 & 1 & 10 & b & 4 & ε \\
13 & Quae hanc sequitur, e tribus. & 4 & 10 & 6 & 0 & a & 4 & ζ \\
14 & Borealis e duabus proximis quae sunt in\newline\hspace*{4mm}\textit{ansa}. & \raisebox{-1\normalbaselineskip}[0pt][0pt]{3} & \raisebox{-1\normalbaselineskip}[0pt][0pt]{30} & \raisebox{-1\normalbaselineskip}[0pt][0pt]{2} & \raisebox{-1\normalbaselineskip}[0pt][0pt]{0} & \raisebox{-1\normalbaselineskip}[0pt][0pt]{a} & \raisebox{-1\normalbaselineskip}[0pt][0pt]{6} & \raisebox{-1\normalbaselineskip}[0pt][0pt]{\textit{e}} \\
15 & Australis e duabus iis. & 4 & 30 & 5 & 0 & a & 6 & \textit{f} \\
16 & Media e tribus quae sunt post \textit{ansam}. & 9 & 30 & 4 & 40 & a & 4 & ν \\
17 & Quae est in nodo duorum filorum lini. & 13 & 40 & 8 & 30 & a & 3 & α \\
18 & Praecedens a septentrionibus nodi. & 11 & 40 & 1 & 20 & b & 4 & ο \\
19 & Media e tribus quae sunt in nodo. & 11 & 50 & 5 & 20 & b & 3 & η \\
20 & Borealis e duabus quae sunt in ore Piscis\newline\hspace*{4mm}borealis. & \raisebox{-1\normalbaselineskip}[0pt][0pt]{13} & \raisebox{-1\normalbaselineskip}[0pt][0pt]{10} & \raisebox{-1\normalbaselineskip}[0pt][0pt]{21} & \raisebox{-1\normalbaselineskip}[0pt][0pt]{45} & \raisebox{-1\normalbaselineskip}[0pt][0pt]{b} & \raisebox{-1\normalbaselineskip}[0pt][0pt]{5} & \raisebox{-1\normalbaselineskip}[0pt][0pt]{\textit{g}} \\
21 & Borealis e tribus quae sunt in extremi-\newline\hspace*{4mm}tate caudae. & \raisebox{-1\normalbaselineskip}[0pt][0pt]{11} & \raisebox{-1\normalbaselineskip}[0pt][0pt]{40} & \raisebox{-1\normalbaselineskip}[0pt][0pt]{9} & \raisebox{-1\normalbaselineskip}[0pt][0pt]{0} & \raisebox{-1\normalbaselineskip}[0pt][0pt]{b} & \raisebox{-1\normalbaselineskip}[0pt][0pt]{4} & \raisebox{-1\normalbaselineskip}[0pt][0pt]{ρ} \\
22 & Praecedens e tribus quae sunt in spina dorsi\newline\hspace*{4mm}huius Piscis. & \raisebox{-1\normalbaselineskip}[0pt][0pt]{6} & \raisebox{-1\normalbaselineskip}[0pt][0pt]{50} & \raisebox{-1\normalbaselineskip}[0pt][0pt]{14} & \raisebox{-1\normalbaselineskip}[0pt][0pt]{20} & \raisebox{-1\normalbaselineskip}[0pt][0pt]{b} & \raisebox{-1\normalbaselineskip}[0pt][0pt]{4} & \raisebox{-1\normalbaselineskip}[0pt][0pt]{ψ¹} \\
23 & Media ex iis. & 7 & 50 & 13 & 0 & b & 4 & ψ² \\
24 & Postrema ex his tribus. & 8 & 50 & 12 & 0 & b & 4 & ψ³ \\
25 & Borealis e duabus quae sunt in ventre. & 13 & 20 & 17 & 0 & b & 4 & υ \\
26 & Australis e duabus iis. & 11 & 0 & 15 & 20 & b & 4 & φ \\
27 & Quae est in spina posteriore prope cau-\newline\hspace*{4mm}dam. & \raisebox{-1\normalbaselineskip}[0pt][0pt]{11} & \raisebox{-1\normalbaselineskip}[0pt][0pt]{10} & \raisebox{-1\normalbaselineskip}[0pt][0pt]{11} & \raisebox{-1\normalbaselineskip}[0pt][0pt]{45} & \raisebox{-1\normalbaselineskip}[0pt][0pt]{b} & \raisebox{-1\normalbaselineskip}[0pt][0pt]{4} & \raisebox{-1\normalbaselineskip}[0pt][0pt]{χ} \\
28 & Et ex iis quae sunt fere sub quadrila-\newline\hspace*{4mm}tero. & \raisebox{-1\normalbaselineskip}[0pt][0pt]{342} & \raisebox{-1\normalbaselineskip}[0pt][0pt]{20} & \raisebox{-1\normalbaselineskip}[0pt][0pt]{2} & \raisebox{-1\normalbaselineskip}[0pt][0pt]{40} & \raisebox{-1\normalbaselineskip}[0pt][0pt]{a} & \raisebox{-1\normalbaselineskip}[0pt][0pt]{4} & \raisebox{-1\normalbaselineskip}[0pt][0pt]{Fl. 27.} \\
\cline{2-8}\end{tabular}\end{center}
{\small
\par\hangindent=12mm\hangafter=1 9. Pt.: « Prima ex iis quae sunt in filo lini, [incipiendo] a cauda » (τῶν κατὰ τὸν λίνον αὐτοῦ ὁ πρῶτος ἀπὸ τῆς οὐρᾶς). Sed etiam L.: « prima stellarum quae sunt in cauda »; pro ἀπὸ legerunt ἐπί.
\par\hangindent=12mm\hangafter=1 10. Long. cod. \textarabic{سنه} 355° pro \textarabic{سند} 354°. — Magnitudo cod. \textarabic{ه} 5 pro \textarabic{و} 6.
\par\hangindent=12mm\hangafter=1 12. Long. ut L.; B., H. et S. 1° 20′.
\par\hangindent=12mm\hangafter=1 13. Lat. ut H. et S.; L. 1° 20′; sed rectius B. 0° 10′ (ϛ′ $\tfrac{1^\circ}{6}$ pro ϛ 6°). — Plaga cod. borealis.
\par\hangindent=12mm\hangafter=1 14. « Ansa » error, ut in nr. 16, pro « voluta (seu curvatura) lini » καμπή. — Plaga cod. borealis.
\par\hangindent=12mm\hangafter=1 15. Plaga cod. borealis.
\par\hangindent=12mm\hangafter=1 16. Plaga cod. borealis.
\par\hangindent=12mm\hangafter=1 18. Lat. ut S.; L. 5° 20′; B. et H. 1° 40′. — Plaga ut L. et S; B. et H. australis.
\par\hangindent=12mm\hangafter=1 23. Lat. ut L. et S.; B. et H. 13° 15′ — Magnitudo cod. \textarabic{ـجـ} 3 pro \textarabic{د} 4.
\par\hangindent=12mm\hangafter=1 25. Male long. H. 12° 50′. — Magnitudo cod. \textarabic{و} 6 pro \textarabic{د} 4.
\par\hangindent=12mm\hangafter=1 28. Long. cod. \textarabic{سنب مو} 352° 46′ pro \textarabic{سمب لـ} 342° 20′.
\par}
\clearpage
\noindent\makebox[\textwidth]{\textbf{168}\hfill{\small C. A. NALLINO}\hfill\textbf{}}\par\vspace{-1mm}\noindent\rule{\textwidth}{.4pt}\par\vspace{2mm}
\begin{center}III.\end{center}
\begin{center}Constellationes ab austro zodiaci.\end{center}
\begin{center}\begin{tabular}{r||p{78mm}|r@{\hspace{2.2mm}}r|r@{\hspace{2.2mm}}r|c|c||l}\cline{2-8}
 & \centering Stellarum descriptio. & \multicolumn{2}{c|}{Longitudo.} & \multicolumn{2}{c|}{Latitudo.} & \parbox{9mm}{\centering\small Plaga Caeli.} & Magnitudo. & \\ \cline{2-8}
 & \multicolumn{7}{c||}{\parbox{150mm}{\centering\rule{0pt}{5mm}\textbf{Cetus} (\textit{qīṭos} « ϰῆτος » sive \textit{sabuʿ al–baḥr} « bellua marina »).}} & \\ \cline{2-8}
1 & Quae est in extremitate nasi. & 28° & 50′ & 7° & 45′ & a & 4 & λ \\
2 & Postrema e tribus quae sunt in gutture\newline\hspace*{4mm}in extremitate maxillae. & \raisebox{-1\normalbaselineskip}[0pt][0pt]{28} & \raisebox{-1\normalbaselineskip}[0pt][0pt]{50} & \raisebox{-1\normalbaselineskip}[0pt][0pt]{12} & \raisebox{-1\normalbaselineskip}[0pt][0pt]{20} & \raisebox{-1\normalbaselineskip}[0pt][0pt]{a} & \raisebox{-1\normalbaselineskip}[0pt][0pt]{3} & \raisebox{-1\normalbaselineskip}[0pt][0pt]{α} \\
3 & Media ex iis, in medio oris. & 23 & 50 & 11 & 30 & a & 3 & γ \\
4 & Praecedens e tribus quae sunt in barba. & 21 & 40 & 14 & 0 & a & 3 & δ \\
5 & Quae est in fronte supra oculum. & \textit{21} & \textit{20} & 8 & 10 & a & 4 & ν \\
6 & Australis lateris posterioris. & 18 & 10 & 27 & 30 & a & 3 & π \\
7 & Media e tribus quae sunt in corpore. & 3 & 10 & 25 & 20 & a & 3 & τ \\
8 & Borealis ex his tribus. & 6 & 10 & 20 & 0 & a & 3 & ζ \\
9 & Posterior e duabus quae sunt in cauda. & 0 & 50 & 15 & 20 & a & 3 & θ \\
10 & Anterior e duabus iis. & 356 & 10 & 15 & 40 & a & 3 & η \\
11 & Borealis quae est in latere anteriore. & 350 & 30 & 13 & 0 & a & 5 magna &  \\
12 & Borealis prope duas quae sunt in extremi-\newline\hspace*{4mm}tate caudae. & \raisebox{-1\normalbaselineskip}[0pt][0pt]{345} & \raisebox{-1\normalbaselineskip}[0pt][0pt]{50} & \raisebox{-1\normalbaselineskip}[0pt][0pt]{9} & \raisebox{-1\normalbaselineskip}[0pt][0pt]{40} & \raisebox{-1\normalbaselineskip}[0pt][0pt]{a} & \raisebox{-1\normalbaselineskip}[0pt][0pt]{3 parva} & \raisebox{-1\normalbaselineskip}[0pt][0pt]{ι} \\
13 & Australis quae est in extremitate caudae. & 346 & 50 & 20 & 20 & a & 3 magna & β \\
 & \multicolumn{7}{c||}{\parbox{150mm}{\centering\rule{0pt}{5mm}\textbf{Orion} (\textit{al–ǵabbār} « gigas »).}} & \\ \cline{2-8}
1 & Nebulosa quae est in capite. & 68 & 10 & 16 & 30 & a & nebulosa & λ \\
2 & Lucida quae est supra humerum dexterum,\newline\hspace*{4mm}et appellatur \textit{mankib al–ǵawzā’} « hu-\newline\hspace*{4mm}« merus \textit{al–ǵawzā’} ». & \raisebox{-2\normalbaselineskip}[0pt][0pt]{73} & \raisebox{-2\normalbaselineskip}[0pt][0pt]{10} & \raisebox{-2\normalbaselineskip}[0pt][0pt]{17} & \raisebox{-2\normalbaselineskip}[0pt][0pt]{0} & \raisebox{-2\normalbaselineskip}[0pt][0pt]{a} & \raisebox{-2\normalbaselineskip}[0pt][0pt]{1 parva} & \raisebox{-2\normalbaselineskip}[0pt][0pt]{α} \\
\cline{2-8}\end{tabular}\end{center}
{\small
\par\hspace*{8mm}\textbf{Cetus}:
\par\hangindent=12mm\hangafter=1 1. Magnitudo cod. \textarabic{ـجـ} 3 pro \textarabic{د} 4.
\par\hangindent=12mm\hangafter=1 4. Male lat. H. 14° 15′.
\par\hangindent=12mm\hangafter=1 5. Long. B. 21° 30′; H. 22° 20′ (quod, \Name{Schiaparelli} teste, verius est); L. 21° 50′; S. 20° 30′. — Lat. cod. \textarabic{مه} 45′ pro \textarabic{ى} 10′.
\par\hangindent=12mm\hangafter=1 9. Lat. ut L. et S.; B. et H. 15° 40′.
\par\hangindent=12mm\hangafter=1 10. Long. cod. \textarabic{سمو} 346° pro \textarabic{سنو} 356°. — Lat. cod. \textarabic{ل} 30′ pro \textarabic{م} 40′.
\par\hangindent=12mm\hangafter=1 11. Magnitudo ut Pts.; L. 5; S. 5 parva vel 6 magna; Arg. 6. — \Name{Baily}: « I cannot identify this star of « the quadrilateral figure, alluded to by Ptolemy: in fact, the figure itself is not easily recognised, « so as to correspond with the description ».
\par\hangindent=12mm\hangafter=1 12. Magnitudo cod. 3 magna, perperam enim in sequenti stella « parvam » pro « magna » scripsit. Pts., S. et Arg 3 parva; L. 3.
\par\hangindent=12mm\hangafter=1 13. Long. cod. \textarabic{سمه} 345° pro \textarabic{سمو} 346°. — Lat. cod. \textarabic{م} 40′ pro \textarabic{لـ} 20′. — Magnitudo cod. 4 parva (cfr. nr. 12); Pts., L. et S. 3 magna; Arg. 2.
\par}
\clearpage
\noindent\makebox[\textwidth]{\textbf{}\hfill{\small AL-BATTANI OPUS ASTRONOMICUM}\hfill\textbf{169}}\par\vspace{-1mm}\noindent\rule{\textwidth}{.4pt}\par\vspace{2mm}
\begin{center}\begin{tabular}{r||p{78mm}|r@{\hspace{2.2mm}}r|r@{\hspace{2.2mm}}r|c|c||l}\cline{2-8}
 & \centering Stellarum descriptio. & \multicolumn{2}{c|}{Longitudo.} & \multicolumn{2}{c|}{Latitudo.} & \parbox{9mm}{\centering\small Plaga Caeli.} & Magnitudo. & \\ \cline{2-8}
3 & Quae est in humero sinistro. & 65° & 10′ & 17° & 30′ & a & 2 parva & γ \\
4 & Stella sub hac quae est in humero sinistro. & 66 & 10 & 18 & 0 & a & 4 parva & A \\
5 & Quae est in cubito dextero. & 75 & 30 & 14 & 30 & a & 4 & μ \\
6 & Borealis ex novem quae sunt in pelle manu\newline\hspace*{4mm}sinistra gestata. & \raisebox{-1\normalbaselineskip}[0pt][0pt]{61} & \raisebox{-1\normalbaselineskip}[0pt][0pt]{40} & \raisebox{-1\normalbaselineskip}[0pt][0pt]{8} & \raisebox{-1\normalbaselineskip}[0pt][0pt]{0} & \raisebox{-1\normalbaselineskip}[0pt][0pt]{a} & \raisebox{-1\normalbaselineskip}[0pt][0pt]{4} & \raisebox{-1\normalbaselineskip}[0pt][0pt]{Fl. 15.} \\
7 & Sexta earum, in parte boreali. & 56 & 0 & 15 & 50 & a & 3 & π¹ \\
8 & Septima, quae est post hanc in parte boreali. & 56 & 0 & 17 & 10 & a & 3 & π³ \\
9 & Octava, quae est post eam in parte boreali. & 56 & 30 & 20 & 20 & a & 3 & π⁵ \\
10 & Nona ex iis quae sunt in pelle, ipsaque est\newline\hspace*{4mm}in parte australi et vocatur \textit{al–miǵlad}\newline\hspace*{4mm}« flagellum ». & \raisebox{-2\normalbaselineskip}[0pt][0pt]{57} & \raisebox{-2\normalbaselineskip}[0pt][0pt]{30} & \raisebox{-2\normalbaselineskip}[0pt][0pt]{21} & \raisebox{-2\normalbaselineskip}[0pt][0pt]{30} & \raisebox{-2\normalbaselineskip}[0pt][0pt]{a} & \raisebox{-2\normalbaselineskip}[0pt][0pt]{3} & \raisebox{-2\normalbaselineskip}[0pt][0pt]{π⁶} \\
11 & Praecedens e tribus quae sunt in zona. & 66 & 30 & 24 & 10 & a & 2 & δ \\
12 & Media ex iis. & 68 & 30 & 24 & 50 & a & 2 & ε \\
13 & Postrema e tribus quae sunt in zona. & 69 & 20 & 25 & 40 & a & 2 & ζ \\
14 & Quae est prope capulum ensis. & 65 & 0 & 25 & 50 & a & 3 & η \\
15 & Borealis ex tribus collectis quae sunt prope\newline\hspace*{4mm}cuspidem ensis. & \raisebox{-1\normalbaselineskip}[0pt][0pt]{67} & \raisebox{-1\normalbaselineskip}[0pt][0pt]{40} & \raisebox{-1\normalbaselineskip}[0pt][0pt]{28} & \raisebox{-1\normalbaselineskip}[0pt][0pt]{40} & \raisebox{-1\normalbaselineskip}[0pt][0pt]{a} & \raisebox{-1\normalbaselineskip}[0pt][0pt]{4} & \raisebox{-1\normalbaselineskip}[0pt][0pt]{\textit{c}} \\
16 & Media ex his tribus. & 67 & 50 & 29 & 10 & a & 3 parva & θ² \\
17 & Praecedens et australis ex his tribus. & 67 & 50 & 29 & 50 & a & 3 & ι \\
18 & Lucida quae est in extremitate pedis sinistri\newline\hspace*{4mm}et appellatur \textit{riǵl al–ǵawzā’} « pes \textit{al–}\newline\hspace*{4mm}« \textit{–ǵawzā’} ». & \raisebox{-2\normalbaselineskip}[0pt][0pt]{61} & \raisebox{-2\normalbaselineskip}[0pt][0pt]{0} & \raisebox{-2\normalbaselineskip}[0pt][0pt]{31} & \raisebox{-2\normalbaselineskip}[0pt][0pt]{30} & \raisebox{-2\normalbaselineskip}[0pt][0pt]{a} & \raisebox{-2\normalbaselineskip}[0pt][0pt]{1} & \raisebox{-2\normalbaselineskip}[0pt][0pt]{β} \\
19 & Stella a septentrionibus huius, supra mal-\newline\hspace*{4mm}leolum. & \raisebox{-1\normalbaselineskip}[0pt][0pt]{62} & \raisebox{-1\normalbaselineskip}[0pt][0pt]{10} & \raisebox{-1\normalbaselineskip}[0pt][0pt]{30} & \raisebox{-1\normalbaselineskip}[0pt][0pt]{15} & \raisebox{-1\normalbaselineskip}[0pt][0pt]{a} & \raisebox{-1\normalbaselineskip}[0pt][0pt]{4} & \raisebox{-1\normalbaselineskip}[0pt][0pt]{τ} \\
20 & Externa sub talo sinistro. & 64 & 30 & 31 & 10 & a & 4 & \textit{e} \\
21 & Quae est sub genu dextero posteriore. & 71 & 20 & 33 & 30 & a & 3 magna & ϰ \\
\cline{2-8}\end{tabular}\end{center}
{\small
\par\hspace*{8mm}\textbf{Orion}:
\par\hangindent=12mm\hangafter=1 4. Long. cod. \textarabic{لـ} 20′ pro \textarabic{ى} 10′. — Lat. cod \textarabic{ح} 8° pro \textarabic{يح} 18°.
\par\hangindent=12mm\hangafter=1 5. Male H. lat. 14° 20′.
\par\hangindent=12mm\hangafter=1 9. Lat. cod. \textarabic{كا} 21° pro \textarabic{ك} 20°.
\par\hangindent=12mm\hangafter=1 10. Nomen \textit{al-miǵlad} mihi aliunde ignotum et ex coniectura restitutum cod. \textarabic{المحله}. — Lat. cod. \textarabic{كط} 29° pro \textarabic{كا} 21°. — Magnitudo ut B., L et Pts.; H. 2; S. 4; Arg. 5 magna.
\par\hangindent=12mm\hangafter=1 12. Lat. cod. \textarabic{ل} 30′ pro \textarabic{ن} 50′.
\par\hangindent=12mm\hangafter=1 13. Magnitudo ut B., Pts., S. et Arg.; H. et L. 3.
\par\hangindent=12mm\hangafter=1 15. Long. cod. \textarabic{ن} 50′ pro \textarabic{م} 40′. — Gradus lat. cod. \textarabic{لـ} 20° pro \textarabic{كح} 28°; minuta lat. ut L. et S; B. et H. 20′.
\par\hangindent=12mm\hangafter=1 16. Long. ut B., L. et S.; male H. 64° 30′. — Minuta lat. ut B. et S.; H. 0′; L. 40′.
\par\hangindent=12mm\hangafter=1 17. Long. ut S.; Pt. 68° 10′.
\par\hangindent=12mm\hangafter=1 18. Long. ut B. et S.; male H. 62° 0′, L. 60° 20′.
\par\hangindent=12mm\hangafter=1 19. Magnitudo L., Pts. et S. 4 magna; Arg. 4.
\par\hangindent=12mm\hangafter=1 21. Magnitudo ut Arg.; L. 3; Pts. et S. 3 parva.
\par}
\par\vfill\hfill 22
\clearpage
\noindent\makebox[\textwidth]{\textbf{170}\hfill{\small C. A. NALLINO}\hfill\textbf{}}\par\vspace{-1mm}\noindent\rule{\textwidth}{.4pt}\par\vspace{2mm}
\begin{center}\begin{tabular}{r||p{78mm}|r@{\hspace{2.2mm}}r|r@{\hspace{2.2mm}}r|c|c||l}\cline{2-8}
 & \centering Stellarum descriptio. & \multicolumn{2}{c|}{Longitudo.} & \multicolumn{2}{c|}{Latitudo.} & \parbox{9mm}{\centering\small Plaga Caeli.} & Magnitudo. & \\ \cline{2-8}
 & \multicolumn{7}{c||}{\parbox{150mm}{\centering\rule{0pt}{5mm}\textbf{Eridanus} (\textit{an–nahr} « fluvius »).}} & \\ \cline{2-8}
1 & Quae est in extremitate pedis Orionis et in\newline\hspace*{4mm}initio fluvii. & \raisebox{-1\normalbaselineskip}[0pt][0pt]{59°} & \raisebox{-1\normalbaselineskip}[0pt][0pt]{30′} & \raisebox{-1\normalbaselineskip}[0pt][0pt]{31°} & \raisebox{-1\normalbaselineskip}[0pt][0pt]{50′} & \raisebox{-1\normalbaselineskip}[0pt][0pt]{a} & \raisebox{-1\normalbaselineskip}[0pt][0pt]{4 magna} & \raisebox{-1\normalbaselineskip}[0pt][0pt]{λ} \\
2 & Lucida magna, quae est ultima stellarum\newline\hspace*{4mm}fluvii. & \raisebox{-1\normalbaselineskip}[0pt][0pt]{11} & \raisebox{-1\normalbaselineskip}[0pt][0pt]{20} & \raisebox{-1\normalbaselineskip}[0pt][0pt]{53} & \raisebox{-1\normalbaselineskip}[0pt][0pt]{30} & \raisebox{-1\normalbaselineskip}[0pt][0pt]{a} & \raisebox{-1\normalbaselineskip}[0pt][0pt]{1} & \raisebox{-1\normalbaselineskip}[0pt][0pt]{θ} \\
3 & Sequens ex quatuor quae sunt in intervallo\newline\hspace*{4mm}[ulteriore]. & \raisebox{-1\normalbaselineskip}[0pt][0pt]{38} & \raisebox{-1\normalbaselineskip}[0pt][0pt]{10} & \raisebox{-1\normalbaselineskip}[0pt][0pt]{32} & \raisebox{-1\normalbaselineskip}[0pt][0pt]{50} & \raisebox{-1\normalbaselineskip}[0pt][0pt]{a} & \raisebox{-1\normalbaselineskip}[0pt][0pt]{3} & \raisebox{-1\normalbaselineskip}[0pt][0pt]{γ} \\
4 & Tertia ante mediam. & 35 & 20 & 28 & 50 & a & 3 & δ \\
5 & Anterior ex quatuor. & 33 & 10 & 28 & 0 & a & 3 & ε \\
6 & Postrema ex quatuor. & 28 & 20 & 25 & 30 & a & 3 & ζ \\
7 & Tertia quae est ante hanc quartam. & 23 & 20 & 23 & 30 & a & 3 & η \\
8 & Quae est in curvatura fluvii et in fine pe-\newline\hspace*{4mm}ctoris Ceti. & \raisebox{-1\normalbaselineskip}[0pt][0pt]{16} & \raisebox{-1\normalbaselineskip}[0pt][0pt]{20} & \raisebox{-1\normalbaselineskip}[0pt][0pt]{32} & \raisebox{-1\normalbaselineskip}[0pt][0pt]{10} & \raisebox{-1\normalbaselineskip}[0pt][0pt]{a} & \raisebox{-1\normalbaselineskip}[0pt][0pt]{4} & \raisebox{-1\normalbaselineskip}[0pt][0pt]{τ¹} \\
9 & Borealis quae est in latere anteriore ex\newline\hspace*{4mm}quatuor quae sunt in trapezio (?). & \raisebox{-1\normalbaselineskip}[0pt][0pt]{32} & \raisebox{-1\normalbaselineskip}[0pt][0pt]{30} & \raisebox{-1\normalbaselineskip}[0pt][0pt]{41} & \raisebox{-1\normalbaselineskip}[0pt][0pt]{20} & \raisebox{-1\normalbaselineskip}[0pt][0pt]{a} & \raisebox{-1\normalbaselineskip}[0pt][0pt]{4} & \raisebox{-1\normalbaselineskip}[0pt][0pt]{τ⁶} \\
10 & Postrema ex his quatuor. & 35 & 50 & 43 & 20 & a & 4 & τ⁹ \\
11 & Posterior e duabus quae sunt post curva-\newline\hspace*{4mm}turam et appellatur \textarabic{التربه}. & \raisebox{-1\normalbaselineskip}[0pt][0pt]{39} & \raisebox{-1\normalbaselineskip}[0pt][0pt]{20} & \raisebox{-1\normalbaselineskip}[0pt][0pt]{53} & \raisebox{-1\normalbaselineskip}[0pt][0pt]{50} & \raisebox{-1\normalbaselineskip}[0pt][0pt]{a} & \raisebox{-1\normalbaselineskip}[0pt][0pt]{4} & \raisebox{-1\normalbaselineskip}[0pt][0pt]{υ⁵} \\
12 & Postrema e tribus quae sunt post has duas. & 29 & 0 & 53 & 0 & a & 4 & υ³ \\
13 & Borealis e duabus inter se oppositis. & 45 & 20 & 53 & 20 & a & 4 & υ⁶ \\
14 & Australis e duabus iis. & 46 & 10 & 51 & 45 & a & 4 & υ \\
\cline{2-8}\end{tabular}\end{center}
{\small
\par\hspace*{8mm}\textbf{Eridanus}:
\par\hangindent=12mm\hangafter=1 1. Long. cod. \textarabic{صط} 69° pro \textarabic{نط} 59°. — Lat. cod. \textarabic{ل} 30° pro \textarabic{لا} 31°. — Magnitudo L., Pts, S. et Arg. 4 (sine « magna »).
\par\hangindent=12mm\hangafter=1 2. Long. ut B., L. et S.; H. incertus utrum 18° 50′ an 358° 40′. — Lat. cod. \textarabic{لج} 33° pro \textarabic{نج} 53°. — \Name{Baily}, cui \Name{Schiaparelli} assentit, iure dicit: « Most of the commentators on Ptolemy’s catalogue have « supposed this star to be \textit{Achernar} (1): but neither the longitude nor latitude of any of the copies « will agree with the position of that star; and moreover \textit{Achernar} was not visible at Alexandria. « The magnitude has probably changed since Ptolemy’s time ». Est ϑ Eridani, 225 Lacaille.
\par\hangindent=12mm\hangafter=1 3. Cod. « sequens quatuor ..... », omissa particula \textarabic{من} « ex ». — Lat. ut B., L. et S; H. 33° 10′.
\par\hangindent=12mm\hangafter=1 6. Long. cod. \textarabic{ل} 30′ pro \textarabic{لـ} 20′.
\par\hangindent=12mm\hangafter=1 7. Lat. ut quidam codices Graeci; L. et S 23° 50′; B. et H. 24° 30′.
\par\hangindent=12mm\hangafter=1 9. Cod. \textarabic{في الجوزا} quod non intelligo; fortasse \textarabic{في الجواز} « in vado »? Pt. ἐν τραπεζίῳ.
\par\hangindent=12mm\hangafter=1 11. Nomen cod. \textit{at-turbah} significat « tumulum, sepulchrum »; sed lectio procul dubio falsa. L. « beemum ».
\par\hangindent=12mm\hangafter=1 12. Minuta lat. ut B., L. et S; H. 10′.
\par\hangindent=12mm\hangafter=1 13. Gradus lat. cod. \textarabic{صج} 63° pro \textarabic{نج} 53°, quos H. et S. habent; sed vera lectio est 50° 20′ ut apud B. et L., atque ut e descriptione Ptolemaica patet.
\par\hangindent=12mm\hangafter=1 14. Lat. cod. \textarabic{ن} 50′ pro \textarabic{مه} 45′.
\par}
\par\vspace{2mm}\begin{center}\rule{40mm}{.4pt}\end{center}{\small (1) I. e. α Eridani. Ita etiam \Name{Schjellerup}.}
\clearpage
\noindent\makebox[\textwidth]{\textbf{}\hfill{\small AL-BATTANI OPUS ASTRONOMICUM}\hfill\textbf{171}}\par\vspace{-1mm}\noindent\rule{\textwidth}{.4pt}\par\vspace{2mm}
\begin{center}\begin{tabular}{r||p{78mm}|r@{\hspace{2.2mm}}r|r@{\hspace{2.2mm}}r|c|c||l}\cline{2-8}
 & \centering Stellarum descriptio. & \multicolumn{2}{c|}{Longitudo.} & \multicolumn{2}{c|}{Latitudo.} & \parbox{9mm}{\centering\small Plaga Caeli.} & Magnitudo. & \\ \cline{2-8}
 & \multicolumn{7}{c||}{\parbox{150mm}{\centering\rule{0pt}{5mm}\textbf{Lepus} (\textit{al–arnab}).}} & \\ \cline{2-8}
1 & Borealis in latere anteriore ex quatuor quae\newline\hspace*{4mm}sunt in \textit{dorso}. & \raisebox{-1\normalbaselineskip}[0pt][0pt]{60°} & \raisebox{-1\normalbaselineskip}[0pt][0pt]{\textit{30}′} & \raisebox{-1\normalbaselineskip}[0pt][0pt]{35°} & \raisebox{-1\normalbaselineskip}[0pt][0pt]{0′} & \raisebox{-1\normalbaselineskip}[0pt][0pt]{a} & \raisebox{-1\normalbaselineskip}[0pt][0pt]{5} & \raisebox{-1\normalbaselineskip}[0pt][0pt]{ι} \\
2 & Borealis in latere posteriore ex iis. & 62 & 30 & 35 & 40 & a & 5 & ν \\
3 & [Australis ex iis in latere anteriore]. & 61 & 0 & 36 & 30 & a & 5 & ϰ \\
4 & Quae est in mento. & 60 & 20 & 39 & 15 & a & 4 magna & μ \\
5 & Quae est ante extremitatem pedis anterio-\newline\hspace*{4mm}ris sinistri. & \raisebox{-1\normalbaselineskip}[0pt][0pt]{57} & \raisebox{-1\normalbaselineskip}[0pt][0pt]{20} & \raisebox{-1\normalbaselineskip}[0pt][0pt]{45} & \raisebox{-1\normalbaselineskip}[0pt][0pt]{15} & \raisebox{-1\normalbaselineskip}[0pt][0pt]{a} & \raisebox{-1\normalbaselineskip}[0pt][0pt]{4 magna} & \raisebox{-1\normalbaselineskip}[0pt][0pt]{ε} \\
6 & Quae est in medio corporis. & 67 & 0 & 41 & 30 & a & 3 & α \\
7 & Quae est sub ventre. & 65 & 30 & 44 & 20 & a & 3 & β \\
8 & Borealis e duabus quae sunt in pedibus po-\newline\hspace*{4mm}sterioribus. & \raisebox{-1\normalbaselineskip}[0pt][0pt]{72} & \raisebox{-1\normalbaselineskip}[0pt][0pt]{10} & \raisebox{-1\normalbaselineskip}[0pt][0pt]{44} & \raisebox{-1\normalbaselineskip}[0pt][0pt]{0} & \raisebox{-1\normalbaselineskip}[0pt][0pt]{a} & \raisebox{-1\normalbaselineskip}[0pt][0pt]{4 magna} & \raisebox{-1\normalbaselineskip}[0pt][0pt]{δ} \\
9 & Australis e duabus iis. & 70 & 10 & 45 & 50 & a & 4 magna & γ \\
10 & Quae est in dorso. & 71 & 10 & 38 & 20 & a & 4 magna & ζ \\
11 & Quae est in extremitate caudae. & 73 & 50 & 38 & 10 & a & 4 magna & η \\
 & \multicolumn{7}{c||}{\parbox{150mm}{\centering\rule{0pt}{5mm}\textbf{Canis maior} (\textit{al–kalb}).}} & \\ \cline{2-8}
1 & Stella quae est in auribus. & 90 & 50 & 35 & 0 & a & 4 & θ \\
2 & Quae est in capite. & 92 & 30 & 36 & 30 & a & 5 & μ \\
3 & Lucida quae est in ore Canis et vocatur\newline\hspace*{4mm}\textit{ash–shiʿrā al–yamāniyyah}. & \raisebox{-1\normalbaselineskip}[0pt][0pt]{88} & \raisebox{-1\normalbaselineskip}[0pt][0pt]{50} & \raisebox{-1\normalbaselineskip}[0pt][0pt]{39} & \raisebox{-1\normalbaselineskip}[0pt][0pt]{10} & \raisebox{-1\normalbaselineskip}[0pt][0pt]{a} & \raisebox{-1\normalbaselineskip}[0pt][0pt]{1} & \raisebox{-1\normalbaselineskip}[0pt][0pt]{α} \\
4 & Quae est in extremitate pedis anterioris. & 82 & 10 & 41 & 20 & a & 3 & β \\
\cline{2-8}\end{tabular}\end{center}
{\small
\par\hspace*{8mm}\textbf{Lepus}:
\par\hangindent=12mm\hangafter=1 1. « In dorso » falsum est; Pt. ceterique « in auribus ». — Minuta long. cod. fortasse falsa; L. et S. 50′; B. et H. 10′.
\par\hangindent=12mm\hangafter=1 3. In columna nominum siderum scriba hanc stellam praetergressus est, sed eius numeros in columnis reliquis servavit. Unde factum est ut numeri reliquarum stellarum Leporis et priorum quinque Canis maioris non stellis suis sed stellis sequentibus convenirent. Ut veram congruentiam demum restitueret, insulsus scriba in columna nominum, inter quintam et sextam stellam Canis maioris falso inseruit: « Quae est in initio femoris dexteri ». Lacuna ita impleta, numeri et stellae sequentes inter se conveniunt. Textum Arabicum ut in cod. est edidi; in versione formam genuinam restitui.
\par\hangindent=12mm\hangafter=1 6. Long. cod. \textarabic{صا} 61° pro \textarabic{صز} 67°.
\par\hangindent=12mm\hangafter=1 7. Long. ut S.; B., H. et L. 66° 0′. — Lat. cod. \textarabic{ن} 50′ pro \textarabic{لـ} 20′; male H. 41° 20′.
\par\hangindent=12mm\hangafter=1 9. Lat. cod. \textarabic{مه} 45′ pro \textarabic{ن} 50′.
\par\hangindent=12mm\hangafter=1 10. Lat. cod. \textarabic{م} 40′ pro \textarabic{لـ} 20′.
\par\hspace*{8mm}\textbf{Canis maior}:
\par\hangindent=12mm\hangafter=1 1. Long. cod. \textarabic{ل} 30′ pro \textarabic{ن} 50′. — De quinque prioribus stellis cfr. quae dixi ad nr. 3 Leporis.
\par\hangindent=12mm\hangafter=1 3. Long. cod. \textarabic{قح} 108° pro \textarabic{فح} 88°. — Lat. cod. \textarabic{نط} 59° pro \textarabic{لط} 39°.
\par\hangindent=12mm\hangafter=1 4. Long. cod. \textarabic{قب} 102° pro \textarabic{فب} 82°.
\par}
\clearpage
\noindent\makebox[\textwidth]{\textbf{172}\hfill{\small C. A. NALLINO}\hfill\textbf{}}\par\vspace{-1mm}\noindent\rule{\textwidth}{.4pt}\par\vspace{2mm}
\begin{center}\begin{tabular}{r||p{78mm}|r@{\hspace{2.2mm}}r|r@{\hspace{2.2mm}}r|c|c||l}\cline{2-8}
 & \centering Stellarum descriptio. & \multicolumn{2}{c|}{Longitudo.} & \multicolumn{2}{c|}{Latitudo.} & \parbox{9mm}{\centering\small Plaga Caeli.} & Magnitudo. & \\ \cline{2-8}
5 & Quae est in initio femoris sinistri. & 97° & 50′ & 48° & 45′ & a & 3 parva & δ \\
6 & Quae est sub ventre inter duo femora. & 94 & 50 & 51 & 30 & a & 3 & ε \\
7 & Quae est in pede dextero. & 80 & 50 & 53 & 45 & a & 3 & ζ \\
8 & Quae est in cauda. & 103 & 20 & 50 & 40 & a & 3 parva & η \\
 & \multicolumn{7}{l}{\hspace*{4mm}Et ex iis quae non sunt ei in figura:} & \\
9 & Posterior e duabus lucidis. & 70 & 10 & 59 & 40 & a & 2 & β Columbae. \\
10 & Praecedens e duabus iis. & 67 & 10 & 57 & 40 & a & 2 & α Columbae. \\
 & \multicolumn{7}{c||}{\parbox{150mm}{\centering\rule{0pt}{5mm}\textbf{Canis minor} (\textit{muqaddam al–kalb} « ante Canem positus, προκύων »).}} & \\ \cline{2-8}
1 & Quae est in collo. & 96 & 10 & 14 & 0 & a & 4 & β \\
2 & Lucida quae est post eam, et appellatur\newline\hspace*{4mm}\textit{ash–shiʿrā ash–sha’miyyah}. & \raisebox{-1\normalbaselineskip}[0pt][0pt]{100} & \raisebox{-1\normalbaselineskip}[0pt][0pt]{20} & \raisebox{-1\normalbaselineskip}[0pt][0pt]{16} & \raisebox{-1\normalbaselineskip}[0pt][0pt]{10} & \raisebox{-1\normalbaselineskip}[0pt][0pt]{a} & \raisebox{-1\normalbaselineskip}[0pt][0pt]{1} & \raisebox{-1\normalbaselineskip}[0pt][0pt]{α} \\
 & \multicolumn{7}{c||}{\parbox{150mm}{\centering\rule{0pt}{5mm}[\textit{Desideratur in codice folium, quo continebantur stellae Navis atque Hydrae,\\et primae stellae Crateris}].}} & \\ \cline{2-8}
 & \multicolumn{7}{c||}{\parbox{150mm}{\centering\rule{0pt}{5mm}\textbf{Reliquae stellae Crateris} (\textit{aqrāṭīros} « ϰρατήρ », sive \textit{al–ka’s} « poculum »).}} & \\ \cline{2-8}
1 & Quae est in labro boreali oris. & 160 & 30 & 13 & 40 & a & 4 & ε \\
2 & Quae est in ansa australi. & 170 & 20 & 16 & 10 & a & 4 & η \\
3 & Quae est in ansa boreali. & 162 & 50 & 11 & 50 & a & 4 & θ \\
 & \multicolumn{7}{c||}{\parbox{150mm}{\centering\rule{0pt}{5mm}\textbf{Corvus} (\textit{al–ghurāb}).}} & \\ \cline{2-8}
1 & Quae est in rostro Corvi, prope Hydram\newline\hspace*{4mm}(\textit{ash–shuǵāʿ}). & \raisebox{-1\normalbaselineskip}[0pt][0pt]{176} & \raisebox{-1\normalbaselineskip}[0pt][0pt]{30} & \raisebox{-1\normalbaselineskip}[0pt][0pt]{21} & \raisebox{-1\normalbaselineskip}[0pt][0pt]{40} & \raisebox{-1\normalbaselineskip}[0pt][0pt]{a} & \raisebox{-1\normalbaselineskip}[0pt][0pt]{3} & \raisebox{-1\normalbaselineskip}[0pt][0pt]{α} \\
\cline{2-8}\end{tabular}\end{center}
{\small
\par\hangindent=12mm\hangafter=1 5. Male H lat. 48° 50′. — Magnitudo ut L.; Pts. et S. 3; Arg. 2.
\par\hangindent=12mm\hangafter=1 7. Long. cod. \textarabic{ق} 100° pro \textarabic{ف} 80°.
\par\hangindent=12mm\hangafter=1 8. Lat. cod. \textarabic{ل} 30° pro \textarabic{ن} 50°. — Magnitudo cod. « nebulosa parva »; L. et Pts. 3; S. 3 parva; Arg. 3 magna.
\par\hspace*{8mm}\textbf{Canis minor}:
\par\hangindent=12mm\hangafter=1 1. Lat. cod. \textarabic{ند} 54° pro \textarabic{يد} 14°.
\par\hangindent=12mm\hangafter=1 2. Minuta long. ut L. et S.; B. et H. 40′. — Lat. cod. \textarabic{نو} 56° pro \textarabic{يو} 16°.
\par\hspace*{8mm}\textbf{Crater}:
\par\hangindent=12mm\hangafter=1 2. Long. cod. \textarabic{قعا} 171° pro \textarabic{قع} 170°. — Magnitudo ut L.; Pts. et S. 4 parva; Arg. 5.
\par\hangindent=12mm\hangafter=1 3. Long. ut S.; B. 162° 30′; H. male 174° 20′; L. ...° 50′. — Lat. cod. \textarabic{ل} 30′ pro \textarabic{ن} 50′. — Magnitudo cod. \textarabic{ه} 5 pro \textarabic{د} 4.
\par\hspace*{8mm}\textbf{Corvus}:
\par\hangindent=12mm\hangafter=1 1. Long. cod. \textarabic{قفو} 186° pro \textarabic{قعو} 176°.
\par}
\clearpage
\noindent\makebox[\textwidth]{\textbf{}\hfill{\small AL-BATTANI OPUS ASTRONOMICUM}\hfill\textbf{173}}\par\vspace{-1mm}\noindent\rule{\textwidth}{.4pt}\par\vspace{2mm}
\begin{center}\begin{tabular}{r||p{78mm}|r@{\hspace{2.2mm}}r|r@{\hspace{2.2mm}}r|c|c||l}\cline{2-8}
 & \centering Stellarum descriptio. & \multicolumn{2}{c|}{Longitudo.} & \multicolumn{2}{c|}{Latitudo.} & \parbox{9mm}{\centering\small Plaga Caeli.} & Magnitudo. & \\ \cline{2-8}
2 & Quae est in collo prope caput. & 175° & 30′ & 19° & 40′ & a & 3 & ε \\
3 & Quae est in ala anteriore dextera. & 174 & 40 & 14 & 50 & a & 3 & γ \\
4 & Praecedens e duabus quae sunt in ala po-\newline\hspace*{4mm}steriore. & \raisebox{-1\normalbaselineskip}[0pt][0pt]{177} & \raisebox{-1\normalbaselineskip}[0pt][0pt]{50} & \raisebox{-1\normalbaselineskip}[0pt][0pt]{12} & \raisebox{-1\normalbaselineskip}[0pt][0pt]{30} & \raisebox{-1\normalbaselineskip}[0pt][0pt]{a} & \raisebox{-1\normalbaselineskip}[0pt][0pt]{3} & \raisebox{-1\normalbaselineskip}[0pt][0pt]{δ} \\
5 & Posterior e duabus iis. & 178 & 10 & 11 & 45 & a & 4 & η \\
6 & Quae est in extremitate pedis prope Hydram. & 181 & 40 & 18 & 10 & a & 3 & β \\
 & \multicolumn{7}{c||}{\parbox{150mm}{\centering\rule{0pt}{5mm}\textbf{Centaurus} (\textit{Qinṭāwuros} sive imago hominis equique, et appellatur etiam\\\textit{az–ẓulmān} « struthiocameli mares »).}} & \\ \cline{2-8}
1 & Australis e quatuor quae sunt in capite. & 201 & 40 & 21 & 40 & a & 5 magna & \textit{g} \\
2 & Borealis ex iis. & 201 & 10 & 18 & 50 & a & 5 magna & \textit{h} \\
3 & Praecedens e duabus mediis reliquis. & 200 & 20 & 20 & 30 & a & 4 magna & \textit{i} \\
4 & Posterior ex his duabus, et est reliqua ex\newline\hspace*{4mm}quatuor. & \raisebox{-1\normalbaselineskip}[0pt][0pt]{201} & \raisebox{-1\normalbaselineskip}[0pt][0pt]{10} & \raisebox{-1\normalbaselineskip}[0pt][0pt]{20} & \raisebox{-1\normalbaselineskip}[0pt][0pt]{0} & \raisebox{-1\normalbaselineskip}[0pt][0pt]{a} & \raisebox{-1\normalbaselineskip}[0pt][0pt]{5 magna} & \raisebox{-1\normalbaselineskip}[0pt][0pt]{\textit{k}} \\
5 & Quae est in humero anteriore sinistro. & 197 & 20 & 25 & 40 & a & 3 & ι \\
6 & Quae est in humero posteriore dextero. & 206 & 50 & 22 & 30 & a & 3 & θ \\
7 & E duabus reliquis quae est in summo surculo. & 213 & 10 & 18 & 15 & a & 4 magna & π \\
8 & Praecedens e tribus quae sunt in latere\newline\hspace*{4mm}dextero. & \raisebox{-1\normalbaselineskip}[0pt][0pt]{204} & \raisebox{-1\normalbaselineskip}[0pt][0pt]{30} & \raisebox{-1\normalbaselineskip}[0pt][0pt]{28} & \raisebox{-1\normalbaselineskip}[0pt][0pt]{20} & \raisebox{-1\normalbaselineskip}[0pt][0pt]{a} & \raisebox{-1\normalbaselineskip}[0pt][0pt]{4 magna} & \raisebox{-1\normalbaselineskip}[0pt][0pt]{τ} \\
9 & Media ex iis. & 205 & 10 & 29 & 20 & a & 4 magna & υ \\
10 & Posterior ex his tribus. & 206 & 20 & 28 & 0 & a & 4 magna & φ \\
11 & Quae est in brachio dextero. & 207 & 30 & 26 & 30 & a & 4 magna & \textit{m} \\
12 & Quae est in parte anteriore brachii dexteri. & 214 & 0 & 25 & 15 & a & 3 & ϰ \\
13 & Lucida quae est in initio lateris sinistri. & 209 & 10 & 33 & 30 & a & 3 & λ \\
14 & Posterior ex obscuris a septentrionibus eius. & 208 & 50 & 31 & 0 & a & 5 & \textit{n} \\
15 & Praecedens ex obscuris a septentrionibus\newline\hspace*{4mm}eius. & \raisebox{-1\normalbaselineskip}[0pt][0pt]{208} & \raisebox{-1\normalbaselineskip}[0pt][0pt]{0} & \raisebox{-1\normalbaselineskip}[0pt][0pt]{30} & \raisebox{-1\normalbaselineskip}[0pt][0pt]{20} & \raisebox{-1\normalbaselineskip}[0pt][0pt]{a} & \raisebox{-1\normalbaselineskip}[0pt][0pt]{5} & \raisebox{-1\normalbaselineskip}[0pt][0pt]{χ} \\
\cline{2-8}\end{tabular}\end{center}
{\small
\par\hangindent=12mm\hangafter=1 2. Long. cod. \textarabic{قفح} 188° pro \textarabic{قعه} 175°.
\par\hangindent=12mm\hangafter=1 4. Long. cod. \textarabic{ل} 30′ pro \textarabic{ن} 50′.
\par\hangindent=12mm\hangafter=1 5. Lat. cod. \textarabic{ن} 50′ pro \textarabic{ى} 10′. — Magnitudo \textarabic{د} 4 pro \textarabic{ـحـ} = \textarabic{ـجـ} 3.
\par\hspace*{8mm}\textbf{Centaurus}:
\par\hangindent=12mm\hangafter=1 1. Long. cod. \textarabic{لـ} 20′ pro \textarabic{م} 40′.
\par\hangindent=12mm\hangafter=1 1-4. Magnitudines ut Pts.; L. et S. « magnam » omittunt.
\par\hangindent=12mm\hangafter=1 4. Cod. « et est secunda ex quatuor »; scilicet \textarabic{الثاني} pro \textarabic{الباقي}.
\par\hangindent=12mm\hangafter=1 6. Long. cod. \textarabic{ل} 30′ pro \textarabic{ن} 50′.
\par\hangindent=12mm\hangafter=1 7. Magnitudo L., Pts. et S. 4, sine « magna ».
\par\hangindent=12mm\hangafter=1 10. Magnitudo ut L. et Pts.; S. 4, sine « magna ».
\par\hangindent=12mm\hangafter=1 13. Magnitudo ut S; L. et Pts. 3 magna.
\par\hangindent=12mm\hangafter=1 14. Magnitudo cod. \textarabic{د} 4 pro \textarabic{ه} 5.
\par\hangindent=12mm\hangafter=1 15. Male H. lat. 33° 0′
\par}
\clearpage
\noindent\makebox[\textwidth]{\textbf{174}\hfill{\small C. A. NALLINO}\hfill\textbf{}}\par\vspace{-1mm}\noindent\rule{\textwidth}{.4pt}\par\vspace{2mm}
\begin{center}\begin{tabular}{r||p{78mm}|r@{\hspace{2.2mm}}r|r@{\hspace{2.2mm}}r|c|c||l}\cline{2-8}
 & \centering Stellarum descriptio. & \multicolumn{2}{c|}{Longitudo.} & \multicolumn{2}{c|}{Latitudo.} & \parbox{9mm}{\centering\small Plaga Caeli.} & Magnitudo. & \\ \cline{2-8}
16 & Postrema e tribus quae sunt in dorso a\newline\hspace*{4mm}dextera. & \raisebox{-1\normalbaselineskip}[0pt][0pt]{197°} & \raisebox{-1\normalbaselineskip}[0pt][0pt]{0′} & \raisebox{-1\normalbaselineskip}[0pt][0pt]{40°} & \raisebox{-1\normalbaselineskip}[0pt][0pt]{0′} & \raisebox{-1\normalbaselineskip}[0pt][0pt]{a} & \raisebox{-1\normalbaselineskip}[0pt][0pt]{3} & \raisebox{-1\normalbaselineskip}[0pt][0pt]{μ} \\
17 & Praecedens e duabus sibi invicem vicinis in\newline\hspace*{4mm}femore dextero equi. & \raisebox{-1\normalbaselineskip}[0pt][0pt]{193} & \raisebox{-1\normalbaselineskip}[0pt][0pt]{50} & \raisebox{-1\normalbaselineskip}[0pt][0pt]{46} & \raisebox{-1\normalbaselineskip}[0pt][0pt]{10} & \raisebox{-1\normalbaselineskip}[0pt][0pt]{a} & \raisebox{-1\normalbaselineskip}[0pt][0pt]{3} & \raisebox{-1\normalbaselineskip}[0pt][0pt]{β} \\
18 & Praecedens e duabus quae sunt sub ventre\newline\hspace*{4mm}equi. & \raisebox{-1\normalbaselineskip}[0pt][0pt]{207} & \raisebox{-1\normalbaselineskip}[0pt][0pt]{30} & \raisebox{-1\normalbaselineskip}[0pt][0pt]{43} & \raisebox{-1\normalbaselineskip}[0pt][0pt]{0} & \raisebox{-1\normalbaselineskip}[0pt][0pt]{a} & \raisebox{-1\normalbaselineskip}[0pt][0pt]{2} &  \\
19 & Posterior e duabus iis. & 208 & 50 & 43 & 45 & a & 3 &  \\
20 & Quae est in poplite dextero prope pedem. & 201 & 10 & 51 & 10 & a & 2 & ν \\
21 & Quae est in ungula pedis sinistri. & 202 & 20 & 55 & 20 & a & 2 & ζ \\
22 & Lucida quae est in extremitate pedis dex-\newline\hspace*{4mm}teri anterioris et vocatur \textit{riǵl al–faras}\newline\hspace*{4mm}« pes equi ». & \raisebox{-2\normalbaselineskip}[0pt][0pt]{199} & \raisebox{-2\normalbaselineskip}[0pt][0pt]{30} & \raisebox{-2\normalbaselineskip}[0pt][0pt]{41} & \raisebox{-2\normalbaselineskip}[0pt][0pt]{10} & \raisebox{-2\normalbaselineskip}[0pt][0pt]{a} & \raisebox{-2\normalbaselineskip}[0pt][0pt]{1} & \raisebox{-2\normalbaselineskip}[0pt][0pt]{α} \\
23 & Quae est in genu sinistro a parte pedis. & 215 & 20 & 45 & 20 & a & 2 & γ \\
24 & Quae est in malleolo dextero. & 206 & 30 & 51 & 40 & a & 2 & ξ \\
25 & Sexta ex iis quae sunt in pede posteriore\newline\hspace*{4mm}dextero. & \raisebox{-1\normalbaselineskip}[0pt][0pt]{205} & \raisebox{-1\normalbaselineskip}[0pt][0pt]{50} & \raisebox{-1\normalbaselineskip}[0pt][0pt]{49} & \raisebox{-1\normalbaselineskip}[0pt][0pt]{10} & \raisebox{-1\normalbaselineskip}[0pt][0pt]{a} & \raisebox{-1\normalbaselineskip}[0pt][0pt]{4} & \raisebox{-1\normalbaselineskip}[0pt][0pt]{ε} \\
26 & Quae est sub medio poplitis sinistri. & 197 & 30 & 55 & 10 & a & 4 & \textit{f} \\
 & \multicolumn{7}{c||}{\parbox{150mm}{\centering\rule{0pt}{5mm}\textbf{Lupus} (\textit{as–sabuʿ} « bellua »).}} & \\ \cline{2-8}
1 & Quae est in extremitate pedis posterioris. & 219 & 10 & 24 & 50 & a & 3 & ο \\
\cline{2-8}\end{tabular}\end{center}
{\small
\par\hangindent=12mm\hangafter=1 16. Long. cod. \textarabic{قضج} 193° pro \textarabic{قضز} 197°.
\par\hangindent=12mm\hangafter=1 17. Long. cod. \textarabic{ل} 30′ pro \textarabic{ن} 50′.
\par\hangindent=12mm\hangafter=1 18. Long. cod. \textarabic{لـ} 20′ pro \textarabic{ل} 30′. — \Name{Baily}, in adnotationibus ad Catalogum \Name{Ulugh Beg}: « I cannot « identify this star ». \Name{Schjellerup} ε Centauri, quam \Name{Baily} e contra putat nr. 25 Catalogi Albateniani respondere.
\par\hangindent=12mm\hangafter=1 19. \Name{Ṣūfī} p. 244 et 249 ait nullam hoc Caeli loco stellam exsistere. \Name{Baily}, in adnotationibus ad Catalogum \Name{Ulugh Beg}: « There is no star of any considerable magnitude near the spot indicated by Pto- « lemy, and unobserved by him [i. e. Ulugh Beg], except 1088 in Lacaille’s catalogue; but neither « the longitude nor latitude will agree with Ptolemy’s description or position ».
\par\hangindent=12mm\hangafter=1 20. Reliqui Arabes eam \textit{riǵl qinṭāwuros} « pedem Centauri » vocant. — Long. ut L., sed ex errore in signo zodiacali provenit; legerunt enim 9° 30′ Librae pro 9° 30′ Scorpii (= 229° 30′) ut B et S. habent. H. 4° 10′ Scorpii (= 224° 10′). Errorem Albatenii et plurium codicum Almagesti memorat \Name{Ṣūfī} p. 32, qui testatur in eundem illapsos esse auctores \textit{Tabulae probatae}. — Male H. lat. 42° 20′.
\par\hangindent=12mm\hangafter=1 23. Long. cod. \textarabic{ريو \AbjadZero{}} 216° 0′ pro \textarabic{ريه ك} 215° 20′.
\par\hangindent=12mm\hangafter=1 24. Long. cod. \textarabic{رر} 207° pro \textarabic{رو} 206°.
\par\hangindent=12mm\hangafter=1 25. Pt. ὁ ἐκτὸς ὑπὸ τὸν δεξιὸν ὀπισθιόποδα « quae extra [est], sub dextero pede posteriore »; sed versio qua al–Battānī usus est legerat in Graeco ἕκτος i. e. « sexta » (cfr. ad nr. 1 Coronae australis). Idem error H., qui vertit « la sixième ». — Magnitudo cod. \textarabic{و} 6 pro \textarabic{د} 4; B., H., Pts. et S. 4; L. 3. — Perperam \Name{Schjellerup} ϑ Centauri, quae est nr. 6.
\par\hangindent=12mm\hangafter=1 26. Magnitudo ut L. et Pts.; B. et H. 2; S. 3 parva; hodie 3. Situs Ptolemaicus magno errore affectus. \Name{Schjellerup} stellam putat esse δ Crucis.
\par\hspace*{8mm}\textbf{Lupus}:
\par\hangindent=12mm\hangafter=1 1. Nomen ut H. et S.; rectius B. non ὀπισθίου « posterioris » sed ἐμπροσθίου « anterioris ».
\par}
\clearpage
\noindent\makebox[\textwidth]{\textbf{}\hfill{\small AL-BATTANI OPUS ASTRONOMICUM}\hfill\textbf{175}}\par\vspace{-1mm}\noindent\rule{\textwidth}{.4pt}\par\vspace{2mm}
\begin{center}\begin{tabular}{r||p{78mm}|r@{\hspace{2.2mm}}r|r@{\hspace{2.2mm}}r|c|c||l}\cline{2-8}
 & \centering Stellarum descriptio. & \multicolumn{2}{c|}{Longitudo.} & \multicolumn{2}{c|}{Latitudo.} & \parbox{9mm}{\centering\small Plaga Caeli.} & Magnitudo. & \\ \cline{2-8}
2 & Quae est in medio poplitis posterioris. & 217° & 0′ & 29° & 10′ & a & 3 & α \\
3 & Borealis e tribus quae sunt in extremitate\newline\hspace*{4mm}caudae. & \raisebox{-1\normalbaselineskip}[0pt][0pt]{214} & \raisebox{-1\normalbaselineskip}[0pt][0pt]{10} & \raisebox{-1\normalbaselineskip}[0pt][0pt]{29} & \raisebox{-1\normalbaselineskip}[0pt][0pt]{20} & \raisebox{-1\normalbaselineskip}[0pt][0pt]{a} & \raisebox{-1\normalbaselineskip}[0pt][0pt]{4 magna} & \raisebox{-1\normalbaselineskip}[0pt][0pt]{ϰ} \\
4 & Australis e duabus quae sunt in pede an-\newline\hspace*{4mm}teriore. & \raisebox{-1\normalbaselineskip}[0pt][0pt]{218} & \raisebox{-1\normalbaselineskip}[0pt][0pt]{20} & \raisebox{-1\normalbaselineskip}[0pt][0pt]{11} & \raisebox{-1\normalbaselineskip}[0pt][0pt]{30} & \raisebox{-1\normalbaselineskip}[0pt][0pt]{a} & \raisebox{-1\normalbaselineskip}[0pt][0pt]{4 magna} & \raisebox{-1\normalbaselineskip}[0pt][0pt]{ε} \\
5 & Borealis e duabus quae sunt in collo. & 230 & 30 & 15 & 20 & a & 4 magna & μ \\
6 & Borealis e duabus quae sunt in pede ante-\newline\hspace*{4mm}riore. & \raisebox{-1\normalbaselineskip}[0pt][0pt]{217} & \raisebox{-1\normalbaselineskip}[0pt][0pt]{40} & \raisebox{-1\normalbaselineskip}[0pt][0pt]{10} & \raisebox{-1\normalbaselineskip}[0pt][0pt]{0} & \raisebox{-1\normalbaselineskip}[0pt][0pt]{a} & \raisebox{-1\normalbaselineskip}[0pt][0pt]{4 magna} & \raisebox{-1\normalbaselineskip}[0pt][0pt]{δ} \\
 & \multicolumn{7}{c||}{\parbox{150mm}{\centering\rule{0pt}{5mm}\textbf{Ara} (\textit{al–miǵmarah} « thuribulum » sive \textit{al–murīḥ} « odores spargens »).}} & \\ \cline{2-8}
1 & Borealis e duabus quae sunt in base thu-\newline\hspace*{4mm}ribuli. & \raisebox{-1\normalbaselineskip}[0pt][0pt]{248} & \raisebox{-1\normalbaselineskip}[0pt][0pt]{50} & \raisebox{-1\normalbaselineskip}[0pt][0pt]{22} & \raisebox{-1\normalbaselineskip}[0pt][0pt]{40} & \raisebox{-1\normalbaselineskip}[0pt][0pt]{a} & \raisebox{-1\normalbaselineskip}[0pt][0pt]{5} & \raisebox{-1\normalbaselineskip}[0pt][0pt]{γ} \\
2 & Quae est in medio partis superioris i. e. in\newline\hspace*{4mm}loco ignis. & \raisebox{-1\normalbaselineskip}[0pt][0pt]{247} & \raisebox{-1\normalbaselineskip}[0pt][0pt]{20} & \raisebox{-1\normalbaselineskip}[0pt][0pt]{26} & \raisebox{-1\normalbaselineskip}[0pt][0pt]{30} & \raisebox{-1\normalbaselineskip}[0pt][0pt]{a} & \raisebox{-1\normalbaselineskip}[0pt][0pt]{4 magna} & \raisebox{-1\normalbaselineskip}[0pt][0pt]{α} \\
3 & Australis e duabus sibi invicem proximis\newline\hspace*{4mm}quae sunt in flamma. & \raisebox{-1\normalbaselineskip}[0pt][0pt]{246} & \raisebox{-1\normalbaselineskip}[0pt][0pt]{20} & \raisebox{-1\normalbaselineskip}[0pt][0pt]{34} & \raisebox{-1\normalbaselineskip}[0pt][0pt]{10} & \raisebox{-1\normalbaselineskip}[0pt][0pt]{a} & \raisebox{-1\normalbaselineskip}[0pt][0pt]{4 magna} & \raisebox{-1\normalbaselineskip}[0pt][0pt]{ζ} \\
4 & Quae est in extremitate flammae. & 242 & 0 & 34 & 0 & a & 4 &  \\
5 & [Australis e duabus quae sunt in base thu-\newline\hspace*{4mm}ribuli.] & \raisebox{-1\normalbaselineskip}[0pt][0pt]{254} & \raisebox{-1\normalbaselineskip}[0pt][0pt]{10} & \raisebox{-1\normalbaselineskip}[0pt][0pt]{25} & \raisebox{-1\normalbaselineskip}[0pt][0pt]{45} & \raisebox{-1\normalbaselineskip}[0pt][0pt]{a} & \raisebox{-1\normalbaselineskip}[0pt][0pt]{4} &  \\
 & \multicolumn{7}{c||}{\parbox{150mm}{\centering\rule{0pt}{5mm}\textbf{Corona australis} (\textit{al–iklīl al–ǵanūbī}).}} & \\ \cline{2-8}
1 & Praecedens e sex quae sunt in curvatura\newline\hspace*{4mm}australi. & \raisebox{-1\normalbaselineskip}[0pt][0pt]{260} & \raisebox{-1\normalbaselineskip}[0pt][0pt]{20} & \raisebox{-1\normalbaselineskip}[0pt][0pt]{21} & \raisebox{-1\normalbaselineskip}[0pt][0pt]{30} & \raisebox{-1\normalbaselineskip}[0pt][0pt]{a} & \raisebox{-1\normalbaselineskip}[0pt][0pt]{4} & \raisebox{-1\normalbaselineskip}[0pt][0pt]{α} \\
\cline{2-8}\end{tabular}\end{center}
{\small
\par\hangindent=12mm\hangafter=1 2. Long. cod. \textarabic{ريو م} 216° 40′ pro \textarabic{رير \AbjadZero{}} 217° 0′.
\par\hangindent=12mm\hangafter=1 4. Lat. ut L. et S.; B. et H. 11° 50′. — Magnitudo ut L.; Pts. 4; S. 6 parva.
\par\hangindent=12mm\hangafter=1 6. Male H. long. 218° 40′. — Magnitudo ut L. et Pts.; S. 5 parva.
\par\hspace*{8mm}\textbf{Ara} (1):
\par\hangindent=12mm\hangafter=1 2. Long. ut L. et S.; B. 247° 30′; H. 249° 30′.
\par\hangindent=12mm\hangafter=1 4. Lat. ut L. et S.; B. et H. 34° 15′.
\par\hangindent=12mm\hangafter=1 5. Scriptor cod. nomen stellae e neglegentia omisit. — Long. ut H.; B. 254° 20′; S. 251° 30′ false; L 221° 30′ (i. e. 11° 30′ Scorpii loco Sagittarii). — Lat. cod. \textarabic{كا} 21° pro \textarabic{كه} 25°.
\par\hspace*{8mm}\textbf{Corona australis}:
\par\hangindent=12mm\hangafter=1 1. Pt.: τῆς νοτίου περιφερείας ὁ προηγούμενος ἕκτος « peripheriae australis sexta praecedens ». Versiones Almagesti quibus S. et \Name{Ulugh Beg} usi sunt, in Graeco ἐκτός « extra » pro ἕκτος legerant (cfr. ad nr. 25 Centauri): « praecedens extra peripheriam australem ».
\par}
\par\vspace{2mm}\begin{center}\rule{40mm}{.4pt}\end{center}{\small (1) Litterae Bayer ad hanc constellationem pertinentes, teste \Name{Baily} in tabula erratorum (et p. 176, ad Catalogum Halleyanum stellarum australium), fere omnes « are prob bly erroneous and it is very difficult « to apply them correctly ». \Name{Schjellerup} eas hoc modo indicat: σ, α, γ, ζ, ϑ. Numeri in Catalogo Lacaille (\textit{Caelum australe stelliferum}, Parisiis 1763), iuxta \Name{Baily}, sunt: 1443, 1436, 1422, 1386, 1480 Lac.}
\clearpage
\noindent\makebox[\textwidth]{\textbf{176}\hfill{\small C. A. NALLINO}\hfill\textbf{}}\par\vspace{-1mm}\noindent\rule{\textwidth}{.4pt}\par\vspace{2mm}
\begin{center}\begin{tabular}{r||p{78mm}|r@{\hspace{2.2mm}}r|r@{\hspace{2.2mm}}r|c|c||l}\cline{2-8}
 & \centering Stellarum descriptio. & \multicolumn{2}{c|}{Longitudo.} & \multicolumn{2}{c|}{Latitudo.} & \parbox{9mm}{\centering\small Plaga Caeli.} & Magnitudo. & \\ \cline{2-8}
2 & Quarta ex his sex. & 266° & 0′ & 20° & 0′ & a & 4 & β \\
3 & Quae eam sequitur ante genu Sagittarii. & 267 & 20 & 18 & 30 & a & 5 & η \\
4 & Lucida quae hanc sequitur a septentrionibus. & 268 & 10 & 17 & 10 & a & 2 & θ \\
5 & Stella a septentrionibus huius lucidae. & 268 & 0 & 16 & 0 & a & 4 & γ \\
6 & Praecedens e duabus obscuris. & 265 & 50 & 14 & 50 & a & 6 & ν \\
7 & Reliqua e duabus obscuris. & 263 & 0 & 14 & 40 & a & 5 & ι \\
8 & Stella quae hanc ipsam praecedit. & 260 & 50 & 15 & 50 & a & 5 & ϰ \\
9 & Reliqua ab austro huius. & 260 & 20 & 18 & 30 & a & 5 & λ \\
 & \multicolumn{7}{c||}{\parbox{150mm}{\centering\rule{0pt}{5mm}\textbf{Piscis australis} (\textit{al–ḥūt al–ǵanūbī}).}} & \\ \cline{2-8}
1 & Quae est in ore Piscis in extremitate aquae. & . . . & . . & . . & . . & . . & . . . . . & α \\
2 & [Praecedens e tribus in curvatura australi\newline\hspace*{4mm}capitis.] & \raisebox{-1\normalbaselineskip}[0pt][0pt]{311} & \raisebox{-1\normalbaselineskip}[0pt][0pt]{50} & \raisebox{-1\normalbaselineskip}[0pt][0pt]{20} & \raisebox{-1\normalbaselineskip}[0pt][0pt]{20} & \raisebox{-1\normalbaselineskip}[0pt][0pt]{a} & \raisebox{-1\normalbaselineskip}[0pt][0pt]{4} & \raisebox{-1\normalbaselineskip}[0pt][0pt]{β} \\
3 & Quae hanc sequitur. & 315 & 20 & 22 & 15 & a & 4 & γ \\
4 & Tertia postrema e tribus his praecedentibus. & 316 & 30 & 22 & 30 & a & 4 & δ \\
5 & Quae est in gutture Piscis. & 315 & 30 & 16 & 15 & a & 4 magna & ε \\
6 & Australis quae est in spina dorsi australi. & 306 & 20 & 19 & 30 & a & 5 & μ \\
7 & Posterior e duabus quae sunt in ventre. & 312 & 20 & 15 & 10 & a & 5 & ζ \\
8 & Praecedens e duabus iis. & 310 & 0 & 14 & 40 & a & 4 & λ \\
9 & Postrema e tribus quae sunt in spina dorsi\newline\hspace*{4mm}boreali. & \raisebox{-1\normalbaselineskip}[0pt][0pt]{306} & \raisebox{-1\normalbaselineskip}[0pt][0pt]{20} & \raisebox{-1\normalbaselineskip}[0pt][0pt]{15} & \raisebox{-1\normalbaselineskip}[0pt][0pt]{0} & \raisebox{-1\normalbaselineskip}[0pt][0pt]{a} & \raisebox{-1\normalbaselineskip}[0pt][0pt]{4} & \raisebox{-1\normalbaselineskip}[0pt][0pt]{η} \\
10 & Media ex his tribus. & 303 & 0 & 16 & 30 & a & 4 & θ \\
11 & Praecedens ex his tribus. & 302 & 10 & 18 & 10 & a & 4 & ι \\
12 & Quae est in extremitate caudae. & 301 & 20 & 22 & 15 & a & 4 & ϰ \\
\cline{2-8}\end{tabular}\end{center}
{\small
\par\hangindent=12mm\hangafter=1 4. Magnitudo Pt. 4, S. 5; sed cfr. nr. 72 tabulae statuum fixarum ad annum 1211.
\par\hangindent=12mm\hangafter=1 5. Long. cod. \textarabic{رعب} 272° pro \textarabic{رصح} 268° quae L. et S. habent; B. et H. 267° 30′.
\par\hangindent=12mm\hangafter=1 7. Lat. ut H., L. et S.; B. 15° 50′.
\par\hangindent=12mm\hangafter=1 8. Long. cod. \textarabic{ل} 30′ pro \textarabic{ن} 50′.
\par\hangindent=12mm\hangafter=1 9. Lat. cod. \textarabic{ن} 50′ pro \textarabic{ل} 30′; B., H. et L. 18° 30′; S. 8° 0′.
\par\hspace*{8mm}\textbf{Piscis australis}:
\par\hangindent=12mm\hangafter=1 1. Est α Piscis austr., quae iam in constellatione Aquarii, nr. 18, a Pt. et Albatenio descripta est. Itaque al–Battānī eius numeros hic omisit, ut quoque L. fecit. Scriptor cod., causam lacunae ignorans, numeros secundae stellae tribuit primae, et nomen secundae stultissime omisit, ne vacuum aliquid appareret.
\par\hangindent=12mm\hangafter=1 3. Lat. cod. \textarabic{ى} 10′ pro \textarabic{ل} 30′.
\par\hangindent=12mm\hangafter=1 5. Lat. cod. \textarabic{ى} 10′ pro \textarabic{يه} 15′. — Magnitudo L., Pts., S. et Arg. 4, sine « magna ».
\par\hangindent=12mm\hangafter=1 6. Lat. cod. \textarabic{ط} 9° pro \textarabic{ىط} = \textarabic{يط} 19°.
\par\hangindent=12mm\hangafter=1 7. Male H. long. 313° 20′. — Magnitudo cod. \textarabic{د} 4 pro \textarabic{ه} 5 (S. 6 parva).
\par\hangindent=12mm\hangafter=1 8. Magnitudo cod. \textarabic{ـجـ} 3 pro \textarabic{د} 4 (S. 5 parva).
\par\hangindent=12mm\hangafter=1 9. Lat. cod. \textarabic{يب} 12° pro \textarabic{يه} 15°.
\par}
\clearpage
\noindent\makebox[\textwidth]{\textbf{}\hfill{\small AL-BATTANI OPUS ASTRONOMICUM}\hfill\textbf{177}}\par\vspace{-1mm}\noindent\rule{\textwidth}{.4pt}\par\vspace{2mm}
\begin{center}\begin{tabular}{r||p{78mm}|r@{\hspace{2.2mm}}r|r@{\hspace{2.2mm}}r|c|c||l}\cline{2-8}
 & \centering Stellarum descriptio. & \multicolumn{2}{c|}{Longitudo.} & \multicolumn{2}{c|}{Latitudo.} & \parbox{9mm}{\centering\small Plaga Caeli.} & Magnitudo. & \\ \cline{2-8}
 & \multicolumn{7}{l}{\hspace*{4mm}Et ex iis quae non sunt ei in figura:} & \\
13 & Praecedens e tribus lucidis. & 289° & 10′ & 22° & 20′ & a & 3 parva &  \\
14 & Media ex his tribus. & 292 & 20 & 22 & 10 & a & 3 parva &  \\
15 & Postrema ex his tribus. & 295 & 10 & 21 & 0 & a & 3 parva &  \\
16 & Obscura quae est ante eam. & 293 & 10 & 20 & 50 & a & 5 &  \\
\cline{2-8}\end{tabular}\end{center}
{\small
\par\hangindent=12mm\hangafter=1 13-16. Quatuor hae stellae, quae re vera ad quintam magnitudinem pertinent, a S. et \Name{Ulugh Beg} aspectae non sunt, et in catalogis omissae. — Nr. 13 est fortasse 1694 Lacaille; nr. 14 fortasse 1717 Lac. nr. 16 fortasse 1704 Lac. — Long. nr. 14 in cod. \textarabic{رصب ى} 262° 10′ pro \textarabic{رضب لـ} 292° 20′. — Lat. nr. 15 ut L.; B. et H. 21° 10′. Magnitudo cod. \textarabic{د} 4 pro \textarabic{ـحـ} = \textarabic{ـجـ} 3. — Magnitudo nr. 16 in cod. \textarabic{د} 4 pro \textarabic{ه} 5.
\par}
\par\vfill\hfill 23
\end{document}
